\documentclass[11pt]{article}

\usepackage[margin=1in]{geometry}
\usepackage{amsmath}
\usepackage{amssymb}
\usepackage{hyperref}
\usepackage{booktabs}
\usepackage{listings}
\usepackage{xcolor}
\usepackage{parskip}
\usepackage{array}

\lstdefinestyle{pystyle}{
  basicstyle=\ttfamily\small,
  keywordstyle=\color{blue!70!black},
  commentstyle=\color{gray},
  stringstyle=\color{orange!70!black},
  showstringspaces=false,
  breaklines=true,
  frame=single,
  language=Python,
}
\lstset{style=pystyle}

\title{DMRGPY Physics User Guide}
\author{}
\date{}

\sloppy

\begin{document}
\maketitle

This guide explains what DMRGPY computes from a \textbf{physics} point
of view --- the quantities each method returns, the formula behind
them, and when to reach for which method. It intentionally says nothing
about the solver backends (DMRG vs.\ ED, C++ vs.\ Python vs.\ Julia) or
code layout --- see \texttt{docs/documentation.tex} for that. Every
\texttt{Many\_Body\_Chain} subclass (\texttt{Spin\_Chain},
\texttt{Fermionic\_Chain}, \texttt{Bosonic\_Chain}, \ldots) exposes the
same physics-facing API described here, with \texttt{mode="DMRG"|"ED"}
selectable on almost every call so a result can be cross-checked against
an exact reference on small systems.

\tableofcontents

\section{Physical models and Hilbert spaces}

Every model is a chain of $n$ local Hilbert spaces $\mathcal
H=\bigotimes_{i=1}^n\mathcal H_i$, on which a Hamiltonian and
observables are built out of local operators.

\begin{center}
\footnotesize
\begin{tabular}{@{}p{0.28\textwidth}p{0.32\textwidth}p{0.3\textwidth}@{}}
\toprule
Chain class & Local Hilbert space & Key operators \\
\midrule
\texttt{spinchain.\allowbreak Spin\_Chain} & spin-$S$, $S\in\{\tfrac12,1,\tfrac32,2,\tfrac52,3\}$ per site & $S^x_i,S^y_i,S^z_i$ \\
\texttt{fermionchain.\allowbreak Fermionic\_Chain} & spinless fermion (occupied/empty) & $c_i,c_i^\dagger,n_i=c_i^\dagger c_i$, JW string $F_i$ \\
\texttt{fermionchain.\allowbreak Majorana\_Chain} & Majorana fermion & Majorana operators built from \texttt{Fermionic\_Chain} \\
\texttt{fermionchain.\allowbreak Spinful\_\allowbreak Fermionic\_Chain} & spin-$\tfrac12$ fermion (4 states), two interleaved spinless sites per physical site & $c_{i\sigma},c^\dagger_{i\sigma},n_{i\sigma}$; derived spin and pairing operators (see text) \\
\texttt{fermionchain.\allowbreak Spinful\_\allowbreak Fermionic\_\allowbreak Chain\_\allowbreak Native} & same physics, genuinely 4-dimensional local space, one site per physical site (\texttt{itensor\_version=3} only) & identical operator lists/formulas as \texttt{Spinful\_Fermionic\_Chain} \\
\texttt{bosonchain.\allowbreak Bosonic\_Chain} & truncated boson Fock space, $n_i\in\{0,\ldots,\text{maxnb}_i-1\}$, per-site dimension \texttt{maxnb} (default 4) settable via \texttt{Bosonic\_Chain(n,\allowbreak maxnb=[\ldots])} & $a_i,a_i^\dagger,n_i$, occupation projectors $\hat n_i^{(k)}$ for $k=0,\ldots,\text{maxnb}_i-1$ \\
\texttt{parafermionchain.\allowbreak Parafermionic\_\allowbreak Chain} & $\mathbb Z_N$ parafermion (clock model), $N\in\{2,3,4\}$ & clock/shift operators $\sigma_i,\tau_i$, composite $\chi_i,\psi_i$ \\
\texttt{mixedchain.\allowbreak Mixed\_Spin\_\allowbreak Fermion\_Chain} & mixes spin-$S$ sites and spinful-fermion locations in the same chain & spin location: native $S^x_i,S^y_i,S^z_i$; fermion location: $c_{i\sigma},c^\dagger_{i\sigma},n_{i\sigma}$ plus derived spin/pairing operators \\
\bottomrule
\end{tabular}
\normalsize
\end{center}

Spinful fermionic chains are built by \emph{interleaving} two spinless
fermionic sites per physical site (site $2i$ = spin up, site $2i+1$ =
spin down) rather than by a genuinely 4-dimensional local space, so that
the same Jordan-Wigner machinery used for spinless fermions applies
unchanged; \texttt{Spinful\_Fermionic\_Chain} wraps this bookkeeping for
you.

\texttt{Spinful\_Fermionic\_Chain\_Native} is the alternative built
directly on a genuinely 4-dimensional local space (ITensor v3's own
\texttt{Electron}/\texttt{Hubbard} site type) instead: one
tensor-network site per physical site, with the same operator lists and
identical physics/sign convention as \texttt{Spinful\_Fermionic\_Chain}
-- the two classes are drop-in equivalent for any given Hamiltonian,
cross-checked to agree exactly under ED and to DMRG tolerance under
\texttt{itensor\_version=3} (the only DMRG backend this class wires
up). Despite halving the site count, it is \emph{not} generally faster
in practice: two-site DMRG's per-sweep cost is driven by the local
dimension of the two-site block being diagonalized, which grows faster
($4\times4=16$ per pair of native sites, versus $2\times2=4$ per pair of
interleaved sites) than the site-count halving saves. The same
disadvantage held up under every other regime checked too: strong
on-site coupling, long-range/power-law hopping, two-site and
one-site+subspace-expansion (\texttt{"TDVP\_GSE"}) real-time evolution,
and the KPM dynamical correlator itself -- see
\texttt{fermionchain.py}'s class docstring for the full rundown. One
case does flip in its favor, though: the 4-point correlator tensor
$\langle c^\dagger_i c_j c^\dagger_k c_l\rangle$
(\texttt{mps.MPS.get\_four\_correlation\_tensor()}, \S5) is a Python
loop of independent \emph{static} overlaps rather than an iterative
two-site search, so it does not pay the two-site combined-local-
dimension penalty above. Both classes support a \texttt{ctmode="full"}
C++-accelerated path in addition to the generic
\texttt{ctmode="explicit"} one --
\texttt{Spinful\_Fermionic\_Chain\_Native} gets its own
(\texttt{Chain::four\_correlation\_tensor\_spinful()}, using ITensor's
own automatic Jordan-Wigner insertion on the flavor-resolved operator
names, since ITensor's \texttt{ElectronSite} has no bare
\texttt{"Cdag"}/\texttt{"C"} the plain version needs). Measured
($n=3,4,5,6,12$ orbitals), \texttt{Spinful\_Fermionic\_Chain\_Native}'s
\texttt{ctmode="full"} is the fastest of all four combinations at every
size tried, including $n=12$ (24 flat modes: $\sim620$s vs $\sim890$s
for \texttt{Spinful\_Fermionic\_Chain}'s own \texttt{ctmode="full"}, a
$\sim30\%$ win). Prefer \texttt{ctmode="full"} for this class whenever
it's available (it always is, for \texttt{itensor\_version=3}).
Otherwise prefer \texttt{Spinful\_Fermionic\_Chain}; no other
calculation tried so far makes the native-site class faster.

\texttt{Mixed\_Spin\_Fermion\_Chain} is for models that need a literal
local moment next to a conduction-electron site (e.g.\ Kondo-lattice-like
Hamiltonians), rather than the large-$U$ two-fermion-site trick
\texttt{Spinful\_Fermionic\_Chain}/\texttt{spinfermionchain.py} use to
emulate a spin. Its \texttt{sitesin} constructor argument is a list with
one entry per logical location, each either a spin label (\texttt{"1/2"},
\texttt{"1"}, \ldots{} as in \texttt{Spin\_Chain}) or a fermion marker
(\texttt{"F"}); a fermion location expands internally to a spin-up/spin-down
site pair, exactly like \texttt{Spinful\_Fermionic\_Chain}. All operator
lists (\texttt{Sx}/\texttt{Sy}/\texttt{Sz},
\texttt{Cup}/\texttt{Cdagup}/\texttt{Cdn}/\texttt{Cdagdn}/\texttt{Nup}/
\texttt{Ndn}/\texttt{Ntot}/\texttt{Delta}) are indexed by \emph{logical}
location, not by physical site -- the fermion-only operators read as the
literal integer 0 at spin locations, since they have no meaning there.
Currently only \texttt{itensor\_version=3} (and \texttt{"python"}) are
supported; see \texttt{mixedchain.py}'s module docstring for why
\texttt{itensor\_version=2} isn't yet, and
\texttt{examples/fermion\_models/mixed\_spin\_fermion\_chain} for a worked Kondo-lattice
example.

\texttt{Bosonic\_Chain(n, maxnb=[\ldots])} takes a per-site local
dimension list (defaulting to \texttt{[4]*n}, i.e.\ up to 3 bosons/site);
ED always honors it exactly (\texttt{pyboson/boson.py}). On the DMRG
side, \texttt{itensor\_version=3} and \texttt{itensor\_version="python"}
both thread a non-default \texttt{maxnb} through to the tensor-network
site itself (encoded as the site type code $100+\text{maxnb}_i$, see
\texttt{mpscpp3/get\_sites.h}/\texttt{extra/bosonfour.h} and, on the
pure-Python side, \texttt{pyitensor/sites/boson.py}'s
\texttt{get\_boson\_site()} factory) -- \texttt{itensor\_version=2} and
the Julia backend still only understand the single fixed 4-level boson
site regardless of what \texttt{maxnb} requests, so a non-default
\texttt{maxnb} should be run under \texttt{itensor\_version=3} (the
default) or \texttt{"python"} for DMRG/ED results to actually agree;
see \texttt{examples/boson\_models/boson\_maxnb\_v3\_VS\_ED}.

\section{Building a Hamiltonian and observables}

Any Hamiltonian or observable is written exactly as its
second-quantized or spin-operator expression, using operator lists
indexed by site, e.g.\ for the spin-$\tfrac12$ Heisenberg chain

\begin{equation}
H=\sum_{i=1}^{n-1}J\,\mathbf S_i\cdot\mathbf S_{i+1}=\sum_{i=1}^{n-1}J\left(S_i^xS_{i+1}^x+S_i^yS_{i+1}^y+S_i^zS_{i+1}^z\right)
\end{equation}

\begin{lstlisting}
from dmrgpy import spinchain
spins = ["S=1/2" for i in range(30)]
sc = spinchain.Spin_Chain(spins)
h = 0
for i in range(29):
    h = h + sc.Sx[i]*sc.Sx[i+1] + sc.Sy[i]*sc.Sy[i+1] + sc.Sz[i]*sc.Sz[i+1]
sc.set_hamiltonian(h)
\end{lstlisting}

Any sum of products of these operators is a valid \texttt{MultiOperator},
so essentially any local (or long-range) Hamiltonian on the given
lattice can be built term by term: single-ion anisotropy
$D(S_i^z)^2$, Dzyaloshinskii-Moriya terms $\mathbf D\cdot(\mathbf
S_i\times\mathbf S_{i+1})$, biquadratic exchange $(\mathbf S_i\cdot
\mathbf S_{i+1})^2$, hopping $t\sum_i(c_i^\dagger c_{i+1}+\text{h.c.})$,
Hubbard interaction $U\sum_i n_{i\uparrow}n_{i\downarrow}$, and so on ---
see Section~\ref{sec:cookbook} for concrete Hamiltonians.

\textbf{Algebra on already-built operators.} \texttt{sc.toMPO(h)} compiles
a \texttt{MultiOperator} into a \texttt{StaticOperator} (an already-built
matrix product operator, \texttt{itensor\_version} 2, 3, and
\texttt{"python"} only --- not \texttt{"julia\_live"} yet). Two
\texttt{StaticOperator}s can then be combined directly with \texttt{+},
\texttt{-}, unary \texttt{-}, and scalar \texttt{*}/\texttt{/}, without
going back through the symbolic \texttt{MultiOperator} form:

\begin{lstlisting}
A = sc.toMPO(sc.Sz[0])
B = sc.toMPO(sc.Sx[0]*sc.Sx[1])
C = A + 2*B - A          # still a StaticOperator
\end{lstlisting}

This is a compressed direct sum at the tensor-network level (like
ITensorMPS.jl's \texttt{+(::MPO, ::MPO)}), useful for combining operators
that only exist as already-built \texttt{StaticOperator}s (e.g.\ two
independently constructed products or exponentials); for the common case
of combining Hamiltonians before ever building an MPO, summing the
underlying \texttt{MultiOperator}s directly (as \texttt{h = h + ...}
above) remains preferred.

\section{Ground-state properties}

\textbf{Ground-state energy}

\begin{equation}
E_0=\langle\mathrm{GS}|H|\mathrm{GS}\rangle=\min_{|\psi\rangle}\frac{\langle\psi|H|\psi\rangle}{\langle\psi|\psi\rangle}
\end{equation}

\begin{lstlisting}
e0 = sc.gs_energy()          # E0
wf = sc.get_gs()             # the |GS> wavefunction itself
\end{lstlisting}

\textbf{The DMRG sweep schedule: bond-dimension ramp.}
\lstinline|sc.maxm| is the \emph{target} bond dimension, not the one
every sweep runs at. By default (\lstinline|sc.bond_ramp = True|) the
ground-state sweep schedule spends the first
\lstinline|sc.bond_ramp_fraction| of its \lstinline|sc.nsweeps| sweeps
growing the bond dimension geometrically from
\lstinline|sc.bond_ramp_start| up to \lstinline|sc.maxm|, and holds it at
\lstinline|sc.maxm| for the rest:

\begin{equation}
m_i=\Big\lceil m_\mathrm{start}\Big(\frac{m_\mathrm{max}}{m_\mathrm{start}}\Big)^{i/n_r}\Big\rfloor\ (i<n_r),
\qquad m_i=m_\mathrm{max}\ (i\ge n_r),
\qquad n_r=\lfloor n_\mathrm{sweeps}\,f_\mathrm{ramp}\rfloor
\end{equation}

Two-site DMRG costs $\mathcal O(m^3)$ per bond, while the early sweeps
mostly just locate the right variational subspace and gain little from a
large $m$ --- so running them small makes them nearly free, and the
expensive full-\lstinline|maxm| sweeps then start from an already-good
state. At the default \lstinline|bond_ramp_fraction = 0.5| the second
half of the schedule --- and hence the returned energy --- always runs at
the full \lstinline|sc.maxm|, so the ramp is a pure scheduling change:
the same ground state, less time.

\begin{lstlisting}
sc.bond_ramp = True             # on by default
sc.bond_ramp_start = 10         # bond dimension of the first sweep
sc.bond_ramp_fraction = 0.5     # fraction of the sweeps spent ramping
sc.bond_ramp_noise_decay = 0.1  # noise decay per ramping sweep
sc.bond_ramp = False            # ...or restore a flat schedule at sc.maxm
\end{lstlisting}

The noise term (\lstinline|sc.noise|, White's density-matrix
perturbation, which keeps DMRG from stalling in a local minimum) is tied
to the same schedule: it starts at \lstinline|sc.noise|, decays by
\lstinline|sc.bond_ramp_noise_decay| per ramping sweep, and is switched
off entirely once the schedule reaches \lstinline|sc.maxm|, so the final,
converged sweeps are noise-free. With \lstinline|bond_ramp = False| the
original schedule is used instead --- full \lstinline|sc.noise| for the
first half of the sweeps, none for the second.

A \emph{warm} start is never truncated by the ramp: when
\lstinline|gs_energy()| is re-entered with a wavefunction already in hand
(\lstinline|set_initial_wf|, or simply a previous
\lstinline|gs_energy()| call's own solution), the ramp's starting bond
dimension is floored at whatever that state already carries, so a re-run
can only improve the energy. The ramp applies to the ground-state solve
on all three DMRG backends (\lstinline|itensor_version| 2, 3 and
\lstinline|"python"|); \lstinline|"julia_live"| keeps its own schedule.
See \lstinline|examples/groundstate/bond_dimension_ramp| for a 30-site
inhomogeneous Heisenberg--Hubbard chain timing ramp against flat.

\textbf{Targeting a quantum-number sector.} By default DMRG searches the
whole Hilbert space and returns the global ground state, at whatever
particle number or total magnetization that happens to have.
\lstinline|set_conserved_sector| instead confines the entire calculation
--- starting state, every sweep, the returned wavefunction --- to one
sector of a conserved quantity:

\begin{equation}
E_0(q)=\min_{|\psi\rangle\,:\,\hat Q|\psi\rangle=q|\psi\rangle}
\frac{\langle\psi|H|\psi\rangle}{\langle\psi|\psi\rangle},\qquad
[\hat Q,H]=0
\end{equation}

\begin{lstlisting}
fc.set_conserved_sector(Nf=6)        # exactly 6 particles
sc.set_conserved_sector(Sz=0)        # total Sz = 0
hc.set_conserved_sector(Nf=8, Sz=0)  # native Hubbard chain: both at once
fc.set_conserved_sector()            # no arguments -> back to the full space
\end{lstlisting}

Which quantities are available follows from the chain's site types:

\begin{center}
\footnotesize
\begin{tabular}{@{}p{0.55\textwidth}p{0.35\textwidth}@{}}
\toprule
Chain & Conserved quantities \\
\midrule
\texttt{Fermionic\_Chain} & \texttt{Nf} \\
\texttt{Spinful\_\allowbreak Fermionic\_Chain} (Jordan--Wigner) & \texttt{Nf}; also \texttt{Sz} under \lstinline|mode="ED"| \\
\texttt{Spinful\_\allowbreak Fermionic\_\allowbreak Chain\_\allowbreak Native} (native Hubbard sites) & \texttt{Nf}, \texttt{Sz}, or both \\
every spin chain (\texttt{Spin\_Chain}, any $S$) & \texttt{Sz} \\
\texttt{Bosonic\_Chain} & \texttt{Nb} \\
parafermion chains & none \\
\bottomrule
\end{tabular}
\end{center}

\lstinline|Sz| is in ITensor's integer $2S^z$ units, so
\lstinline|Sz=0| is $S^z_\mathrm{tot}=0$ and \lstinline|Sz=2| is
$S^z_\mathrm{tot}=1$. Asking for only part of what a site type offers
means exactly what it says: a native Hubbard chain with \lstinline|Nf|
alone fixes the particle number while leaving $S^z$ free, so spin-flip
terms stay legal.

This is the direct route to an addition spectrum $E_0(N)$, the charge gap
$E_0(N{+}1)+E_0(N{-}1)-2E_0(N)$, or a magnetization curve --- each point
a genuine ground-state energy of its own sector, rather than something
extracted by tuning a chemical potential or a field.

On \lstinline|itensor_version=3| its cost relative to the unconstrained
search depends on how big the solve is: the quantum numbers restore the
block sparsity the default mode gives up, but they also add per-block
bookkeeping, and only the former scales. Measured on Heisenberg chains
(BLAS pinned to one thread), sector mode is \textbf{0.6x} --- i.e.
slower --- at $n=20$, \lstinline|maxm=60|; \textbf{2.0x} faster at
$n=40$, \lstinline|maxm=100|; and \textbf{4.2x} faster at $n=60$,
\lstinline|maxm=200|, always for the same energy. Expect a small penalty
on toy systems and a real speedup on the ones that actually cost
something. On \lstinline|itensor_version="python"| the mechanism is
different --- that backend stores its tensors densely and uses the quantum
numbers only to \emph{label} basis states, so there is no block sparsity
to gain --- but a sector is not a tax either. Measured on the same 40-site
Heisenberg chain (\lstinline|maxm=100|, 15 sweeps, one P-core, BLAS pinned
to one thread, 3 repeats): \lstinline|itensor_version=3| takes 2.85\,s
dense and 2.05\,s in a sector, \lstinline|itensor_version="python"|
11.04\,s dense and 8.91\,s in a sector. So the sector run is ${\sim}1.2$x
\emph{faster} than the same backend's dense run, and ${\sim}4.3$x slower
than \lstinline|itensor_version=3|'s sector run --- about the same 4x
pure-Python-vs-compiled factor as everywhere else, i.e. the sector costs
nothing extra relative to it. The pure-Python speedup has a different cause
than v3's: confining the state lets it converge at a lower bond dimension
($\chi=52$ against $94$ for the same energy), which more than pays for the
charge penalty's extra MPO bond dimension. That balance depends on the
sweep count --- at 6 sweeps, before the sector run has settled to its lower
$\chi$, it is the slower of the two --- so treat ``about the same, give or
take 20\%'' as the honest summary rather than a reliable speedup. What a
sector buys on this backend is the targeting itself: the sector's own
ground state, which a dense search cannot reach at all.

Two consequences to be aware of, both reported rather than left to be
discovered. First, every operator built on the chain while a sector is
set must itself conserve the requested quantities --- a non-conserving
one raises a \lstinline|ValueError| naming it, whether it is in the
Hamiltonian, in a \lstinline|vev|, or a dynamical-correlator vertex:

\begin{lstlisting}
sc.set_conserved_sector(Sz=0)
sc.vev(sc.Sz[0]*sc.Sz[1])   # fine
sc.vev(sc.Sx[0])            # ValueError: ... changes ... by Sz=2 ...
\end{lstlisting}

A Hamiltonian written as $S^xS^x+S^yS^y+S^zS^z$ is fine even though no
single term of it conserves $S^z$: that expansion is recognized and its
$S^+S^+$/$S^-S^-$ strings cancel exactly. Second, the sector invalidates
the Hamiltonian and ground state already held by the backend, so the next
\lstinline|gs_energy()| re-solves from scratch.

Sector targeting is implemented by three solvers: DMRG on
\lstinline|itensor_version=3|, DMRG on
\lstinline|itensor_version="python"|, and \lstinline|mode="ED"| (see
below). DMRG on \lstinline|itensor_version=2| and on
\lstinline|"julia_live"| has no quantum numbers at all, and a sector-mode
chain deliberately raises rather than letting one of those answer, since a
solver without quantum numbers would silently return the \emph{global}
ground state instead.

The two DMRG backends implement it differently, which matters in one place.
\lstinline|itensor_version=3| confines the calculation structurally: its
tensors are block-sparse over quantum numbers, so an amplitude outside
the sector has nowhere to be stored. \lstinline|itensor_version="python"|
keeps dense storage and confines the calculation instead by adding a
charge penalty $\lambda\sum_k(\hat Q_k-q_k)^2$ to the operator its
\emph{variational} solves (ground state, excited states, band edges)
minimize. That penalty is identically zero on the target sector, so it
changes no reported number --- \lstinline|gs_energy()| reports $\langle
H\rangle$ under the plain Hamiltonian regardless --- and it is not
optional: a dense SVD mixes rows across charge blocks at the $10^{-16}$
level on every truncation, and a variational sweep amplifies that leak
toward whichever sector is lower in energy. With the penalty switched
off, an $n=12$ chain asked for $N_f=2$ with an attractive interaction
converges to the \emph{full} band instead; with it on, it reproduces
sector-restricted ED to $10^{-8}$. If a solve ever does end up outside
the requested sector, it raises rather than reporting the energy of the
wrong one. \lstinline|chain.set_sector_penalty(lam)| overrides
the default strength (derived from the Hamiltonian's own coefficient
scale) on that backend.

One accuracy caveat, on that backend only: \emph{excited} states inside a
sector are confined exactly (every one of them has the requested charge to
$10^{-8}$) but converge less accurately than on
\lstinline|itensor_version=3|. The overlap-penalty excited-state solver is
already the least accurate part of the pure-Python backend without any
sector, and the charge penalty widens the local effective Hamiltonian's
spectral range on top of that. Measured on a 6-site Heisenberg chain at
$S^z=0$, the second and third levels of the sector came out $0.2$--$0.4$
above their exact values, where \lstinline|itensor_version=3| reproduces
them exactly. Ground-state energies, expectation values, correlators and
time evolution are unaffected --- those match sector-restricted ED to
$10^{-8}$ on both backends. See
\lstinline|examples/backend_comparison/sector_v3_VS_python| for both
backends' addition spectra checked against sector-restricted ED, and for
the charge-penalty threshold measured directly.

What works under a sector: ground state and excited states, expectation
values and static correlators, entanglement entropies, real-time
evolution and dynamical correlators via the default
\lstinline|tevol_method="TDVP"| (and \lstinline|"TDVP_GSE"|) --- all
subject to the conserving-operator rule above. What does not, and says
so: METTS (\lstinline|metts_vev|,
\lstinline|metts_dynamical_correlator|), and the infinite-chain
algorithms (iDMRG, VUMPS). METTS averages over an ensemble sampled from
every sector at once, so confining the chain cannot restrict it; the
infinite-chain algorithms assemble tensors by hand in a way that is
dense-only under QN indices. They refuse rather than produce a wrong
number. \lstinline|tevol_method="TEBD"| is the one entry that differs by
backend: it refuses on \lstinline|itensor_version=3| (its gate assembly
sums bond gates before their fluxes agree, which aborts inside ITensor)
and works normally on \lstinline|itensor_version="python"|, whose gates
are dense and simply inherit the Hamiltonian's charge conservation. See
\lstinline|examples/groundstate/sector_targeted_groundstate| for the
addition spectrum and charge gap of a $t$--$V$ chain, checked against ED.

\textbf{Sector targeting under} \lstinline|mode="ED"|\textbf{.} ED
implements the same API, and by the simplest mechanism of the three: a
conserved charge is diagonal in ED's product basis, so a sector is a
\emph{set of basis states}, and confining a calculation to it means taking
the corresponding submatrix of every assembled operator. Same call, same
names, same $2S^z$ units, same answers --- checked sector by sector
against \lstinline|itensor_version=3| in
\lstinline|tests/test_sector_conservation_ed.py|.

\begin{lstlisting}
fc = fermionchain.Fermionic_Chain(8)
fc.mode = "ED"                   # this chain is solved by ED
fc.set_hamiltonian(h)
fc.set_conserved_sector(Nf=3)
fc.gs_energy(mode="ED")          # the three-particle ground state
fc.vev(sum(fc.N), mode="ED")     # exactly 3.0
\end{lstlisting}

Three things about it differ in kind rather than in result:

\begin{itemize}
\item \textbf{It restricts, it does not shrink.} The full Hilbert space is
  still assembled before the mask is applied, so a sector buys a smaller
  eigenproblem, not a smaller construction. What limits ED is the full
  space either way.
\item \textbf{An unreachable target can surface at the first calculation}
  rather than at \lstinline|set_conserved_sector|, because the ED object
  is built lazily. When the chain's DMRG backend carries the same quantum
  numbers it still validates the request immediately, as before; the lazy
  case is the ED-only quantum number below.
\item \textbf{\lstinline|promote_to_dense()| is DMRG-only.} It exists
  because a sector re-solve is expensive there; in ED it is not, so the
  way to leave a sector and keep working is
  \lstinline|set_conserved_sector()| and re-solve.
\end{itemize}

ED is also the only solver that can target a total-$S^z$ sector of
\texttt{Spinful\_\allowbreak Fermionic\_Chain}, the Jordan--Wigner
spinful chain: its DMRG representation is $2n$ \emph{spinless} fermionic
sites, which carry \lstinline|Nf| and know nothing about spin, while its
ED representation has an explicit up/down mode per orbital. Set
\lstinline|chain.mode = "ED"| before asking for it; a DMRG call on that
chain afterwards raises rather than answering with an \lstinline|Nf|-only
sector.

Because ED has quantum numbers now, a sector also no longer forbids the
automatic DMRG-to-ED fallbacks: a chain whose C++ extension was never
compiled, or an \lstinline|itensor_version=3| chain too short for
ITensor's two-site DMRG ($n<3$), answers the sector correctly through ED
instead of refusing.

The conserving-operator rule is the same under ED, and enforced for the
same reason: restricting an operator to the sector is exact for a static
expectation value but identically \emph{zero} for one that changes the
charge, so \lstinline|vev(fc.C[0])| inside a fixed-$N_f$ sector raises
instead of reporting a clean, wrong zero.
\lstinline|examples/groundstate/ed_sector_addition_spectrum| sweeps the
addition spectrum and charge gap of a $t$--$V$ chain this way.

\textbf{Leaving a sector without losing the state.} The
conserving-operator rule above rules out exactly the quantities one
usually wants \emph{from} a fixed-$N$ ground state:
$c_i|\mathrm{GS}\rangle$, a one-body density matrix built from it, a
photoemission weight, a pairing quench. Clearing the sector with
\lstinline|set_conserved_sector()| does not help on its own --- it also
throws the state away, so the next \lstinline|gs_energy()| re-solves
unconstrained and answers with the global ground state.
\lstinline|promote_to_dense()| is the other option: it leaves sector mode
while \emph{keeping} the state, converted exactly to its dense equivalent
on the chain's ordinary site indices.

\begin{lstlisting}
fc.set_conserved_sector(Nf=6)          # sweeps confined to 6 particles
fc.gs_energy()
fc.applyoperator(fc.C[3], fc.wf0)      # ValueError: c changes Nf by -1
fc.promote_to_dense()                  # same state, ordinary indices
wf = fc.applyoperator(fc.C[3], fc.wf0) # now legal
fc.vev(sum(fc.N), wf=wf)               # -> 5
\end{lstlisting}

Nothing is re-solved, truncated or approximated: a QN-conserving MPS is
the same wavefunction as its dense counterpart, only stored
block-sparsely, and promotion just scatters the blocks back into full
tensors. What it costs is the block sparsity itself, so promote
\emph{after} the expensive sweeps, not before. On
\lstinline|itensor_version="python"| the state was stored densely all
along and promotion only relabels its site indices, so there is nothing
to lose by promoting early --- but nothing to gain either. The
Hamiltonian and the
band-edge caches are rebuilt on the next call that needs them, while the
ground-state energy and wavefunction are kept --- a bare
\lstinline|gs_energy()| afterwards therefore still returns the sector's
energy rather than re-solving; call \lstinline|restart()| if an
unconstrained re-solve is what you want. A wavefunction Python is already
holding is not reached by \lstinline|promote_to_dense()| and needs
\lstinline|wf = fc.promote_mps(wf)| of its own.

Promotion always rebases onto the chain's \emph{own} site indices, kept
from construction, so states promoted out of different sectors at
different times remain comparable:

\begin{lstlisting}
fc.set_conserved_sector(Nf=n//2)   ; fc.gs_energy() ; fc.promote_to_dense()
wf_N = fc.wf0.copy()
fc.set_conserved_sector(Nf=n//2-1) ; fc.gs_energy() ; fc.promote_to_dense()
wf_Nm = fc.wf0.copy()
Z = abs(fc.overlap(wf_Nm, fc.applyoperator(fc.C[i], wf_N)))**2  # photoemission weight
\end{lstlisting}

\lstinline|promote_to_dense()| is available on
\lstinline|itensor_version=3| and \lstinline|itensor_version="python"|,
like \lstinline|set_conserved_sector| itself, and does nothing if no
sector is set. Handing a chain a wavefunction built under a different
sector setting raises instead of quietly contracting the wrong indices.
See
\lstinline|examples/groundstate/sector_promotion_to_dense| for the
$t$--$V$ chain's CDW profile, one-body density matrix and site-resolved
photoemission weight
$Z_i=|\langle \mathrm{GS}_{N-1}|c_i|\mathrm{GS}_N\rangle|^2$, all
computed this way and checked against sector-restricted ED.


\textbf{Expectation values and moments.} For any operator $O$ built the
same way as $H$,

\begin{equation}
\langle O\rangle=\langle\mathrm{GS}|O|\mathrm{GS}\rangle,\qquad
\langle O^n\rangle=\langle\mathrm{GS}|O^n|\mathrm{GS}\rangle
\end{equation}

\begin{lstlisting}
mz = [sc.vev(sc.Sz[i]).real for i in range(n)]        # local magnetization profile
e2 = sc.vev(h, npow=2)                                  # <H^2>, e.g. for fluctuations
\end{lstlisting}

\textbf{Energy fluctuation} (a measure of how sharply the DMRG/ED state
is an eigenstate, and physically the variance of $H$ in the prepared
state):

\begin{equation}
\delta E=\sqrt{\langle H^2\rangle-\langle H\rangle^2}
\end{equation}

\begin{lstlisting}
de = sc.gs_energy_fluctuation()
\end{lstlisting}

\textbf{Static two-point correlators.} \texttt{sc.vev(sc.Sz[0]*sc.Sz[i])}
gives $\langle S^z_0 S^z_i\rangle$ directly; \texttt{correlator.get\_correlator}
provides shorthand names for common correlators over a list of site
pairs, e.g.\ \texttt{"SS"} for the full dot product

\begin{equation}
\langle\mathbf S_i\cdot\mathbf S_j\rangle=\langle S_i^xS_j^x\rangle+\langle S_i^yS_j^y\rangle+\langle S_i^zS_j^z\rangle
\end{equation}

or fermionic correlators like \texttt{"cdc"} ($\langle c_i^\dagger
c_j\rangle$), \texttt{"density"}/\texttt{"densitydensity"} ($\langle
n_in_j\rangle$), and pairing correlators \texttt{"delta"}/
\texttt{"deltadeltad"}. These static correlators are the equal-time
limit of the dynamical correlators in Section~\ref{sec:dynamical}, and
are what you Fourier-transform to get an equal-time structure factor
$S(q)=\sum_{i,j}e^{iq(i-j)}\langle O_iO_j\rangle$.

\section{Excited states and energy gaps}
\label{sec:excited}

\textbf{Low-lying spectrum.}

\begin{lstlisting}
es = sc.get_excited(n=4)          # [E0, E1, E2, E3]
gap = sc.get_gap()                # E1 - E0
\end{lstlisting}

Both \texttt{mode="DMRG"} and \texttt{mode="ED"} return one entry per
eigen\emph{state}, so a degenerate level appears once per member in both and
the two agree index by index. \texttt{get\_gap()} is therefore $\sim 0$,
correctly, whenever the ground level is degenerate --- ask for enough levels
that the multiplet and the state you want both fit (with a three-fold
degenerate ground state, the first genuine excitation is \texttt{n=4}).
Degenerate members converge to slightly different energies under DMRG, so
when that spread approaches the splitting you are after, identify states by a
quantum number (e.g. $\langle S^2\rangle$) rather than by position in the
list.

\texttt{get\_excited\_states(n, purify=True)} additionally returns the
wavefunctions; \texttt{purify=True} re-diagonalizes $H$ in the
Gram-Schmidt orthogonalized subspace spanned by the raw excited-state
MPS, correcting for near-degenerate states that DMRG's iterative solvers
can otherwise mix.

Excited states beyond the ground state are found with an overlap-penalty
method (each state is optimized against $H+w\sum_k|\psi_k\rangle\langle\psi_k|$
for the states already found), which can occasionally converge to a
spurious, non-eigenstate stationary point instead of the true excited
state (\texttt{itensor\_version} 2, 3, and \texttt{"python"} are all
susceptible). Every call to \texttt{get\_excited\_states}/\texttt{get\_excited}
on a DMRG backend checks each returned state's energy fluctuation
$\langle H^2\rangle-\langle H\rangle^2$ against the ground state's own
(both already computed as a byproduct of the search) and emits a
\texttt{UserWarning} if a state's fluctuation is far above that
reference --- a cheap, always-on sanity check, though not a fix; a
warned-about state's energy and wavefunction should be treated with
caution (e.g.\ cross-checked against \texttt{mode="ED"} on a smaller
system, or recomputed with a larger \texttt{scale}/more sweeps).

\textbf{Sector (charge) gaps.} For a conserved quantity $A$ with
$[H,A]=0$ (e.g.\ total particle number $\hat N=\sum_i n_i$, or total
$S^z$), the gap to the lowest state with $\langle A\rangle$ shifted by
$d$ from the ground-state sector is obtained by adding a
Lagrange-multiplier penalty and increasing $\lambda$ until the
constraint is satisfied:

\begin{equation}
H_\lambda=H+\lambda\big(A-\langle A\rangle_0-d\big)^2,\qquad
\Delta_A(d)=E_\lambda-E_0
\end{equation}

\begin{lstlisting}
gap_charge = fc.get_charge_gap(d=2)   # gap to add/remove a pair of particles
\end{lstlisting}

This is exactly how a superconducting/charging gap is extracted in a
finite fermionic chain: $d=2$ probes the energy cost of adding a pair of
particles (Cooper-pair-like excitation) rather than a single particle,
which avoids odd/even-parity finite-size artifacts.

\textbf{Example: the Haldane gap.} For the spin-1 Heisenberg chain, the
famous Haldane gap between the (nearly four-fold degenerate, due to
fractionalized $S=\tfrac12$ edge states) ground-state manifold and the
first bulk excitation is

\begin{lstlisting}
es = sc.get_excited(n=6)
haldane_gap = es[4] - es[0]     # skip the 4 near-degenerate edge states
\end{lstlisting}

\textbf{Non-Hermitian Hamiltonians.} For $H\neq H^\dagger$ (complex
hopping amplitudes, gain/loss terms, PT-symmetric models, \ldots) the
``ground state'' convention throughout dmrgpy is the eigenvalue with the
smallest real part. Eigenvalues come in left/right eigenpairs,
\begin{equation*}
H\,|\psi_R\rangle=E\,|\psi_R\rangle,\qquad
H^\dagger|\psi_L\rangle=E^{*}|\psi_L\rangle,\qquad
\langle\psi_L|\psi_R\rangle=1 .
\end{equation*}
On every session backend (ITensor v3, ITensor v2, and the pure-Python
engine --- \lstinline|itensor_version| 2, 3 or \lstinline|"python"|),
\lstinline|gs_energy()| solves this with a genuine non-Hermitian DMRG
(NH-DMRG) --- a port of
ITensorNHDMRG.jl\footnote{\url{https://github.com/tipfom/ITensorNHDMRG.jl}}
in its default configuration: independent Arnoldi solves of the two-site
eigenproblem for $H$ and $H^\dagger$ (``onesided'' solver, targeting the
smallest real part), with both MPS truncated by the same isometry
obtained from the Hermitian average $\rho=(\rho_L+\rho_R)/2$ of the
left/right reduced density matrices (the ``fidelity'' algorithm of
Yamamoto \emph{et al.}, PRB \textbf{105}, 205125 (2022)). The full
eigenpair is available directly:

\begin{lstlisting}
e, psil, psir = sc.nhdmrg()     # E, left and right eigenvector MPS
sc.gs_energy()                  # same E; stores psir as the ground state
\end{lstlisting}

Because a non-Hermitian ``energy'' is not a variational bound,
\lstinline|nhdmrg()| certifies convergence through both eigen-residuals
$\lVert H|\psi_R\rangle-E|\psi_R\rangle\rVert$ and
$\lVert H^\dagger|\psi_L\rangle-E^{*}|\psi_L\rangle\rVert$ and re-runs
from a fresh random state (up to \lstinline|ntries| times) if either
stalls; \lstinline|krylovdim|/\lstinline|restarts| tune the per-bond
Arnoldi effort. The MPS Arnoldi route (\lstinline|get_excited_states|,
now used for non-Hermitian \emph{excited} states, $n\ge 2$) remains
available on every backend, but is typically several orders of magnitude
less accurate than NH-DMRG at comparable cost --- see
\texttt{examples/non\_hermitian/nhdmrg\_VS\_ED\_VS\_arnoldi}, which
cross-checks NH-DMRG on all three ported backends against exact
diagonalization and the Arnoldi route on an interacting fermionic chain
with a staggered imaginary potential.

\lstinline|itensor_version="julia_live"| also implements
\lstinline|nhdmrg()|, but is not a port: it calls the real
ITensorNHDMRG.jl package, with two adaptations applied in
\texttt{mpsjulialive/nhdmrg.jl} (see its header for the full derivation
of each). First, that package's ``adjoint'' sweep is against $H^{T}$
rather than $H^\dagger$, so its left vector solves the transpose
eigenvalue equation and is complex-conjugated here to match dmrgpy's
convention. Second, it does not tie its left solve to whichever
eigenvalue its right solve picked, so a spectrum with a complex-conjugate
pair tied for the smallest real part (the generic
PT-symmetric/Hatano-Nelson situation) could return a left and a right
vector belonging to \emph{different} eigenvalues; this is broken by
re-solving against $e^{i\theta}H$, which leaves every eigenvector
untouched and rotates the spectrum just enough to separate the pair's
real parts. The rotation can itself re-target a different eigenvalue when
$|\mathrm{Im}\,\lambda|$ is large, so the untied run's eigenvalue is kept
as an anchor and a tie-break result is accepted only if it reproduces it;
otherwise the attempt is failed rather than answered. Neither adaptation is observable on a
complex-\emph{symmetric} $H$, which is what most textbook non-Hermitian
models (and every other non-Hermitian example in this repository) happen
to be --- see
\texttt{examples/non\_hermitian/nhdmrg\_julia\_asymmetric\_VS\_ED} for a
chain with asymmetric hopping, where both are.

The biorthogonal pair $|\psi_R\rangle,|\psi_L\rangle$
is also what feeds the non-Hermitian dynamical correlator,
\lstinline|get_dynamical_correlator(submode="KPM")| for $H\neq H^\dagger$
--- see Section~\ref{sec:dynamical}.

Two independent MPS Arnoldi implementations are available, both
matrix-free (they only ever apply $H$ to a wavefunction, never build a
matrix) and both usable in ED mode or DMRG mode on any backend.
\lstinline|dmrgpy.mpsalgebra.lowest_energy_non_hermitian_iram|
(\texttt{algebra/arpacktk.py}) is an Implicitly Restarted Arnoldi Method
(IRAM), adapted from
ARPACK\footnote{\url{https://bitbucket.org/chaoyang2013/arpack}}'s
\lstinline|znaupd|/\lstinline|znaup2|/\lstinline|znaitr|/\lstinline|znapps|
--- the same exact-shift polynomial-filter restart ARPACK's own
\lstinline|eigs|-style solvers use, re-derived here for dmrgpy's own
MPS/ED wavefunction objects instead of flat arrays (no BLAS/LAPACK
Fortran dependency). It is the \textbf{default} MPS Arnoldi solver
(\lstinline|dmrgpy.mpsalgebra.lowest_energy_non_hermitian|, and the
route \lstinline|get_excited_states| uses for non-Hermitian $H$ with
$n\ge2$, \lstinline|excited_states_non_hermitian| in
\texttt{excited.py}), since it compresses and reuses its existing Krylov
subspace instead of rebuilding it from scratch on every restart, needing
noticeably fewer $H|\psi\rangle$ applications to reach the same
tolerance on most spectra.
\lstinline|dmrgpy.mpsalgebra.lowest_energy_non_hermitian_arnoldi|
(\texttt{algebra/arnolditk.py}, a restarted Arnoldi method with explicit
Rayleigh-Ritz reseeded restarts) remains available directly for
comparison --- the two can trade places on hard (near-degenerate)
spectra, see \texttt{examples/non\_hermitian/arnoldi\_vs\_iram\_benchmark},
which benchmarks both head to head (Op-count, wall time, accuracy) in
both ED and DMRG mode --- and is still used internally for Krylov
operators with a projector baked in (\texttt{fermionchain.py}'s
\lstinline|Spinon_Chain.get_gs|, which IRAM's interface doesn't (yet)
support).

\texttt{algebra/arpacktk.py} also implements ARPACK's shift-invert mode
(\lstinline|dmrgpy.mpsalgebra.mpsiram_shift_invert|, ARPACK's mode 3,
$\mathrm{OP}=(A-\sigma I)^{-1}$), used to find the eigenvalues closest
to a target energy $e$ rather than an extremal one ---
\texttt{degeneracy.py}'s \lstinline|eigenvalue_degeneracy| (used by
\lstinline|gs_degeneracy(mode="DMRG")| on non-Hermitian Hamiltonians) is
built on it. Since dmrgpy has no exact MPO inverse ---
\lstinline|self.applyinverse| is itself only an iterative
correction-vector solve --- this departs from ARPACK's own mode-3
assumption that $\mathrm{OP}$'s eigenvalues are \emph{exactly}
$1/(\lambda-\sigma)$: only the IRAM restart \emph{direction} (which
Krylov directions the implicit shifts filter away) comes from
$\mathrm{OP}$'s own cheap Hessenberg matrix (\lstinline|which="LM"|,
since the largest $|\mathrm{OP}\text{-eigenvalue}|$ is the
$H$-eigenvalue closest to $e$); convergence and the reported eigenpairs
are instead recomputed every outer iteration from $H$'s own exact
(still cheap, $O(\mathrm{ncv}^2)$) representation on the Krylov basis
$\mathrm{OP}$ builds --- the same accommodation arnolditk's own
\lstinline|mode="ShiftInv"| path already made.

\texttt{algebra/arpacktk.py} also implements ARPACK's mode 2 --- the
generalized eigenproblem $A|\psi\rangle=\lambda M|\psi\rangle$ for a
Hermitian positive-definite $M$
(\lstinline|dmrgpy.mpsalgebra.mpsiram_generalized|/
\lstinline|generalized_excited_states|, $\mathrm{OP}=M^{-1}A$, $B=M$).
Unlike mode 1/3, the Krylov basis must be orthonormal in the
$M$-weighted inner product $\langle u,v\rangle_M=\langle u|M|v\rangle$
rather than the plain one, so building it uses a separate routine
(\lstinline|arnoldi_extend_generalized|) rather than adding a
conditional $M$ to the mode-1/3 one. $\mathrm{OP}$ again goes through
\lstinline|self.applyinverse| (the same approximate-inverse caveat as
mode 3, here with the trivial shift $\sigma=0$) --- but unlike mode 3,
$\mathrm{OP}$'s own Hessenberg-matrix eigenvalues stay directly
meaningful (there is no shift to undo), so convergence/selection follow
the same pattern as plain \lstinline|mpsiram|, just built with the
$M$-inner product. Verified against \lstinline|scipy.linalg.eigh|'s
generalized Hermitian-definite eigensolver in both ED and DMRG mode.

\textbf{Generalized-eigenvalue DMRG.} For the same problem
$H|\psi\rangle=\lambda A|\psi\rangle$ ($A$ Hermitian positive definite),
\lstinline|gs_energy_generalized(A)| solves it with a genuine DMRG sweep
instead of a Krylov method, needing no approximate operator inverse at
all (unlike \lstinline|mpsiram_generalized| above, which must invert $M$
iteratively since dmrgpy has no exact MPO inverse):

\begin{lstlisting}
lam = fc.gs_energy_generalized(a_metric_operator)   # smallest lambda
\end{lstlisting}

The trick is a self-consistent Lagrange multiplier: minimizing
$\langle\psi|H|\psi\rangle$ subject to the metric normalization
$\langle\psi|A|\psi\rangle=1$ has stationarity condition
$(H-\lambda A)|\psi\rangle=\mu|\psi\rangle$ for multiplier $\lambda$ ---
the \emph{ordinary} ($\mu$,$\psi$) eigenproblem of the plain-normalized
shifted operator $H-\lambda A$, exactly what a standard two-site DMRG
sweep already finds. At $\mu=0$ this is precisely
$H|\psi\rangle=\lambda A|\psi\rangle$, so each outer iteration (i) builds
the MPO $H-\lambda A$ from the current $\lambda$ estimate, (ii) runs one
ordinary DMRG sweep against it, then (iii) updates $\lambda$ to the
freshly-swept state's generalized Rayleigh quotient
$\langle\psi|H|\psi\rangle/\langle\psi|A|\psi\rangle$ --- one outer
iteration per \lstinline|Sweeps| schedule entry, so bond dimension ramps
exactly as an ordinary \lstinline|gs_energy()| run's own schedule does.
$A=\mathrm{Id}$ reduces this exactly to plain ground-state DMRG.

Implemented on the pure-Python (pyitensor) backend
(\lstinline|itensor_version="python"|, i.e. after
\lstinline|chain.setup_python()|) ---
\lstinline|dmrgpy.pyitensor.dmrg.dmrg_generalized| is the underlying
routine ---, on compiled ITensor v3 (\lstinline|itensor_version=3|,
i.e. after \lstinline|chain.setup_cpp(version=3)|), where
\lstinline|Chain::gs_energy_generalized|
(\texttt{mpscpp3/chain\_session.h}) runs the identical algorithm against
ITensor v3's own \lstinline|dmrg()|/\lstinline|Sweeps|/\lstinline|sum()|
instead of pyitensor's hand-rolled two-site sweep, and on the live Julia
backend (\lstinline|itensor_version="julia_live"|, i.e. after
\lstinline|chain.setup_julia()|), where \lstinline|get_gs_generalized|
(\texttt{mpsjulialive/generalized.jl}) runs it once more against
ITensorMPS.jl's own
\lstinline|dmrg()|/\lstinline|Sweeps|/\lstinline|add()|.
\lstinline|itensor_version=2| (mpscpp2) doesn't
have this session method yet and raises \lstinline|NotImplementedError|.
See \texttt{examples/groundstate/dmrg\_generalized\_benchmark}, which
heads all three solvers up against each other on the same
interacting-fermion-chain test problem: needing no approximate inverse,
both DMRG routes are far more accurate than ARPACK mode 2 (whose own
correction-vector solve also gets more expensive, and less accurate, as
the chain grows --- two orders of magnitude slower and $10^{9}$ times
less accurate on an 8-site chain), and v3 is itself consistently
$\sim2$--$4\times$ faster than pyitensor at these sizes.
\lstinline|julia_live| matches both DMRG routes to machine precision and,
once warm, lands between them: at $n=8$, v3 0.16\,s,
\lstinline|julia_live| 0.20\,s, pyitensor 0.50\,s. Its \emph{first} call
in a session is dominated by Julia's JIT compilation ($\sim46$\,s), so a
single-point timing of this backend measures compilation, not the solver
--- read its second and later rows.

\textbf{Non-Hermitian generalized-eigenvalue DMRG.}
\lstinline|gs_energy_generalized| also accepts a non-Hermitian $H$:
\lstinline|self.hamiltonian| is checked for Hermiticity, and a
non-Hermitian one is transparently dispatched to a dedicated solver
(\texttt{nhdmrg.py}'s \lstinline|nhdmrg_generalized|) instead of raising
--- same calling convention, same "smallest real part" convention as
non-Hermitian \lstinline|gs_energy()|/NH-DMRG (\S4) above:

\begin{lstlisting}
lam = fc.gs_energy_generalized(a_metric_operator)   # complex lambda, smallest Re
\end{lstlisting}

The metric $A$ must still be Hermitian positive definite (mirroring
\lstinline|mpsiram_generalized|'s own $M$ precondition --- and, following
ARPACK's own convention there, the \emph{primary} operator needs no
Hermiticity assumption at all, only the metric does). The
self-consistent Lagrange-multiplier trick generalizes directly: since
$A$ is Hermitian, $(\lambda A)^\dagger=\bar\lambda A$, so minimizing (in
the NH-DMRG sense --- smallest real part, not a variational bound)
subject to the metric \emph{biorthogonal} normalization
$\langle\psi_L|A|\psi_R\rangle=1$ gives stationarity condition
$(H-\lambda A)|\psi_R\rangle=\mu|\psi_R\rangle$,
$(H-\lambda A)^\dagger|\psi_L\rangle=\bar\mu|\psi_L\rangle$ --- the
\emph{ordinary} ($\langle\psi_L|\psi_R\rangle=1$-normalized) NH-DMRG
eigenproblem of the shifted pair
$(H-\lambda A,\,H^\dagger-\bar\lambda A)$, exactly what one ordinary
NH-DMRG sweep already finds. At $\mu=0$ this is precisely
$H|\psi_R\rangle=\lambda A|\psi_R\rangle$, so each outer iteration (i)
builds $H-\lambda A$ and $H^\dagger-\bar\lambda A$ from the current
(complex) $\lambda$ estimate, (ii) runs one ordinary NH-DMRG sweep
against them, then (iii) updates $\lambda$ to the freshly-swept pair's
generalized \emph{biorthogonal} Rayleigh quotient
$\langle\psi_L|H|\psi_R\rangle/\langle\psi_L|A|\psi_R\rangle$ in place of
the Hermitian case's plain one. $A=\mathrm{Id}$ reduces this exactly to
plain NH-DMRG.

Implemented on the pure-Python (pyitensor) backend
(\lstinline|dmrgpy.pyitensor.nhdmrg.nhdmrg_generalized|), on compiled
ITensor v3 (\lstinline|Chain::nhdmrg_generalized|,
\texttt{mpscpp3/chain\_session.h}, a line-for-line port against this
file's own hand-rolled two-site sweep) and on the live Julia backend
(\lstinline|get_gs_generalized_nhdmrg|,
\texttt{mpsjulialive/generalized.jl}, the same outer loop wrapped around
real ITensorNHDMRG.jl sweeps) --- \texttt{mpscpp2}
(\lstinline|itensor_version=2|) still raises
\lstinline|NotImplementedError| for a non-Hermitian $H$, no analogous
session method there. Since \lstinline|nhdmrg_generalized|'s two-site
sweep never calls ITensor v3's own \lstinline|dmrg()| (it is hand-rolled
directly against a restarted Arnoldi solve and manual ITensor
contractions, unlike the Hermitian path above), it does \emph{not} need
the Hermitian path's short-chain guard: chains shorter than 3 sites work
fine here on \lstinline|itensor_version=3|. See
\texttt{examples/non\_hermitian/nhdmrg\_generalized\_benchmark}, which
heads all three routes up against each other on a non-Hermitian
interacting-fermion-chain test problem (a staggered imaginary on-site
potential): ARPACK mode 2 needs no adaptation at all for a non-Hermitian
primary operator (its $M$-positive-definite precondition was always the
only one), so this is a fair like-for-like comparison, and the same
pattern as the Hermitian benchmark holds --- needing no approximate
inverse, both DMRG routes are far more accurate and (as the chain grows)
up to two orders of magnitude faster than ARPACK mode 2. \lstinline|julia_live| also matches to machine precision here, sitting
just ahead of pyitensor once warm ($n=8$: v3 1.16\,s,
\lstinline|julia_live| 2.26\,s, pyitensor 2.42\,s) with the same
first-call JIT caveat as the Hermitian benchmark. Unlike the
Hermitian case, v3 isn't consistently faster than pyitensor here (roughly
on par at the sizes benchmarked) --- NH-DMRG's per-bond cost already pays
for \emph{two} Arnoldi solves (right block and its adjoint) regardless of
backend, narrowing the compiled-vs-pure-Python gap relative to plain
ground-state DMRG's single local diagonalization per bond.

\textbf{Caveat (both Hermitian and non-Hermitian
\lstinline|gs_energy_generalized|).} Afterward, \lstinline|fc.wf0|/
\lstinline|fc.e0| hold the eigenstate/eigenvalue of the
\emph{generalized} problem, not a plain eigenstate of
\lstinline|fc.hamiltonian| alone --- every other method that reads
\lstinline|fc.wf0| as an ordinary ground state
(\lstinline|get_excited_states()|, any dynamical/KPM correlator, \ldots)
has no way to tell the difference and will silently build on the wrong
reference state if called afterward. Call
\lstinline|gs_energy_generalized()| as the last step of a calculation, or
recompute a genuine ground state (\lstinline|gs_energy()|/
\lstinline|nhdmrg()|) first if you need one of those other methods too.

\textbf{Effective low-energy Hamiltonians.} Given the $n$ lowest
eigenstates $\{|\psi_k\rangle\}$ and the projector
$P=\sum_k|\psi_k\rangle\langle\psi_k|$ onto the manifold they span, the
projected Hamiltonian $PHP$ can be fitted as a real linear combination of
the same projections of a chosen set of operators $O_a$:

\begin{equation}
PHP\simeq\sum_a J_a\,PO_aP,\qquad
J_a=\arg\min_{J}\Big\|PHP-\sum_a J_aPO_aP\Big\|
\end{equation}

and the fitted $J_a$ are the effective couplings.

\begin{lstlisting}
# fit against an explicit list of operators -> one coefficient each
J = fc.get_heff(operators=[fc.Sx[0]*fc.Sx[1], fc.Sy[0]*fc.Sy[1],
                           fc.Sz[0]*fc.Sz[1]], n=4, mode="ED")

# ...or against all pairwise products of the chain's own Sx/Sy/Sz
from dmrgpy import effectivehamiltonian
coef = effectivehamiltonian.get_effective_hamiltonian_couplings(fc, n=4)
latex = sc.get_effective_hamiltonian(n=1)   # the same fit, as a latex string
\end{lstlisting}

\lstinline|n| selects the manifold, \lstinline|tol| is the magnitude
below which a fitted coupling is dropped, and \lstinline|operators|
overrides the default basis (\lstinline|get_projection_operators|, every
$S^x/S^y/S^z$ of the chain). \lstinline|method="single"| builds
$PO_aPO_bP$ --- the product of the two \emph{projected} operators, one
projection per operator --- while \lstinline|method="full"| builds
$P(O_aO_b)P$, the projection of the product, at one projection per
\emph{pair}. These are the same operator only when the manifold is closed
under $O_a$; they are not two implementations of one thing.

Two ways this calculation can be meaningless, both of which are now
refused rather than returned:

\begin{itemize}
  \item \textbf{$n$ must not cut a degenerate multiplet.} The retained
    states are then not a symmetry-invariant subspace and the fitted
    couplings are physically meaningless (on a Hubbard dimer, whose low
    manifold is a singlet plus a three-fold triplet, \lstinline|n=3| gave
    couplings ${\sim}50\times$ the correct ones). A \lstinline|ValueError|
    names the degeneracy and the values of $n$ that do cut cleanly.
  \item \textbf{The operator basis must be linearly independent on that
    manifold.} Otherwise the coefficients can be shifted by any
    null-space vector without changing the fit, so no unique set of
    couplings exists; the raw candidate set is typically very much
    overcomplete (37 candidates spanning a rank-16 space on the dimer). A
    \lstinline|ValueError| reports the rank and the number of null
    directions.
\end{itemize}

The fit itself is a linear least-squares solve, so it is exact and
reproducible; the eigenstates are L\"owdin-orthonormalized first when
they need it, which DMRG's overlap-penalty excited states generally do.
See \lstinline|examples/groundstate/effective_hamiltonian|, which
recovers the exact $J\cos2\phi$/$J\sin2\phi$ twist of a Hubbard dimer's
exchange under a Peierls phase $\phi$.


\section{Entanglement and quantum information}

\textbf{Entanglement entropy of a real-space bipartition.} Cutting the
chain at bond $(i,i+1)$, the Schmidt decomposition of the ground state
is $|\mathrm{GS}\rangle=\sum_\alpha\lambda_\alpha|\alpha\rangle_L\otimes
|\alpha\rangle_R$ and the von Neumann entanglement entropy of either
half is

\begin{equation}
S_i=-\sum_\alpha\lambda_\alpha^2\log\lambda_\alpha^2=-\mathrm{Tr}\big[\rho_L\log\rho_L\big],\qquad
\rho_L=\mathrm{Tr}_R|\mathrm{GS}\rangle\langle\mathrm{GS}|
\end{equation}

\begin{lstlisting}
s = wf.get_bond_entropy(i, i+1)
\end{lstlisting}

\texttt{get\_site\_entropy(i)}/\texttt{get\_pair\_entropy(i,j)} instead
build a reduced density matrix from local \emph{projectors} (e.g.\
$S^z=\pm\tfrac12$ for spin-$\tfrac12$, occupation number for fermions)
rather than a full MPS bond cut, which lets you ask about the
entanglement of a single site or a pair of (possibly non-adjacent)
sites with the rest of the system.

\textbf{Mutual information} between two sites (or subsystems) $i,j$:

\begin{equation}
I(i,j)=S_i+S_j-S_{ij}
\end{equation}

\begin{lstlisting}
mi = wf.get_mutual_information(1, 2)
\end{lstlisting}

quantifies total (classical + quantum) correlation between $i$ and $j$,
and is often used to map out effective couplings or emergent degrees of
freedom (e.g.\ between the two edge spins of a Haldane chain).

\textbf{CFT central charge.} For a critical (gapless) chain, the
entanglement entropy of a cut at position $\ell$ in an open chain of
length $L$ follows the Calabrese-Cardy formula:

\begin{equation}
S(\ell)=\frac{c}{6}\log\!\left[\frac{2L}{\pi a}\sin\frac{\pi\ell}{L}\right]+\text{const.}
\end{equation}

DMRGPY computes $S(\ell)$ at every bond and least-squares fits this
formula to extract the central charge $c$ --- the universal number that
identifies the underlying conformal field theory (e.g.\ $c=\tfrac12$
for the critical Ising chain, $c=1$ for a free boson / XX chain):

\begin{lstlisting}
c = wf.get_CFT_central_charge()
\end{lstlisting}

\textbf{Single-particle correlation matrix and orbital entanglement.}
For fermionic systems, the correlation matrix

\begin{equation}
C_{ij}=\langle c_i^\dagger c_j\rangle
\end{equation}

is directly accessible (\texttt{sc.get\_correlation\_matrix()}), and
diagonalizing it, $C=U\,n\,U^\dagger$, gives natural-orbital occupations
$n_\alpha\in[0,1]$ (\texttt{get\_correlated\_orbitals}) and an
orbital-resolved entanglement entropy

\begin{equation}
S=-\sum_\alpha\Big[n_\alpha\log n_\alpha+(1-n_\alpha)\log(1-n_\alpha)\Big]
\end{equation}

(\texttt{get\_correlation\_entropy}). \texttt{get\_highorder\_correlation\_matrix}
and \texttt{get\_four\_correlation\_tensor} give the corresponding
two-particle correlators $\langle c_i^\dagger c_j^\dagger c_l
c_k\rangle$ and $\langle c_i^\dagger c_j c_k^\dagger c_l\rangle$.

\texttt{get\_four\_correlation\_tensor(ctmode=...)} has three
implementations: \texttt{ctmode="explicit"} (backend-agnostic Python loop
of \texttt{vev()}s, always available), \texttt{ctmode="full"} (native
per-element AutoMPO build -- C++ for \texttt{itensor\_version} \texttt{2}/
\texttt{3}, pure Python for \texttt{"python"} -- builds and applies an
independent AutoMPO per $(i,j,k,l)$ tuple), and \texttt{ctmode="sweep"}
(\texttt{itensor\_version} \texttt{3} or \texttt{"python"}, plain
non-native-spinful fermionic sites only) -- a single-sweep,
environment-reuse implementation following the algorithmic idea of
\href{https://github.com/ITensor/ITensorCorrelators.jl}{ITensorCorrelators.jl}:
rather than rebuilding a fresh MPO for every tuple, it reuses partial
tensor-network contractions across the whole $(N,N,N,N)$ tensor. Agrees
with \texttt{ctmode="full"} and ED to machine precision / solver
tolerance on both backends, and is substantially faster: at $n=12$,
9.3\,s $\to$ 1.2\,s under \texttt{itensor\_version=3} and 37.7\,s $\to$
2.0\,s under \texttt{"python"} (measured single-threaded; see the
BLAS-threads section), i.e.\ roughly $8\times$ and $19\times$, against a
\texttt{ctmode="full"} that is slower still.

Those numbers reflect a fix worth knowing about if you read the older
docstrings. The tensor splits into pairwise-distinct-index entries (handled
by the sweep) and repeated-index entries, and the latter used to fall back
to the same per-tuple AutoMPO build \texttt{ctmode="full"} uses. There are
only $O(n^3)$ of those against the sweep's $O(n^4)$, which is why they were
described as subdominant -- but a smaller count times a far more expensive
per-tuple cost still dominates, and measured, that fallback was 96\% of the
whole runtime at $n=12$. Every one of those operators is a product of four
$C^\dagger/C$ factors on at most 3 distinct sites, i.e.\ \emph{local}, so
both backends now fold them directly instead of compiling an MPO over the
whole chain. \texttt{accelerate} (default \texttt{True}, accepted on every
\texttt{ctmode}) still only affects the repeated-index entries for
\texttt{ctmode="sweep"} --- the pairwise-distinct sweep has no equivalent
conjugate-pair saving to skip (see either backend's own docstring) -- so
don't expect the usual $\sim2\times$ win \texttt{accelerate} gives
\texttt{ctmode="full"}. See
\texttt{examples/staticcorrelators/four\_correlation\_tensor\_sweep\_VS\_full},
\texttt{Chain::four\_correlation\_tensor\_sweep}
(\texttt{mpscpp3/chain\_session.h}) and \texttt{pyitensor/chain.py}'s port
for the algorithm.

\texttt{ctmode="fold"} (\texttt{itensor\_version="python"} or \texttt{3}) evaluates
every tuple as a \emph{local} operator fold -- no MPO is built for any of
them. It is flavour-agnostic: it reads each mode's (operator name, site) off
the chain's own \texttt{C}/\texttt{Cdag}, so it covers spinless chains and
\texttt{Spinful\_Fermionic\_Chain\_Native} alike. That matters because native
spinful sites previously had only \texttt{ctmode="explicit"} under
\texttt{"python"}, which builds an MPO and sweeps the whole chain per tuple;
\texttt{fold} is exact against it to machine precision and 8--12$\times$
faster (5 sites / 10 modes: 9.8\,s $\to$ 0.9\,s), and is now the default
there. The C++ backend has the same port, where it replaces a
\texttt{ctmode="full"} that was itself \emph{slower} than the pure-Python
fold: 2.8\,s $\to$ 0.7\,s at 5 sites, a $4\times$ win and an
algorithm-beats-language result worth remembering. It does not reuse environments \emph{across} tuples the way
\texttt{ctmode="sweep"} does, so for a plain spinless chain prefer
\texttt{"sweep"}.

\texttt{ctmode="batched"} (\texttt{itensor\_version="python"}, spinless and
native spinful alike) computes the same tensor with every transfer
contraction \emph{batched} over the tuples that share it, and is the default
under \texttt{"python"} since it is exact against both \texttt{"sweep"} and
\texttt{"fold"} to machine precision and 15--28$\times$ faster. The
reorganization is worth understanding, because it fixes two separate
problems at once. \texttt{"sweep"} and \texttt{"fold"} both evaluate one
tuple's worth of MPS transfer at a time -- a chain of $\chi\times\chi$
contractions -- and there are $O(n^4)$ of them, each far too small to keep
BLAS busy; and the \emph{repeated}-index tuples (those whose four indices
are not pairwise distinct) get no environment reuse at all under
\texttt{"sweep"}, so despite being only $O(n^3)$ of them they measured at
61--65\% of its total runtime at $n=12\ldots20$. \texttt{"batched"} writes
each tuple in site-sorted order as a sequence of at most four (local matrix,
parity) steps; two tuples agreeing on their first $r$ steps share a partial
environment whatever sites those steps landed on, so the whole $(N,N,N,N)$
tensor collapses onto a trie of a few dozen matrix sequences, each holding
one array of environments batched over the site combinations that realize
it. One site of the sweep is then two large GEMMs per trie node instead of
tens of thousands of small contractions, and distinct and repeated tuples go
through the same machinery. Measured single-threaded on a spinless chain at
\texttt{maxm=20}, full tensor: $n=16$, 10.65\,s (\texttt{"sweep"}) /
6.69\,s (ITensor v3, C++) / 0.38\,s (\texttt{"batched"}); $n=30$, 117.3\,s /
72.4\,s / 7.4\,s. For native spinful the comparison is against
\texttt{"fold"}: 1.29\,s $\to$ 0.22\,s at 5 sites and 3.02\,s $\to$ 0.31\,s
at 6, the crossover sitting at 4 sites (below that the trie build is a fixed
cost the tensor is too small to amortize). The batch axis is also what makes
this calculation worth putting on a GPU at all, and it is the one place in
this library where the device wins \emph{below} the usual $\chi\sim120$--$160$
crossover -- because the arithmetic per array operation comes from the tuple
batch rather than from bond dimension. Measured on an H200 against one Xeon
core at $n=30$ (warm, i.e.\ XLA's ${\sim}22$\,s of per-shape compilation
already paid): 0.97/0.95/0.98\,s at \texttt{maxm} $=$ 20/40/80 against
1.89/9.44/22.81\,s on the host, i.e.\ $2.0\times$, $9.9\times$ and
$23.3\times$ -- the device time does not move at all across the range, so
every one of those is a lower bound. Setting
\texttt{backend.set\_backend("jax")} immediately before the call is enough;
the ground state itself can stay on NumPy, since the MPS is converted once.
Full tables and the cold-run caveat in
\texttt{docs/gpu\_cpu\_performance.md}.

\texttt{ctmode=None} (the default) auto-selects the fastest method
actually available for the wavefunction's backend/chain type:
\texttt{"batched"} whenever it applies (\texttt{itensor\_version="python"},
either flavour), else \texttt{"sweep"} whenever it applies, else \texttt{"full"} whenever it
applies (any \texttt{itensor\_version} in \texttt{2}/\texttt{3}/
\texttt{"python"}, or \texttt{Spinful\_Fermionic\_Chain\_Native} under
\texttt{itensor\_version=3} via its own
\texttt{four\_correlation\_tensor\_spinful()}), else the always-correct
\texttt{"explicit"} fallback (e.g.\ for
\texttt{itensor\_version="julia\_live"}, which has no
\texttt{"full"}/\texttt{"sweep"} implementation). Passing a
\texttt{ctmode} explicitly is still a hard request -- it raises rather
than silently falling back if that method is not available for the
wavefunction at hand. \texttt{Spinful\_Fermionic\_Chain\_Native} always
needs \texttt{ctmode="full"} (or the default, which resolves to it): the
sweep has no native-spinful-site counterpart yet.

\section{Dynamical (frequency-dependent) correlators}
\label{sec:dynamical}

The central quantity is a retarded correlator (Green's function)
between two operators $A,B$,

\begin{equation}
G_{AB}(\omega)=\langle\mathrm{GS}|A\,\frac{1}{\omega-H+E_0+i\delta}\,B|\mathrm{GS}\rangle,\qquad
S_{AB}(\omega)=-\frac1\pi\,\mathrm{Im}\,G_{AB}(\omega)
\end{equation}

where $\delta$ is a small broadening. Choosing $A=B=S^z_i$ at the same
site gives the local dynamical spin structure factor
$S^{zz}_{ii}(\omega)$ (what a local probe like NMR/ESR couples to);
choosing $A=S^z_i$, $B=S^z_j$ at different sites and Fourier-transforming
over $i-j$ gives the momentum-resolved dynamical structure factor
$S(q,\omega)$ measured in inelastic neutron scattering. All five
submodes below compute $G_{AB}(\omega)$ (or equivalently
$S_{AB}(\omega)$); they differ only in \emph{how}, and therefore in what
energy range/resolution/cost trade-off they offer:

\begin{lstlisting}
(x, y) = sc.get_dynamical_correlator(mode="DMRG", submode="KPM",
                                      name=(sc.Sz[0], sc.Sz[0]))
\end{lstlisting}

\textbf{What \texttt{name=} accepts.} Either a documented correlator
string together with the site indices it applies to
(\texttt{name="ZZ", i=0, j=1}, and likewise
\texttt{"XX"}/\texttt{"YY"}/\texttt{"+-"}/\texttt{"cdc"}/
\texttt{"densitydensity"}/\ldots), or an explicit pair $(A,B)$ of
operators. The pair may be \texttt{MultiOperator}s, as above, or the
\emph{already-built} operators \texttt{toMPO()} returns
(see \emph{Algebra on already-built operators} above): a
\texttt{StaticOperator} under
\texttt{mode="DMRG"}, an \texttt{EDOperator} under \texttt{mode="ED"}.
Products of those are accepted too, which is the natural way to write a
composite channel such as
$c^\dagger_{0\uparrow}c^{\phantom\dagger}_{1\uparrow}$:

\begin{lstlisting}
Cdagup = [fc.toMPO(A, mode=mode) for A in fc.Cdagup]
Cup    = [fc.toMPO(A, mode=mode) for A in fc.Cup]
A01 = Cdagup[0]*Cup[1]
(x, y) = fc.get_dynamical_correlator(name=(A01, A01.get_dagger()), mode=mode)
\end{lstlisting}

Handing over an already-built operator saves rebuilding it on every
call, which is what makes it worth doing across a frequency sweep or a
parameter scan. It works for \texttt{mode="ED"} (every submode) and for
\texttt{submode="EX"}/\texttt{"CVM"}/\texttt{"ROOTN"} under
\texttt{mode="DMRG"} --- the paths that consume an operator by
\emph{applying} it. The remaining
submodes (\texttt{"KPM"}, \texttt{"TD"}, \texttt{"TDZ"},
\texttt{"SECTOR"}, the METTS finite-temperature correlator of
\S\ref{sec:thermal}, and \texttt{cvm\_solver="variational"})
instead rebuild the operator inside the backend from its symbolic term
list, so they need the \texttt{MultiOperator} itself and raise a
\texttt{TypeError} saying so if given \texttt{toMPO()} output.

\textbf{\texttt{submode="KPM"} --- Kernel Polynomial Method.} $H$ is
linearly rescaled into $\tilde H=(H-b)/a$ so its full spectrum lies in
$[-1,1]$ (\texttt{kpm\_scale}, default 0.7, sets how much margin is left
inside $[-1,1]$), then $G_{AB}$'s spectral function is expanded in
Chebyshev polynomials $T_m(\tilde H)$:

\begin{equation}
\mu_m=\langle\mathrm{GS}|A\,T_m(\tilde H)\,B|\mathrm{GS}\rangle,\qquad
S_{AB}(\tilde\omega)=\frac{1}{\pi\sqrt{1-\tilde\omega^2}}\left[\mu_0+2\sum_{m=1}^{N-1}g_m\,\mu_m\,T_m(\tilde\omega)\right]
\end{equation}

with Jackson-kernel damping coefficients $g_m$ suppressing Gibbs
ringing from the truncation at $N$ moments (\texttt{kpm\_n\_scale}
scales $N$, i.e.\ the energy resolution, relative to the rescaled
bandwidth). This is the default, general-purpose method: robust, works
across the whole spectrum at once, cost grows only linearly with the
number of moments.
The band edges entering the rescaling are obtained variationally: the
lower edge reuses the ground-state energy, and the upper edge runs a
deliberately reduced-effort DMRG on $-H$ (few sweeps at modest bond
dimension) --- it is only a spectral \emph{bound}, protected by the
\texttt{kpm\_scale} margin, not a physical result, and a variational
underestimate only shrinks the number of moments. If the bound is ever
too tight for the chosen \texttt{kpm\_scale}, the moment recursion
detects the resulting exponential divergence and aborts with an
explicit error rather than returning a silently wrong spectrum.

\textbf{Reconstruction kernel (\texttt{kernel=}) --- how the moments
become a spectrum.} The formula above is only one way to turn
$\{\mu_m\}$ into $S_{AB}$. \texttt{kernel="jackson"} (the default) is the
classic choice; \texttt{"lorentz"} and \texttt{"plain"} are the usual
alternatives. Every one of them damps the moments, $\mu_m\to
g_m\mu_m$, which is equivalent to convolving the exact spectrum with a
positive kernel of width $\sim\pi/N$ --- and \emph{any} such kernel is
only second-order accurate, so the error at a smooth point decays as
$O(N^{-2})$ no matter how many moments are spent.

\texttt{kernel="hodc"} selects the high-order delta-Chebyshev
reconstruction of Yi, Massatt, Horning, Luskin, Pixley \& Kaye,
\href{https://arxiv.org/abs/2512.03149}{arXiv:2512.03149} (2025), which
breaks that second-order barrier. Instead of damping the moments it
replaces $\delta$ itself by the order-$m$ rational regularization

\begin{equation}
K_\eta(E,x)=-\frac1\pi\sum_{l=1}^{m}\mathrm{Im}\left[\frac{w_l}{E-x+\eta z_l}\right],\qquad
z_l=x_l+i,\qquad \sum_{l}w_l z_l^{\,j}=\delta_{j0}\ \ (j<m)
\end{equation}

--- $m$ complex-weighted Lorentzians of width $\eta$, whose weights are
fixed by a Vandermonde moment-matching system that annihilates the first
$m-1$ terms of the small-$\eta$ expansion, leaving $K_\eta\to\delta$ at
$O(\eta^m)$ rather than $O(\eta)$. Expanding $K_\eta(E,\cdot)$ in
Chebyshev polynomials gives energy-\emph{dependent} coefficients
$\nu_m(E,\eta)$ and the reconstruction

\begin{equation}
S_{AB}(\tilde\omega)=\sum_{m=0}^{N-1}\nu_m(\tilde\omega,\eta)\,\mu_m
\end{equation}

from the \textbf{same, undamped} moments. Nothing in the expensive half
of the calculation changes --- no extra Chebyshev recursion, no backend
change, and \texttt{kernel=} is available on every backend that supports
\texttt{submode="KPM"}. \texttt{hodc\_order} is $m$ (default 6, the
paper's own choice; 1--8 are allowed, and $m=1$ degenerates to plain
Lorentzian broadening), and \texttt{hodc\_eta} is $\eta$ in the same
energy units as \texttt{delta}/\texttt{es}, defaulting to
\texttt{delta} itself --- i.e.\ ``you asked for resolution
\texttt{delta}, you get resolution \texttt{delta}, with an
$O(\delta^m)$ error instead of an $O(\delta^2)$ one''. Two caveats, both
inherited from the method rather than from this implementation:

\begin{itemize}
\item The kernel is \textbf{not positive} for $m>2$, so the
  reconstructed spectrum can dip slightly negative next to a sharp
  feature. That is the price of the higher order.
\item The gain is an asymptotic statement about \emph{smooth} points. It
  is real where the spectrum is a continuum and disappears at a van Hove
  singularity or a band edge, and it disappears entirely for a spectrum
  made of isolated $\delta$-peaks, which is not smooth at any
  resolution.
  \texttt{examples/dynamical\_correlator/hodc\_VS\_jackson\_kernel}
  measures both regimes against an exactly solvable case.
\item \textbf{It is more sensitive to MPS truncation noise than
  Jackson}, and this one is specific to using it inside DMRG rather than
  to the method. The Chebyshev recursion is only conditionally stable
  under truncation: the error in $\mu_m$ grows roughly exponentially in
  $m$, and a damping kernel's $g_m\to0$ as $m\to N$ happens to suppress
  exactly the worst moments. HODC takes them undamped --- its only
  high-$m$ suppression is the $\nu_m\sim q^m\sim e^{-m\eta}$ decay of its
  own coefficients. So at modest \texttt{kpmmaxm} the useful moment count
  is capped by moment accuracy well before it is capped by cost, and if
  an HODC spectrum looks \emph{worse} than the Jackson one the first
  thing to try is a larger \texttt{hodc\_eta} (more damping), not more
  moments. Measured in that example, on an XX chain at
  \texttt{kpmmaxm=32} with $N=1200$ moments whose error has grown to
  $O(10^2)$ by the last of them: at the default $N\eta\approx3.3$ HODC is
  more than 10x \emph{worse} than Jackson, and at $N\eta\approx14$ it is
  several times better (4--8x, drifting between DMRG runs). On exact
  moments the same sweep puts the optimum right at the default, which is
  where it belongs. The same example measures the moment error directly against
  exact moments, which is the cleanest way to find where that horizon
  sits for a given chain.
\end{itemize}

Because the kernel only enters after the moments exist,
\texttt{sc.get\_dynamical\_correlator\_moments(...)} returns the raw
$(\mu,E_\mathrm{min},E_\mathrm{max},a,N,\delta)$ and
\texttt{kpmdmrg.dynamical\_correlator\_from\_moments(...)} reconstructs
from them, so several kernels can be compared without repeating the DMRG
work.

\textbf{Energy truncation (\texttt{itensor\_version="python"} and \texttt{3}) --- narrowing the KPM
window below the full bandwidth.} The default \texttt{kpm\_scale}
rescales $H$ so its \emph{entire} many-body spectrum fits in $[-1,1]$,
which is always safe but wastes resolution: a local operator's
correlator usually has real spectral weight only over a width $W_A$
much smaller than the full bandwidth $W$, so most of the Chebyshev
expansion's $N$ moments are ``spent'' on frequency regions the
correlator never visits. Choosing a narrower window (smaller
\texttt{kpm\_scale}) concentrates the same $N$ onto the physically
relevant region instead --- but without a safeguard, the moment
recursion's own numerical noise lets a little high-energy weight leak
in, and since Chebyshev polynomials are unbounded outside $[-1,1]$, that
leakage grows exponentially over the recursion (exactly the failure the
\texttt{kpm\_scale}-too-tight error above is detecting, just
deliberately triggered by choosing a small window on purpose rather than
by an inaccurate band-edge estimate). \emph{Energy truncation} (Holzner,
Weichselbaum, McCulloch \& von Delft, ``Chebyshev matrix product state
approach for spectral functions'', \href{https://doi.org/10.1103/PhysRevB.83.195115}{PRB 83, 195115 (2011)}, Sec.
III-B) is the fix: after every Chebyshev vector is formed, it is swept
site by site, and at each site a small ($\le$ \texttt{kpm\_truncate\_dK})
Krylov subspace of that site's local effective Hamiltonian is
diagonalized and any component with rescaled energy
$|\varepsilon|\ge$~\texttt{kpm\_truncate\_threshold} is projected out,
before the vector is used for a moment or fed into the next recursion
step. This is a \emph{precautionary} measure, not an exact projector (a
finite Krylov dimension only approximates the local spectrum), so a
small residual above threshold is expected and does not vanish with more
\texttt{kpm\_truncate\_nsweeps} --- matching the original paper's own
characterization of it (its Sec.~V.C).

\begin{lstlisting}
sc.kpm_scale = 0.3            # narrower than the safe ~0.5 floor
sc.kpm_energy_truncate = True # enable energy truncation (default: False)
sc.kpm_truncate_dK = 30       # per-site Krylov subspace dimension
sc.kpm_truncate_nsweeps = 10  # number of truncation sweeps per vector
sc.kpm_truncate_threshold = 1.0  # rescaled-energy cutoff (paper's eps_P)
(x, y) = sc.get_dynamical_correlator(mode="DMRG", submode="KPM", name=(sc.Sz[0], sc.Sz[0]))
\end{lstlisting}

\texttt{kpm\_truncate\_dK} is a \emph{convergence} knob, not a tuning
one: the projection it defines is exact once $d_K$ reaches the dimension
of a site's local space, so raise it until the reconstructed spectrum
stops moving and then stop (larger $d_K$ only costs time, since the
per-site Krylov build is where energy truncation's whole overhead lives).
The paper's own Table~I recommends $d_K=30$. Both backends build that
subspace with full re-orthogonalization; a single Gram--Schmidt pass is
not enough at $d_K\sim 30$, where the Krylov vectors converge toward the
site's locally dominant eigendirection and a singly-orthogonalized basis
silently ceases to be orthonormal (which used to corrupt pyitensor's
projection and move a 6-site Heisenberg chain's spectral peak from
$0.48\,J$ to $1.06\,J$).

Enabling \texttt{kpm\_energy\_truncate} also switches the rescaling
convention itself: instead of centering the window on the full
bandwidth's \emph{midpoint}, it anchors it at the ground state $E_0$
(the paper's own Eq.~21b), placing $E_0$ at one edge of the window and
$E_0+W_s$ ($W_s=(E_{\max}-E_{\min})\cdot$\texttt{kpm\_scale}) at the
other. This matters because a dynamical correlator's Chebyshev vectors
are built by acting $A$/$B$ on the ground state, so the physically
relevant region sits just above $E_0$, not around the spectrum's
geometric middle --- the midpoint-centered window would otherwise clip
the ground state itself out of the window before ever reaching a
genuinely useful, narrower regime. Available for
\texttt{itensor\_version="python"} and \texttt{itensor\_version=3}; there
is no \texttt{itensor\_version=2} port (mpscpp2 has no equivalent
machinery to build a per-site local effective Hamiltonian from, unlike
mpscpp3's \texttt{LocalMPO}/\texttt{diagHermitian}). Requesting it on
\texttt{itensor\_version=2} is silently a no-op (the existing,
always-safe rescaling is used instead). The v3 port
(\texttt{mpscpp3/chain\_session.h}'s
\texttt{kpm\_dynamical\_correlator\_truncated()}) is a wholly independent
method from \texttt{kpm\_dynamical\_correlator()} --- a deliberate design
choice so the existing, always-safe v3 KPM path is never touched by this
feature --- not a branch inside the same function the way the pyitensor
port is; both implement the identical algorithm and agree on the same
physical answer (see \texttt{test\_kpm\_energy\_truncation\_v3\_accuracy.py}'s
cross-backend check). Available only for the $(A,B)$-operator dynamical
correlator, not the lower-level ``arbitrary operator'' KPM
(\texttt{general\_kpm}/\texttt{kpm\_wfa\_wfb}), which has no ground-state
reference to anchor to --- setting the flag has no effect on those, on
any backend (they expand an already-bounded operator rescaled into
$[a,b]$ by \texttt{scale\_operator()}, which needs no safeguard). See
\texttt{examples/dynamical\_correlator/dynamical\_correlator\_kpm\_energy\_truncation}
for a worked example (including the divergence this fixes, reproduced on
purpose), \texttt{src/dmrgpy/pyitensor/kpm\_energy\_truncation.py} for
the pyitensor implementation, and \texttt{mpscpp3/chain\_session.h}'s
own \texttt{kpm\_dynamical\_correlator\_truncated}/
\texttt{kpm\_energy\_truncate} comments for the v3 one.

\textbf{Performance note:} energy truncation is a controlled
\emph{slowdown}, not a speedup --- it buys resolution/feasibility (a
window that would otherwise diverge), not raw speed. Measured directly
(a 4-site chain, \texttt{kpm\_scale} narrowed from the safe 0.7 to
0.65): the pyitensor backend costs $\sim$5.5--12.5$\times$ more per
Chebyshev moment (\texttt{kpm\_truncate\_dK}/\texttt{nsweeps} of 10/3 vs.
the paper's own recommended 30/10, respectively) than the untruncated
path, and v3's native port $\sim$15--32$\times$ more per moment for the
same settings on this small system --- the per-site Krylov-subspace cost
scales with bond dimension, so this overhead grows quickly with system
size (measured directly at 8 sites: $\sim$245$\times$ per moment at
\texttt{kpm\_truncate\_dK=30}/\texttt{nsweeps=10}). Narrowing the window
does reduce the moment count needed for a given frequency resolution,
but that reduction only outweighs the per-moment cost once the
correlator's own spectral width is genuinely much smaller than the full
bandwidth (the paper's own premise for large, extensive systems) --- for
a small test system the two effects roughly cancel or net out to a
slowdown, not a speedup.

\textbf{\texttt{submode="CVM"} --- correction-vector method.} Instead of
a global polynomial expansion, this solves directly for the correction
vector at one frequency $\omega$ at a time, via the positive-definite
linear system

\begin{equation}
\big[(H-\omega-E_0)^2+\eta^2\big]\,x_c=-\eta\,B|\mathrm{GS}\rangle
\end{equation}

(solved by conjugate gradient in MPS form), from which

\begin{equation}
x=(\omega+E_0+i\eta-H)^{-1}B|\mathrm{GS}\rangle=i\,x_c+\frac{H-\omega-E_0}{\eta}\,x_c,\qquad
G_{AB}(\omega)=\langle\mathrm{GS}|A|x\rangle
\end{equation}

Here $\eta$ (\texttt{delta}) is the artificial broadening that
regularizes the resolvent at a real frequency. CVM is more accurate at a
single targeted frequency/energy window (e.g.\ zooming in on a sharp
resonance) than a global KPM expansion, at the cost of re-solving the
linear system for every $\omega$ on the requested grid. Two things keep
that per-$\omega$ cost down. First, everything $\omega$-independent is
computed once per sweep instead of once per point --- in particular the
right-hand side $-\eta B|\mathrm{GS}\rangle$. Second, the CG loop
terminates early once the bond-dimension truncation
(\texttt{cvm\_maxm}) puts a floor under the reachable residual ---
past that floor the iteration cannot improve the tracked best solution
any further (the truncated recurrence in fact diverges), so the solver
returns the best correction vector seen instead of burning the full
\texttt{cvm\_nit} iteration budget (\texttt{cvm\_patience} sets how
many iterations without meaningful improvement conclude the floor is
reached, and \texttt{cvm\_blowup} how far past the best residual the
running one may diverge before stopping). Neither feature changes the
answer.
Each point reports its CG iteration count and best residual; if that
residual stalls far above \texttt{cvm\_tol}, the correction vector is
not converged at this bond dimension and the fix is a larger
\texttt{cvm\_maxm}, not more iterations. (Warm-starting each point's CG
from the neighboring point's correction vector was tried and measured
to \emph{hurt} --- truncated CG from a nearby-but-wrong start can
stagnate at a much worse residual than from the cold start --- so each
point is solved independently.)

\textbf{\texttt{submode="ROOTN"} --- root-$N$ Krylov-space correction
vector.} Implements Nocera \& Alvarez, ``Root-$N$ Krylov-space
correction-vectors for spectral functions with the density matrix
renormalization group''
(\href{https://arxiv.org/abs/2204.03165}{arXiv:2204.03165}). Rather than
building the correction vector
$x(\omega+i\eta)=(\omega-H+E_0+i\eta)^{-1}B|\mathrm{GS}\rangle$ in one
shot, it is built as $N$ sequential fractional-power steps,

\begin{equation}
x^{p/N}(\omega+i\eta)=\Big(\frac{1}{\omega-H+E_0+i\eta}\Big)^{p/N}B|\mathrm{GS}\rangle,\qquad p=1,\dots,N,
\end{equation}

where each step re-seeds a Krylov (Lanczos) subspace of dimension
\texttt{nkry} with the \emph{previous} step's vector, tridiagonalizes $H$
in that subspace, and applies the $1/N$-power resolvent in the resulting
eigenbasis before handing the result forward as the next step's seed.
$N=1$ reduces to the ``conventional'' Krylov-space correction vector
(Nocera, PRE 2016) that the paper compares against. In the paper's
MPS/DMRG setting, \texttt{nkry}'s role is played by the bond dimension
$m$, and building the correction vector in $N$ smaller steps lets the
entanglement grow gradually instead of needing a single very large $m$ at
high target frequency; here, \texttt{nkry} directly caps the Lanczos
subspace size at each step, and the same qualitative effect is observed:
at a fixed, small \texttt{nkry}, increasing $N$ reduces the error against
the exact answer, especially near the top of the many-body bandwidth.
With \texttt{nkry} equal to the full Hilbert space dimension, root-$N$
reproduces the exact answer for any $N$, since the Krylov subspace is
then exact --- a useful self-consistency check of the recursion itself,
independent of the Krylov truncation.

\begin{lstlisting}
(x, y) = sc.get_dynamical_correlator(mode="ED", submode="ROOTN",
                                      name=(sc.Sz[0], sc.Sz[0]),
                                      N=8, nkry=20)
(x, y) = sc.get_dynamical_correlator(mode="DMRG", submode="ROOTN",
                                      name=(sc.Sz[0], sc.Sz[0]),
                                      N=8, nkry=20) # itensor_version in (2,3,"python")
\end{lstlisting}

Two implementations exist, both cross-checked to agree at machine
precision on small chains (see
\texttt{examples/dynamical\_correlator/dynamical\_correlator\_rootn\_ED}
and
\texttt{examples/dynamical\_correlator/dynamical\_correlator\_rootn\_v3\_VS\_ED}):

\begin{itemize}
\item \texttt{mode="ED"} (\texttt{src/dmrgpy/algebra/rootn.py}): the
recursion above against the exact ED Hamiltonian, with the Krylov
subspace built from plain numpy vectors.
\item \texttt{mode="DMRG"} (\texttt{src/dmrgpy/rootndmrg.py},
\texttt{itensor\_version} in \texttt{(2,3,"python")}): the \emph{same}
recursion, but with the Krylov subspace built out of \emph{global} MPS
vectors --- each Lanczos step is a truncated MPO application
(\texttt{self.toMPO(self.hamiltonian)*v}) rather than a local, per-bond
update. This is deliberately \emph{not} how the paper's own Appendix
algorithm works (a multi-target state-averaged DMRG sweep, jointly
representing $|\mathrm{GS}\rangle$, $B|\mathrm{GS}\rangle$,
$\mathrm{Re}(x)$, $\mathrm{Im}(x)$ in one block MPS compressed together at
every bond, so the correction vector's local tensor is built directly
from the ground state's own environment): a first attempt at a more
literal, per-bond local-sweep implementation (in
\texttt{mpscpp3/chain\_session.h}, using \texttt{LocalMPO} and a
Lanczos-based fractional-power update at each bond, modeled on
\texttt{applyExp}/TDVP) gave numerically \emph{wrong} results
(sign-flipping, unstable) when cross-checked against exact ED ---
reapplying the local update bond-by-bond on a single self-referential MPS
does not correctly realize ``apply $f(H)$ once'' globally, unlike TDVP (a
well-defined local Trotter step) or ground-state DMRG (repeated local
energy minimization provably converges to the global minimum). The
global-Krylov approach here avoids that failure mode entirely by only
ever using already-tested whole-MPS primitives (truncated MPO
application, inner products, MPS addition) --- the same primitives
\texttt{submode="CVM"}'s conjugate gradient already relies on --- at the
cost of not reproducing the paper's own bond-dimension bookkeeping across
the four channels, and of every Lanczos step costing a full MPO
application rather than a cheap local tensor contraction.
\end{itemize}

\textbf{\texttt{submode="TD"} --- time-dependent DMRG.} Real-time
evolution gives the correlator directly in the time domain,

\begin{equation}
C(t)=\langle\mathrm{GS}|A(t)B(0)|\mathrm{GS}\rangle=e^{iE_0t}\langle\mathrm{GS}|A\,e^{-iHt}\,B|\mathrm{GS}\rangle
\end{equation}

which is then windowed with a damping/taper factor $w_\delta(t)$ (see
\texttt{damping} below; the default is an exponential
$w_\delta(t)=e^{-\delta t}$, equivalent to a Lorentzian broadening of
width $\delta$ in frequency) and Fourier transformed,

\begin{equation}
S_{AB}(\omega)=\frac1\pi\,\mathrm{Re}\!\int_0^{T}\!dt\;e^{i\omega t}\,C(t)\,w_\delta(t)
\end{equation}

The total simulated time $T$ (\texttt{damping\_periods}/$\delta$) must
be long enough that the damping has suppressed truncation ringing by
$t=T$. Frequency resolution is set by $T$ (via the usual $\Delta\omega
\sim1/T$ time-frequency uncertainty), so this method is best when you
want fine resolution over a \emph{narrow} frequency window, at the cost
of a real dynamical simulation (bond dimension grows with entanglement
generated during the evolution). The propagator used for this evolution
is the same \texttt{sc.tevol\_method} selector described in
\S\ref{sec:quenches}, including \texttt{tevol\_method="TEBD"} ---
cheaper per step than TDVP whenever the Hamiltonian is strictly
nearest-neighbor, e.g.\ a \texttt{Spinful\_Fermionic\_Chain\_Native}
Hubbard chain (the standard interleaved
\texttt{Spinful\_Fermionic\_Chain} is \emph{not} nearest-neighbor after
Jordan-Wigner threading, so TEBD raises \texttt{NotImplementedError}
there). A script that wants the \texttt{"TEBD"} speedup on whichever of
the two representations turns out nearest-neighbor, without hardcoding
that choice per model, can use \texttt{tevol\_method="AUTO"} instead
(\S\ref{sec:quenches}) to fall back to \texttt{"TDVP"} automatically on
the interleaved representation rather than raising.

\emph{Sharpening \texttt{"TD"}'s lineshape (\texttt{damping},
\texttt{predict}).} An exponential window gives $S_{AB}(\omega)$ an
exact Lorentzian lineshape, whose $1/(\omega-\omega_0)^2$ tail is much
heavier away from a peak than \texttt{"KPM"}'s default Jackson-kernel
reconstruction (see
\texttt{docs/td\_dynamical\_correlator\_sharpening\_plan.md} for the
full design, literature, and the empirical comparison behind the
default below). Two independent, composable knobs address this:

\begin{verbatim}
(x, y) = sc.get_dynamical_correlator(mode="DMRG", submode="TD",
        name=(sc.Sz[0], sc.Sz[0]), delta=0.05,
        damping="exp",                # default; "gaussian"/"parzen" also available
        predict=True, lp_order=None, lp_extend_factor=10)  # predict=True is the default
\end{verbatim}

\begin{itemize}
\item \texttt{damping} selects $w_\delta(t)$: \texttt{"exp"} (default),
  the exponential above; \texttt{"gaussian"},
  $w_\delta(t)=e^{-(\delta t)^2/2}$, whose Fourier transform decays as
  $e^{-\omega^2}$ (far faster than the Lorentzian's algebraic tail), at
  the cost of a slightly wider FWHM at the same $\delta$
  ($2\sqrt{2\ln2}\,\delta\approx2.35\delta$ vs the Lorentzian's
  $2\delta$); \texttt{"parzen"}, a taper with compact support that
  vanishes (with zero derivative) exactly at $t=T$, which instead
  targets the Gibbs ringing from truncating $C(t)$ at a finite $T$,
  independent of the peak-broadening tradeoff above.
\item \texttt{predict} (default \texttt{True}) extrapolates the raw,
  undamped $C(t)$ via linear prediction
  (\texttt{dynamicstk.linearprediction.linear\_predict\_extend},
  following White \& Affleck and Barthel, White \& Schollw\"ock,
  arXiv:0901.2342) before any windowing: an autoregressive model of
  order \texttt{lp\_order} (default \texttt{None}, auto-picked as
  $\min(20,\max(4,\lfloor n_t/10\rfloor))$ so it stays safe even for a
  short simulation) is fit to the tail of $C(t)$ and used to synthesize
  \texttt{lp\_extend\_factor} times as many additional samples, so the
  same real TDVP/TEBD simulation yields a sharper spectral function than
  windowing alone --- the standard fix in the DMRG-dynamics literature
  for finite-simulated-time resolution loss, since it needs no
  additional entanglement growth. Pass \texttt{predict=False} to
  recover the exact pre-existing behavior.
  \texttt{lp\_fit\_start\_fraction} (default 0.5) skips the
  corresponding leading fraction of $C(t)$ before fitting, and
  \texttt{lp\_max\_pole\_radius} (default 1.0) reflects any fitted pole
  outside the unit circle back onto it, since the physical correlator
  only has weight for poles on or inside it and a slightly unstable fit
  would otherwise diverge under extrapolation.
\end{itemize}

\texttt{damping="exp"} paired with \texttt{predict=True} is the default
combination because it empirically gave the narrowest, best-centered
peak among every combination tried on a test system with a known exact
gap: pairing \texttt{predict=True} with \texttt{"gaussian"} instead came
out \emph{worse} (its larger intrinsic FWHM at fixed $\delta$ partly
cancels prediction's own narrowing), and pairing it with
\texttt{"parzen"} came out statistically tied with plain \texttt{"exp"},
not better
(\texttt{docs/td\_dynamical\_correlator\_sharpening\_plan.md} has the
numbers). Note that neither knob changes the \emph{asymptotic} algebraic
decay of the tail far from any resonance --- \texttt{"KPM"}'s far tail
still decays much faster, since its Jackson-kernel-damped Chebyshev
reconstruction has no Lorentzian/algebraic tail to begin with; these
knobs sharpen \texttt{"TD"}'s peaks and suppress its near-tail ringing,
not close that specific gap.

Both \texttt{damping} and \texttt{predict}/\texttt{lp\_*} are also
accepted by \texttt{submode="TDZ"} and by the internal $S(k,\omega)$
reduction (\texttt{sxt\_to\_skomega}), since all three share the same
windowing/FFT tail --- but their own defaults are unchanged
(\texttt{damping="exp"}, \texttt{predict=False}), since the empirical
comparison above was only run for \texttt{"TD"}.

\textbf{\texttt{submode="TDZ"} --- complex-time evolution (Cao, Lu,
Stoudenmire \& Parcollet, arXiv:2311.10909).} Real-time evolution
(\texttt{"TD"} above) grows entanglement, so the MPS bond dimension
needed for a given accuracy grows with the simulated time $T$. TDZ
instead evolves along a complex time contour

\begin{equation}
z(t,\alpha_0)=\int_0^t e^{-i\alpha_0 f(t')}\,dt',\qquad f(t)=e^{-t\omega_0},\qquad \omega_0=2\pi/t_{\max}
\end{equation}

Since $\mathrm{Im}\,z(t,\alpha_0)<0$ for $\alpha_0>0$, this
progressively damps high-energy content as it evolves, so the bond
dimension needed for a given accuracy grows far more slowly than under
real-time evolution alone (the original paper reports $\chi\sim20$--$30$
vs $\chi\sim500$--$700$ for comparable accuracy on the Anderson impurity
model). The true real-time ($\alpha_0=0$) correlator is then recovered
order by order via a perturbative Taylor expansion in $\alpha_0$ around
the simulated contour,

\begin{equation}
C(t,0)\approx\phi^{(0)}(t,\alpha_0)+\sum_{n=1}^{n_{\max}}g^{(n)}(t,\alpha_0)
\end{equation}

where $\phi^{(n)}(t,\alpha_0)=\langle H^n B|\mathrm{GS}\rangle\cdot
|\psi(t,\alpha_0)\rangle$ (precomputed once per $n$, reused as a fixed
overlap target at every time step) and $g^{(n)}$ are explicit
combinatorial expressions in $\phi^{(1..n)}$ and the pure contour
integrals $J^{(n)}(t,\alpha_0)=-i\,\partial^n_{\alpha_0}z(t,\alpha_0)$
(see the paper's Appendix B; this implementation hardcodes $n\le4$,
which the paper finds always suffices for $\alpha_0\lesssim0.3$). The
reconstructed $C(t,0)$ is then windowed/Fourier-transformed exactly as
in \texttt{"TD"}.

\begin{verbatim}
(x, y) = sc.get_dynamical_correlator(mode="DMRG", submode="TDZ",
                                      name=(sc.Sz[0], sc.Sz[0]),
                                      alpha0=0.1, n_max=4, dt=0.05)
\end{verbatim}

\texttt{alpha0} is the contour angle parameter (larger reduces the bond
dimension needed further, but requires a larger \texttt{n\_max} to
reconstruct the real axis accurately); \texttt{n\_max} ($\le4$) is the
reconstruction order; \texttt{dt}/\texttt{tmax}/\texttt{nt} set the
underlying time step/duration exactly as in \texttt{"TD"}. Uses
two-site TDVP when available (\texttt{itensor\_version} 3 or
\texttt{"python"}, \texttt{tevol\_method="TDVP"} or \texttt{"AUTO"} ---
\texttt{"AUTO"} reduces to plain \texttt{"TDVP"} here rather than trying
\texttt{"TEBD"} first, since a complex/imaginary time step has no TEBD
counterpart on any backend to try, see \S\ref{sec:quenches}'s own note
--- the paper's own setup), one-site TDVP with global subspace expansion
(\texttt{tevol\_method="TDVP\_GSE"}, see \S\ref{sec:quenches} --- same
\texttt{itensor\_version} support as \texttt{"TDVP"}), or falls back to
the MPO-Taylor propagator otherwise (\texttt{tevol\_method="MPO"} or
\texttt{"TEBD"}, or \texttt{itensor\_version=2}, which has no TDVP) ---
the same TDVP-vs-Taylor choice \texttt{"TD"} already makes. Current
scope: only the ``greater'' branch of the correlator is computed (the
same
simplification \texttt{"TD"} itself already makes), so this is best
used the same way as \texttt{"TD"}: high-resolution work in a narrow
frequency window, now reachable at a lower bond-dimension cost for a
given simulated time.

\textbf{\texttt{submode="EX"} --- exact diagonalization in a truncated
DMRG subspace.} Builds $A$, $B$, $H$ explicitly in the subspace spanned
by the lowest \texttt{nex} DMRG excited states, then
evaluates the exact Lehmann sum in that subspace,

\begin{equation}
G_{AB}(\omega)=\sum_{n=1}^{n_{ex}}\frac{\langle\mathrm{GS}|A|n\rangle\langle n|B|\mathrm{GS}\rangle}{\omega-E_n+i\delta}
\end{equation}

i.e.\ a small, explicit sum over poles at the computed excited-state
energies $E_n$, each with residue given by the transition matrix
elements. Cheap and exact \emph{within} the truncated subspace; only as
good as how many/which excited states were computed. The \texttt{nex}
excited-state MPS are not assumed to be orthonormal (\texttt{dcex.py}
builds their own overlap matrix and solves a generalized eigenvalue
problem $Hc=eSc$ rather than assuming $S=1$), which matters in practice
since Gram-Schmidt over bond-truncated MPS is itself only approximate.

\textbf{\texttt{submode="SECTOR"} --- the Lehmann sum over one
quantum-number sector.} The same explicit pole sum as \texttt{"EX"}, but
the intermediate states are not excited states of the whole Hilbert
space: they are the lowest eigenstates of the \emph{one
conserved-quantum-number sector that $B|\mathrm{GS}\rangle$ actually
lands in}. That sector is not a choice. $B$ shifts every conserved charge
by a fixed amount, so $\langle n|B|\mathrm{GS}\rangle$ vanishes
identically unless $|n\rangle$ carries the ground state's charges plus
$\mathrm{charge}(B)$:

\begin{center}
\begin{tabular}{lll}
\hline
correlator & ground state in & intermediate states in \\
\hline
$(c_i,c_j^\dagger)$ & $N$ particles & $N+1$ \\
$(c_j^\dagger,c_i)$ & $N$ particles & $N-1$ \\
$(S^-_i,S^+_j)$ & $S_z$ & $S_z+1$ (\texttt{Sz}$+2$ in $2S_z$ units) \\
$(S^+_i,S^-_j)$ & $S_z$ & $S_z-1$ \\
$(S^z_i,S^z_j)$ & $S_z$ & $S_z$ \\
\hline
\end{tabular}
\end{center}

\begin{lstlisting}
fc.set_conserved_sector(Nf=6)          # solve at exactly 6 particles
(x, y) = fc.get_dynamical_correlator(submode="SECTOR",
                                     name=(fc.C[0], fc.Cdag[0]), nex=20)
\end{lstlisting}

Each intermediate state is therefore a separately converged DMRG
eigenstate of a \emph{smaller} Hilbert space. Symmetry --- not the
overlap penalty \texttt{"EX"} has to rely on --- is what keeps them out
of the ground state and out of each other's charge channels: they live in
a \emph{different} sector, so the search cannot collapse back into
$|\mathrm{GS}\rangle$. Within the target sector they are still found by
that same penalty method, only in a much smaller and better-conditioned
space, and the usual per-state fluctuation warning still applies. There
is no
broadening/resolution tradeoff, no expansion order, no time window and no
Fourier artifact: the poles come out at their converged energies with
their exact weights, and \texttt{return\_poles=True} returns them as
poles (energies and complex weights) rather than only as a curve.

The single approximation is truncation of the sum at \texttt{nex} states,
and it is \emph{measured} rather than assumed. The returned
\texttt{info["captured"]} is

\begin{equation}
\text{captured}=\frac{\sum_{n\le n_{ex}}|\langle n|B|\mathrm{GS}\rangle|^2}{\langle \mathrm{GS}|B^\dagger B|\mathrm{GS}\rangle},
\end{equation}

the fraction of the spectral weight of $B|\mathrm{GS}\rangle$ the
returned poles account for, and a warning fires below 0.9. \texttt{nex}
is also capped automatically at the target sector's exact dimension, so
asking for more states than exist is clipped rather than chased. Use this
method for a few sharp low-lying excitations; for a broad continuum the
low-lying states carry a small share of the weight and
\texttt{"KPM"}/\texttt{"TD"} are the right tools --- exactly the opposite
tradeoff.

Other keywords: \texttt{sector=} and \texttt{conserve=} override the
reference sector and which quantities are conserved (by default the
chain's own \texttt{set\_conserved\_sector}, or --- with none set --- the
charges measured on the unconstrained ground state, restricted to the
quantities the Hamiltonian actually conserves). Two operators whose
charges do not cancel raise, since the correlator then vanishes
identically by symmetry.

An operator with \emph{no} definite charge is handled rather than
rejected: $S_x$ raises \textbf{and} lowers $S_z$, so
\texttt{name=(Sx,Sx)} has no single target sector, but
$S_x=(S^++S^-)/2$ splits it exactly into two pieces that do, and the
channels whose charges cancel are summed --- an $S_z+1$ and an $S_z-1$
contribution, which is also the physically right decomposition.
\texttt{info["target\_sector"]} is then a list, one entry per channel.
The single-sector primitive \texttt{sectordc.sector\_poles} still
requires a definite charge, since its \texttt{info} names one sector and
carries that sector's matrix elements; use it when you want those (to
build $S(q,\omega)$ out of them, say).
\texttt{itensor\_version=3} and \texttt{itensor\_version="python"} only,
and never falls back to ED: falling back would answer with the
\emph{global} excited states, a different calculation.

\textbf{\texttt{get\_spectral\_function(i, j=None, spin=None, ...)} ---
the single-particle spectral function.} The physics-facing wrapper, since
``give me $A(\omega)$'' should not require assembling the hole part by
hand:

\begin{equation}
\begin{split}
A_{ij}(\omega)=&\sum_n\langle \mathrm{GS}|c_i|n^{N+1}\rangle\langle n^{N+1}|c_j^\dagger|\mathrm{GS}\rangle\,L(\omega-\omega_n^+)\\
&+\sum_m\langle \mathrm{GS}|c_j^\dagger|m^{N-1}\rangle\langle m^{N-1}|c_i|\mathrm{GS}\rangle\,L(\omega-\omega_m^-)
\end{split}
\end{equation}

with $\omega_n^+=E_n^{N+1}-E_0^N-\mu$ and
$\omega_m^-=-(E_m^{N-1}-E_0^N)-\mu$. The chemical potential
$\mu=(E_0^{N+1}-E_0^{N-1})/2$ is subtracted by default
(\texttt{shift="mu"}, \texttt{shift=None} for the raw axis), which is the
convention that formula is normally written in: particle weight then sits
at $\omega>0$ and hole weight at $\omega<0$, separated by the gap. That
is not cosmetic --- in the raw energies both families can land on the
same side of zero, for any Hamiltonian carrying an explicit
chemical-potential term or attractive interactions. \texttt{info} also
returns \texttt{mu}, the charge gap $E_0^{N+1}+E_0^{N-1}-2E_0^N$, and the
two halves separately. On a spinful chain pass
\texttt{spin="up"}/\texttt{"dn"}; the target sectors are then
$(N\pm1,S_z\pm1)$. See
\texttt{examples/dynamical\_correlator/sector\_spectral\_function\_hubbard},
where the Mott gap opens with $U$ and $\mu=U/2$ by particle-hole
symmetry.

\textbf{\texttt{get\_spin\_spectral\_function(i, j=None, ...)} --- the
dynamical spin structure factor, resolved by $S_z$ channel.} $S^{+-}$
from the $S_z+1$ sector, $S^{-+}$ from $S_z-1$, $S^{zz}$ from the ground
state's own sector; returns $S^{zz}+(S^{+-}+S^{-+})/2$ (the same
combination as $S^{xx}+S^{yy}+S^{zz}$) and, with
\texttt{return\_poles=True}, the three channels separately. Keeping them
separate is half the point: which channel a feature lives in is a
selection rule no single broadened curve can show. Note $S^{zz}$'s target
sector \emph{is} the reference sector, so its $n=0$ state is the ground
state itself and the raw Lehmann sum carries an elastic $\omega=0$ pole
of weight $\langle S^z_i\rangle\langle S^z_j\rangle$ --- that is the
correct Lehmann sum (the exact ED reference contains it too) rather than
the connected structure factor; pass \texttt{connected=True} to subtract
it.

One property specific to this method makes momentum-resolved work cheap:
nothing in the expensive half depends on \emph{which} operators are being
correlated, so one set of sector solves serves the entire matrix of
$(i,j)$ pairs. That is cached, so a full $S(q,\omega)$ or $A(k,\omega)$
map costs one pair of sector solves plus $L$ sets of matrix elements ---
measured, a 12-site $S(q,\omega)$ map takes about as long as a single
site's correlator. See
\texttt{examples/dynamical\_correlator/sector\_spin\_structure\_factor},
which reproduces the des Cloizeaux--Pearson lower edge of the two-spinon
continuum.

\textbf{\texttt{submode="maxent"} --- maximum-entropy reconstruction.}
Reconstructs a positive-definite spectral function from a finite set of
moments $\langle(H-E_0)^k\rangle$ using a maximum-entropy method
(\texttt{distribution.get\_distribution\_maxent}), rather than a
Chebyshev expansion --- useful when positivity of the reconstructed
$S(\omega)$ matters more than matching KPM's polynomial-expansion
artifacts.

\textbf{\texttt{submode="KPM"} for non-Hermitian Hamiltonians.} When
$H\neq H^\dagger$ (Section~\ref{sec:excited}), \texttt{submode="KPM"}
automatically routes to a different algorithm, a port of the
non-Hermitian Kernel Polynomial Method (NH-KPM) of
NHKPM.jl\footnote{\url{https://github.com/GUANGZECHEN/NHKPM.jl}}
(Phys.\ Rev.\ Lett.\ \textbf{130}, 100401):

\begin{lstlisting}
(x, y) = sc.get_dynamical_correlator(mode="DMRG", submode="KPM",
                                      name=(sc.Sz[0], sc.Sz[0]), E_max=10)
\end{lstlisting}

The ground state is now the biorthogonal pair
$|\psi_R\rangle,|\psi_L\rangle$ from NH-DMRG (Section~\ref{sec:excited})
rather than a single self-dual $|\mathrm{GS}\rangle$, and the correlator
computed is $\langle\psi_L|A(z)\,B|\psi_R\rangle$. Because the spectrum
of $H$ is complex, the ordinary Chebyshev recursion (which needs a
\emph{real} rescaled spectrum in $[-1,1]$) does not apply; NH-KPM instead
expands the frequency-shifted operator
$\tilde H(z)=(z\,\mathbb{1}-H)/E_{\max}$ using a \emph{coupled}
forward/adjoint recursion built from both $\tilde H(z)$ and $\tilde
H(z)^\dagger$ (see \texttt{src/dmrgpy/algebra/kpm.py}'s
\texttt{get\_mu\_n\_nh}/\texttt{spec\_from\_moments\_nh}, ported
line-for-line from the reference's
\texttt{get\_vn\_NH}/\texttt{get\_spec\_kpm\_NH}). The key practical
consequence is that, unlike the Hermitian case, the moments depend on
$z$ itself, so they are recomputed from scratch at every requested
frequency rather than amortized once over the whole spectrum --- this
mirrors the reference algorithm's own cost profile, and is why NH-KPM is
noticeably more expensive per frequency point than the Hermitian
\texttt{"KPM"} path. \texttt{E\_max} (an upper bound on the spectral
radius of $H$) must be supplied explicitly: unlike the Hermitian case's
variational band-edge estimate (see above), there is no automatic
estimator yet for a non-Hermitian spectral bound. Implemented so far for
the ED backend, \lstinline|itensor_version=3|, and
\lstinline|itensor_version="python"|; \lstinline|itensor_version=2|
raises \lstinline|NotImplementedError|. See
\texttt{examples/non\_hermitian/nhkpm\_v3\_VS\_ED} (ED vs
\lstinline|itensor_version=3|, machine-precision agreement on a small
interacting fermionic chain with a staggered imaginary potential) and
\texttt{examples/non\_hermitian/nhkpm\_python\_VS\_v3\_timing}
(\lstinline|itensor_version=3| vs \lstinline|"python"| on a non-uniform
hopping/non-uniform imaginary-onsite-energy chain, same
machine-precision agreement --- the pure-Python backend runs roughly 2x
slower than v3 for this workload, since NH-KPM's per-frequency moment
recursion is far more matvec-heavy than the Hermitian KPM path).

\textbf{The \lstinline|dex| cutoff --- which states count as ``the''
ground state (\lstinline|mode="ED", submode="ED"|).} The $T=0$ ED sum
above starts from the ground state, but a degenerate ground manifold has
no single one. The ED path therefore takes every eigenstate whose
excitation energy lies below a tolerance \lstinline|dex| (default
$10^{-5}$) as the initial manifold and averages the correlator over it
with \emph{equal weights} $1/n_{\rm ex}$. That is the right thing to do
for a genuinely degenerate manifold, and only for that case:
\lstinline|dex| is a hard step, so a level at $0.99\,\texttt{dex}$
contributes in full while one at $1.01\,\texttt{dex}$ contributes nothing
at all.

Consequently \lstinline|dex| must be chosen relative to the
\emph{splittings being resolved}, not to the broadening $\delta$ or the
frequency grid. The pathological case is a parameter sweep --- e.g.\ a
Zeeman field whose splitting grows through \lstinline|dex| --- where the
result jumps discontinuously as each level crosses the threshold, and the
sweep partially averages over the very splitting it was meant to resolve.
No single \lstinline|dex| works at both ends of such a sweep: it must be
large enough to keep the whole manifold at zero field and small enough to
exclude the split states at finite field. DMRGPY emits a
\lstinline|RuntimeWarning| whenever an eigenvalue lies within a factor of
three of \lstinline|dex| on either side, i.e.\ exactly when the answer
depends on where the cutoff was placed.

The clean way out of that regime is to stop using a cutoff at all and
weight the manifold physically instead: pass a small \lstinline|T| (see
the finite-temperature section), which switches to the
Boltzmann-weighted Lehmann sum. It reproduces the equal-weight degenerate
average smoothly as $T\to0$ --- at exact degeneracy the two agree to
machine precision --- while varying continuously, not in steps, once the
manifold splits. Use \lstinline|dex| for a one-off spectrum at a fixed
parameter, and \lstinline|T| for anything swept. See
\texttt{examples/kondo/atom\_iets\_orbital\_resolved}.

\textbf{Choosing a method:} KPM (default) for a first look at the full
spectrum; CVM or TD when you need high resolution in a specific,
narrow frequency window; TDZ instead of TD when that window also needs
long simulated times/low frequencies, where TD's real-time bond-dimension
growth becomes limiting; EX when a handful of excited states already
capture the physics (e.g.\ a small gapped system); maxent when you want
a guaranteed-positive reconstruction from limited moment data (e.g.\
combined with finite-temperature ED, see Section~\ref{sec:thermal}). For
a non-Hermitian $H$, \texttt{"KPM"} is currently the only submode with a
genuine biorthogonal implementation (see above); the other submodes fall
back to a correction-vector method that assumes $A^\dagger=B$.

\section{Real-time dynamics: quenches}
\label{sec:quenches}

Beyond frequency-domain correlators, DMRGPY directly simulates real-time
unitary evolution $|\psi(t)\rangle=e^{-iHt}|\psi(0)\rangle$ and measures
an observable along the way:

\begin{lstlisting}
from dmrgpy import timedependent
wf0 = sc.get_gs()                          # prepare |GS> of H0
sc.set_hamiltonian(h1)                     # quench to a different H
(ts, sz) = timedependent.evolve_and_measure(sc, operator=sc.Sz[0],
                                             nt=200, dt=1e-2, wf=wf0)
\end{lstlisting}

giving $\langle\psi(t)|S^z_0|\psi(t)\rangle$ as a function of time after
a Hamiltonian quench $H_0\to H_1$ --- the standard non-equilibrium
quantum quench setup: prepare the ground state of one Hamiltonian
(e.g.\ a symmetry-broken/N\'eel-ordered $H_0$), then let it evolve
under a different $H_1$ (e.g.\ the isotropic Heisenberg point) and watch
observables relax, oscillate, or thermalize. \texttt{evolution\_ABA}
similarly lets you apply an operator $A$ as an instantaneous local
quench (e.g.\ flip a spin, or add a particle) before evolving and
measuring $B(t)$. \texttt{timeevolution.imaginary\_exponential} computes
autocorrelation functions directly, $\langle\psi_0|e^{iHt}|\psi_0
\rangle$, without a separate measurement operator.

\textbf{Choosing the propagator: \texttt{sc.tevol\_method}.} Five
options:

\begin{itemize}
  \item \texttt{"TDVP"} (the default) --- two-site TDVP, which grows the
    MPS bond dimension via SVD the same way ground-state DMRG does. Used
    whenever \texttt{itensor\_version} is \texttt{3}, \texttt{"python"} or
    \texttt{"julia\_live"};
    \texttt{itensor\_version=2} falls back to \texttt{"MPO"} (below) even
    with this default.
  \item \texttt{"TDVP\_GSE"} (\texttt{itensor\_version} 3,
    \texttt{"python"} or \texttt{"julia\_live"}, same support as plain
    \texttt{"TDVP"} above)
    --- one-site TDVP preceded, for the first
    \texttt{sc.tdvp\_gse\_sweeps} steps (default 3), by a \emph{global
    subspace expansion} step: a Krylov subspace
    $\{\psi,H\psi,H^2\psi,\dots\}$ of dimension
    \texttt{sc.tdvp\_gse\_krylov\_order} (default 3) is used to enlarge
    the MPS's bond dimension \emph{without changing the state it
    represents}, using a cutoff \texttt{sc.tdvp\_gse\_cutoff} (default
    $10^{-8}$) --- the scheme of Yang \& White,
    \href{https://arxiv.org/abs/2005.06104}{arXiv:2005.06104}/Phys.\ Rev.\
    B 102, 094315 (2020). Most useful when the starting state's bond
    dimension is small (e.g.\ a product-state quench) and one-site TDVP
    alone (which conserves bond dimension exactly) wouldn't be able to
    grow into the entanglement the subsequent evolution generates. On
    \texttt{"julia\_live"} neither half is a dmrgpy port ---
    ITensorMPS.jl ships the expansion itself as
    \lstinline|expand(psi,H; alg="global_krylov")| (citing the same
    paper) and its \lstinline|tdvp| takes \lstinline|nsite=1| directly,
    so \texttt{mpsjulialive/tdvp.jl} only wires the two together; that
    route also handles a bond-dimension-1 (product-state) start fine,
    unlike \texttt{itensor\_version=3} (see
    \texttt{examples/time\_evolution/tdvp\_gse\_VS\_ED\_time\_evolution}'s
    own note on that edge case).
  \item \texttt{"TEBD"} (\texttt{itensor\_version} \texttt{3},
    \texttt{"python"}, or \texttt{"julia\_live"}, and only for a strictly
    nearest-neighbor Hamiltonian --- any term touching 3 or more distinct
    sites raises \texttt{NotImplementedError} (\texttt{"python"}), a
    catchable \texttt{RuntimeError} (\texttt{3}), or a
    \texttt{juliacall.JuliaError} (\texttt{"julia\_live"})) --- the
    standard 2nd-order-Trotter,
    even/odd-bond (``brick-wall'') algorithm: $e^{-iH\,dt}\approx
    e^{-iH_{\rm odd}\,dt/2}\, e^{-iH_{\rm even}\,dt}\,e^{-iH_{\rm
    odd}\,dt/2}$, where $H_{\rm odd}$/ $H_{\rm even}$ sum the bond
    Hamiltonians on odd-/even-indexed bonds (each internally commuting,
    since every bond in one group acts on disjoint sites). Onsite terms
    are split half-and-half onto each site's neighboring bond(s) (full
    weight at a chain boundary). Unlike the TDVP variants, every bond's
    evolution gate $e^{-i\tau h_{\rm bond}}$ is the exact exponential of
    the \emph{bare} local 2-site Hamiltonian, built once up front and
    reused unchanged for every time step --- no per-step Krylov/Lanczos
    work at all --- so \texttt{"TEBD"} is typically cheaper per step than
    TDVP whenever the Hamiltonian qualifies. \texttt{itensor\_version=3}
    builds the gate via ITensor's own \texttt{BondGate} primitive
    (\texttt{mpscpp3/tebd.h}); \texttt{"python"} exponentiates the bare
    2-site matrix directly (\texttt{pyitensor/tebd.py}'s
    \texttt{TEBDEvolver}); \texttt{"julia\_live"} builds the gate as an
    ITensor outer product of per-site operators and exponentiates via
    ITensors.jl's own tensor \texttt{exp()} (\texttt{mpsjulialive/tebd.jl})
    --- all three agree with each other (and with exact diagonalization)
    on a fermionic hopping+onsite benchmark. Unlike the other two
    backends, the Julia port resolves the Jordan-Wigner string itself,
    from scratch, off the \emph{raw} \texttt{"C"/"Cdag"} term list (the
    same one the MPO-based methods there already serialize), rather than
    \texttt{MultiOperator.to\_terms()}'s Jordan-Wigner-\emph{predressed}
    \texttt{"A"/"Adag"/"F"} form the other two backends consume --- real
    ITensors.jl's builtin \texttt{"Fermion"} site type only defines
    \texttt{op("C",..)}/\texttt{op("Cdag",..)}/\texttt{op("F",..)}, not
    dmrgpy's own \texttt{"A"/"Adag"} names.
  \item \texttt{"MPO"} --- a hand-rolled 2nd-order Taylor expansion of
    $e^{-iH\,dt}$ applied as an MPO each step; the only option on
    \texttt{itensor\_version=2} (which has no TDVP or TEBD at all), and
    available (if slower/less accurate for a given bond dimension)
    everywhere else too.
  \item \texttt{"AUTO"} (\texttt{itensor\_version} \texttt{3},
    \texttt{"python"}, or \texttt{"julia\_live"}, same support as
    \texttt{"TEBD"}) --- tries \texttt{"TEBD"} first and transparently
    retries as plain \texttt{"TDVP"} if the Hamiltonian turns out not to
    be strictly nearest-neighbor, instead of raising. This exists because
    ``is my Hamiltonian nearest-neighbor'' is often a property of
    \emph{which terms a caller adds}, not something pinned down once at
    chain-construction time --- a script that starts nearest-neighbor and
    later grows a longer-range term (or is reused across several models)
    would otherwise need its own \texttt{tevol\_method} bookkeeping to
    keep getting the cheaper \texttt{"TEBD"} path whenever it still
    applies. \texttt{"TEBD"} itself deliberately stays a hard opt-in that
    raises rather than silently falling back (see the note below) ---
    \texttt{"AUTO"} is the explicit way to ask for the fallback instead,
    so a caller who really meant \texttt{"TEBD"} and got a surprising
    long-range term still finds out. The check costs at most one
    discarded MPO build on the fallback path: both
    \texttt{bond\_hamiltonians()} implementations
    (\texttt{pyitensor/tebd.py}, \texttt{mpscpp3/tebd.h},
    \texttt{mpsjulialive/tebd.jl}) reject a non-nearest-neighbor term
    before touching the wavefunction at all, so retrying with
    \texttt{"TDVP"} never discards a partial time evolution. Not the
    default, precisely because it isn't a free win: whether TEBD applies
    becomes something you find out from behavior (a quietly different
    integrator, and thus different truncation-error characteristics)
    rather than from reading \texttt{sc.tevol\_method} --- worth it once a
    script's Hamiltonian shape is genuinely dynamic, not worth it as a
    blanket default for scripts that already know their Hamiltonian is
    nearest-neighbor and can just say \texttt{"TEBD"}.
\end{itemize}

\texttt{"julia\_live"} implements \texttt{"TDVP"}, \texttt{"TDVP\_GSE"},
\texttt{"TEBD"}, and \texttt{"AUTO"}; the legacy \texttt{"MPO"} path
raises \lstinline|NotImplementedError| there rather than silently running
plain TDVP instead, so a backend-comparison script can't quietly end up
comparing different integrators.

\begin{lstlisting}
sc.tevol_method = "TDVP_GSE"
sc.tdvp_gse_sweeps = 3
sc.tdvp_gse_krylov_order = 3
sc.tdvp_gse_cutoff = 1e-8
\end{lstlisting}

\begin{lstlisting}
sc.setup_cpp(version=3)    # or sc.setup_python(), or sc.setup_julia()
sc.tevol_method = "TEBD"
\end{lstlisting}

\section{Density of states}

\textbf{Many-body density of states.} The full many-body spectral
density

\begin{equation}
\rho(E)=\sum_n\delta(E-E_n)
\end{equation}

(sum over \emph{all} eigenstates of $H$, not resolved by any operator)
has no ready-made helper --- the closest available tool is
\texttt{get\_distribution}/\texttt{kpmdmrg.general\_kpm}
(Section~\ref{sec:dynamical}), but that computes a ground-state
expectation value $\langle\mathrm{gs}|B\,\delta(X)\,A|\mathrm{gs}\rangle$,
not a full-spectrum trace, so it is not a drop-in replacement for this
quantity.

\textbf{Single-particle local density of states.} For fermionic chains,
the physically distinct quantity usually meant by ``DOS'' (e.g.\ as
measured by scanning tunneling spectroscopy) is the local
single-particle spectral function at site $i$,

\begin{equation}
A_i(\omega)=-\frac1\pi\,\mathrm{Im}\,G^R_{ii}(\omega)
\end{equation}

with $G^R_{ii}(\omega)$ built from $\langle c_i(t)c_i^\dagger(0)\rangle$
($\omega>0$) and $\langle c_i^\dagger(t)c_i(0)\rangle$ ($\omega<0$),
i.e.\ the particle-addition and particle-removal (electron/hole)
branches of the local Green's function concatenated across $\omega=0$.
There is no ready-made helper for this (the module this section used to
point to, \texttt{fermiondos.py}, relied on a
\texttt{get\_dynamical\_correlator} calling convention that no longer
exists and was removed) --- build the two branches directly with
\texttt{get\_dynamical\_correlator} instead, using
\texttt{(fc.C[i], fc.Cdag[i])} for the particle-removal branch and
\texttt{(fc.Cdag[i], fc.C[i])} for the particle-addition branch, then
concatenate.

\section{Finite temperature}
\label{sec:thermal}

Thermal (mixed-state, finite-$T$) expectation values are obtained via
\textbf{purification}: each physical site is paired with an ancilla
site, and the maximally-entangled state of every physical-ancilla pair
(e.g.\ the singlet-forming Heisenberg coupling $\mathbf S_i^{\rm
phys}\cdot\mathbf S_i^{\rm anc}$ as the ``Hamiltonian'' preparing that
state) is exactly the purification of the infinite-temperature
($T=\infty$) physical density matrix $\rho\propto\mathbb 1$.
Imaginary-time evolving this purified state under the \emph{physical}
Hamiltonian $H$,

\begin{equation}
|\Psi(\beta)\rangle=e^{-\beta H/2}|\Psi(0)\rangle,\qquad\beta=1/T
\end{equation}

(stepped via repeated small first-order updates
$|\Psi(\beta+\Delta\beta)\rangle\propto(1-\Delta\beta\,H)|\Psi(\beta)\rangle$,
renormalized at each step) and tracing out the ancillas gives exactly
the thermal (Gibbs) density matrix of the physical chain,
$\rho(T)=\mathrm{Tr}_{\rm anc}|\Psi(\beta)\rangle\langle\Psi(\beta)|
\propto e^{-\beta H}$:

\begin{lstlisting}
from dmrgpy import thermal
tc = thermal.Thermal_Spin_Chain(spins, T=0.1)
\end{lstlisting}

$T=0$ recovers ordinary ground-state DMRG. For small systems (Hilbert
space dimension $\lesssim2000$, \texttt{algebra.maxsize}),
\texttt{get\_correlation\_matrix(T=...)} (ED only,
\texttt{entropytk/correlationentropy.py}'s
\texttt{get\_correlation\_matrix\_finiteT}) instead computes the exact
thermal average directly by brute-force ED,
$\rho=\sum_nP_n|n\rangle\langle n|$, $P_n=Z^{-1}e^{-(E_n-E_0)/T}$ --- a
useful cross-check of the purification approach, following the same
excited-states pattern as \texttt{vev(mode="ED", T=...)} below (full
spectrum whenever it's small enough to diagonalize exactly, an explicit
\texttt{n=} and a \texttt{RuntimeError} on inadequate truncation
otherwise). Automatic default operators (no \texttt{operators=} kwarg)
are only defined for fermionic chains today (\texttt{self.C}); on a
\texttt{Spin\_Chain}/\texttt{Thermal\_Spin\_Chain} pass an explicit
\texttt{operators=[...]} (e.g.\ \texttt{[sc.Sz[i] for i in range(n)]}).

\textbf{\texttt{vev(O, mode="ED", T=...)}} is a second, independent way
to get the same brute-force ED thermal average
$\langle O\rangle=\sum_nP_n\langle n|O|n\rangle$,
$P_n=Z^{-1}e^{-(E_n-E_0)/T}$, for a single operator rather than the whole
correlation matrix (\texttt{edtk/edchain.py}'s \texttt{EDchain.vev}
$\to$ \texttt{vevtk/thermalvev.py}'s \texttt{thermal\_vev\_ex}). The sum
only ever runs over a finite set of exact eigenstates obtained from
\texttt{algebra.lowest\_states}; whenever the full Hilbert space is small
enough to diagonalize exactly anyway (the same \texttt{maxsize=2000}-state
threshold \texttt{algebra.lowest\_states} itself uses),
\texttt{thermal\_vev\_ex} defaults to using the \emph{entire} spectrum,
so the thermal average is exact for any $T$. Above that size only a
sparse, truncated set of the lowest states is available (an explicit
\texttt{n=} kwarg, forwarded from \texttt{vev(...)}, controls how many);
\texttt{thermal\_vev\_ex} then checks the Boltzmann weight of the
highest included state and raises \texttt{RuntimeError} rather than
silently returning a wrong average if that state still carries
non-negligible weight at the requested $T$ --- pass a larger \texttt{n=}
in that case.

\textbf{\texttt{metts\_vev(O, T, ...)}} is a third, independent
finite-temperature method: METTS (Minimally Entangled Typical Thermal
States, E.M. Stoudenmire and S.R. White, \emph{New J. Phys.} \textbf{12},
055026 (2010), arXiv:1002.1305), implemented for
\texttt{itensor\_version="python"} (\texttt{pyitensor/metts.py}),
\texttt{itensor\_version=3} (\texttt{mpscpp3/chain\_session.h}'s
\texttt{Chain::metts\_vev}, a direct port of the same algorithm onto real
ITensor v3, not an independent reimplementation), and
\texttt{itensor\_version="julia\_live"} (\texttt{mpsjulialive/metts.jl}'s
\texttt{metts\_vev}, a value-level port of the same algorithm reusing
\texttt{tdvp.jl}'s \texttt{tdvp\_step} unchanged for the imaginary-time
evolution) --- not for \texttt{itensor\_version=2} (mpscpp2 has no TDVP
module at all, and METTS needs imaginary-time TDVP). On
\texttt{itensor\_version=3}, a chain
shorter than 3 sites raises \texttt{NotImplementedError} rather than
crashing: ITensor v3's two-site TDVP hits the same ``LocalOp is default
constructed'' abort as its two-site \texttt{dmrg()} for such short chains
(see the Architecture section's note on that \texttt{mpscpp3} bug).
Rather than purification's single, growing-entanglement wavefunction or
ED's exact sum over eigenstates, METTS samples a Markov chain of
\emph{unentangled} classical product states (CPS) $|i\rangle$: each is
imaginary-time evolved by half the inverse temperature,
$|\phi_i\rangle=e^{-\beta H/2}|i\rangle/\lVert e^{-\beta H/2}|i\rangle\rVert$
(via imaginary-time TDVP --- \texttt{tdvp\_step} with a purely real,
rather than imaginary, effective time step), then collapsed back down to
a new CPS $|i'\rangle$ by sampling a definite outcome at each site in
turn, with probability $|\langle i'|\phi_i\rangle|^2$, from a single
left-to-right sweep (no need to ever form the full $2^N$ amplitude
vector --- a canonicalized MPS's marginal at one site, conditioned on
those already sampled to its left, is already diagonal). A plain
(unweighted) sample average of $\langle\phi_i|O|\phi_i\rangle$ over the
resulting chain then converges to the true thermal average, because this
specific Markov chain's stationary distribution already carries the
correct Boltzmann weight --- no importance reweighting needed. The
collapse basis is alternated between two choices from one sample to the
next (\texttt{basis\_ops=("Sz","Sx")} by default) since collapsing
repeatedly in the same basis can trap the chain in one symmetry sector
for many steps, inflating autocorrelation (the paper's own finding):

\begin{lstlisting}
mean, stderr = sc.metts_vev(sc.Sz[0], T, nsamples=300, nwarmup=30,
                             dbeta_half_step=0.05, basis_ops=("Sz","Sx"))
\end{lstlisting}

\texttt{nwarmup} discards that many initial Markov-chain steps before
averaging (equilibration); \texttt{dbeta\_half\_step} sets the
imaginary-time TDVP step size for the $e^{-\beta H/2}$ evolution (split
into $\lceil(\beta/2)/\texttt{dbeta\_half\_step}\rceil$ equal steps).
Being a Monte Carlo method, \texttt{metts\_vev} returns
\texttt{(mean, stderr)} rather than an exact value --- \texttt{stderr} is
a naive i.i.d. estimate and, since consecutive METTS samples are
Markov-correlated, is likely optimistic unless
\texttt{dbeta\_half\_step}/\texttt{nwarmup} are generous enough that
samples actually decorrelate. Its main advantage over purification is
that the sampled states stay unentangled classical product states
between imaginary-time evolutions (no ancilla doubling, and bond
dimension never has to carry the entanglement of a single
ever-more-thermalized wavefunction) --- see
\texttt{examples/finite\_temperature/metts\_VS\_exact} for a cross-check
against the exact ED thermal average on a small Heisenberg chain.

\texttt{O} can also be a list/tuple of operators, measured together on
one shared sampled Markov chain instead of resampling from scratch for
each:

\begin{lstlisting}
results = sc.metts_vev([sc.Sz[0], sc.Sz[1], H], T, nsamples=300, nwarmup=30)
# results == [(mean_Sz0, stderr_Sz0), (mean_Sz1, stderr_Sz1), (mean_H, stderr_H)]
\end{lstlisting}

Since the \texttt{nwarmup+nsamples} imaginary-time evolutions dominate
the cost of \texttt{metts\_vev}, not the handful of extra
$\langle\phi|O_k|\phi\rangle$ measurements per sample, batching several
observables this way is far cheaper than calling \texttt{metts\_vev}
once per operator.

For \texttt{itensor\_version="python"}, an \texttt{njobs} keyword
(default 1) runs \texttt{njobs} independent METTS Markov chains in
parallel worker processes and pools their statistics, instead of one
longer sequential chain --- \texttt{nsamples} is split as evenly as
possible across them, each chain gets its own \texttt{nwarmup}
equilibration and an independently-seeded RNG, and their per-chain
\texttt{(mean, stderr, count)} triples are combined into the same
\texttt{(mean, stderr)} a single pooled run over all the raw samples
would give (no raw samples need to cross the process boundary, so this
is exact, not an approximation). Measured directly on a 6-site chain
(\texttt{nsamples=200}, \texttt{nwarmup=20}): wall time went from 34s at
\texttt{njobs=1} down to 14s at \texttt{njobs=8} --- real but sub-linear
speedup, since every extra chain repeats the full \texttt{nwarmup}
rather than sharing it. This is the effective optimization for this
backend: its per-sample cost is dominated by \texttt{pyitensor}'s own
generic-tensor-engine Python overhead (confirmed by profiling), not
shared BLAS work, so splitting across OS processes helps where the
\texttt{kernels.py} JAX/numba contraction-kernel route does not (see
that module's docstring, and \texttt{pyitensor/metts.py}'s own comment,
on why numba is measurably \emph{slower} for METTS specifically --- each
sample restarts from a fresh bond-dimension-1 product state, so a run
exercises many distinct contraction shapes rather than reusing one, and
the fixed per-shape compile tax never amortizes). \texttt{njobs} is not
available for \texttt{itensor\_version=3} or \texttt{"julia\_live"}: each
is a single live in-process session (a C++ object, or a live Julia
process) with no per-worker copy a process pool could hand out;
requesting \texttt{njobs>1} on either raises rather than silently falling
back to \texttt{njobs=1}. It is also unavailable while the pure-Python
backend is running on a device (\texttt{backend.set\_backend("jax")}, see
the GPU section of \texttt{docs/documentation.md}), for a related reason:
the worker processes are started with \texttt{spawn}, so each one
re-imports \texttt{pyitensor} with the default NumPy backend rather than
inheriting this process's device selection, and the chains would silently
run on the host --- so that combination raises too.

\textbf{\texttt{metts\_dynamical\_correlator(name, T, ...)}} extends
METTS from static expectation values to real-time finite-temperature
\emph{dynamical} correlators
$\mathcal{C}_{AB}(t)=\langle A(t)B\rangle_T=\langle e^{iHt}Ae^{-iHt}B\rangle_T$
(Z. Wang, P. McClarty, D. Dankova, A. Honecker and A. Wietek,
``Spectroscopy and complex-time correlations using minimally entangled
typical thermal states'', arXiv:2405.18484, Sec.~II, ``Dynamical METTS
algorithm''), implemented for \texttt{itensor\_version="python"},
\texttt{3}, and \texttt{"julia\_live"}
(\texttt{mpsjulialive/metts.jl}'s \texttt{metts\_dynamical\_correlator}, a
value-level port reusing the same \texttt{tdvp\_step}
\texttt{metts\_vev} already uses, now with a purely real time step for
the real-time evolution of $|v_i(t)\rangle$/$|w_i(t)\rangle$ instead of
the purely imaginary one used for sampling). For
every METTS sample $|\psi_i\rangle$ produced by the
exact same Markov chain \texttt{metts\_vev} already samples
(imaginary-time evolution + sequential-sampling collapse), define
$|v_i(0)\rangle=B|\psi_i\rangle$, $|w_i(0)\rangle=|\psi_i\rangle$,
real-time evolve both independently under $H$ (two-site TDVP), and
measure $\mathcal{C}^i(t)=\langle w_i(t)|A|v_i(t)\rangle$ at each
requested time step. A plain (unweighted) sample average of
$\mathcal{C}^i(t)$ over retained samples converges to
$\mathcal{C}_{AB}(t)$, for the same reason \texttt{metts\_vev}'s own
plain average converges to the thermal average --- no importance
reweighting needed:

\begin{lstlisting}
ts, means, stderrs = sc.metts_dynamical_correlator(
    (sc.Sz[0], sc.Sz[0]), T, nt=100, dt=0.1, nsamples=200, nwarmup=30,
    dbeta_half_step=0.05, basis_ops=("Sz","Sx"))
\end{lstlisting}

\texttt{name} follows the same \texttt{(A,B)} convention as
\texttt{get\_dynamical\_correlator}'s own \texttt{name=} (a string like
\texttt{"ZZ"}, or an explicit
\texttt{(MultiOperator,MultiOperator)} tuple/list) --- $A$,$B$ are used
exactly as given, with no dagger applied to either, matching
\texttt{get\_dynamical\_correlator(mode="ED", submode="ED", T=...)}'s own
convention (see below) so the two are directly comparable. \texttt{nt},
\texttt{dt} set the (uniformly spaced) real-time measurement grid
$t=0,\Delta t,\dots,(n_t-1)\Delta t$; \texttt{nsamples}, \texttt{nwarmup},
\texttt{dbeta\_half\_step}, \texttt{basis\_ops}, \texttt{seed},
\texttt{niter}, \texttt{njobs} all mean exactly what they mean for
\texttt{metts\_vev} (same shared Markov chain, same caveats on
\texttt{stderrs} being a Markov-correlated, likely-optimistic naive
estimate). \texttt{tdvp\_niter} separately bounds the Krylov iterations
used for the \emph{real-time} evolution of
$|v_i(t)\rangle$/$|w_i(t)\rangle$ specifically (default 50), independent
of \texttt{niter}'s bound on the imaginary-time sampling step (default
30) --- the two generally warrant different settings since
$|v_i(t)\rangle$/$|w_i(t)\rangle$ typically become more entangled over
the course of real-time evolution than the METTS samples
$|\psi_i\rangle$ themselves ever do (for \texttt{itensor\_version="julia\_live"},
\texttt{niter}/\texttt{tdvp\_niter} are accepted for signature parity but
silently ignored, same as \texttt{metts\_vev}'s own \texttt{niter} ---
ITensorMPS.jl's \texttt{tdvp()} manages its own internal Krylov dimension
with no exposed per-step iteration-count knob). No Fourier transform/windowing is
performed internally: \texttt{metts\_dynamical\_correlator} returns the
raw time-domain samples/statistics, matching \texttt{evolution\_DC}'s own
\texttt{(ts,cs)} convention --- apply a window (e.g.\ a Hann window, as
the reference paper recommends) and FFT separately if a
frequency-domain spectral function is wanted.

For a direct, exact reference to validate against,
\texttt{get\_dynamical\_correlator(mode="ED", submode="ED", T=...,
name=...)} computes the finite-temperature dynamical correlator's
spectral function via a full Boltzmann-weighted Lehmann sum over
\emph{every} ED eigenstate,
\begin{equation}
\mathcal{C}_{AB}(\omega)=\frac{1}{\mathcal{Z}}\sum_{n,m}e^{-\beta E_n}\langle n|A|m\rangle\langle m|B|n\rangle\,[\text{kernel at }\omega=E_m-E_n],
\end{equation}
the finite-$T$ generalization of the existing $T=0$
\texttt{submode="ED"} near-degenerate-ground-state sum to every
eigenstate weighted by its exact Boltzmann factor (exact, since it
starts from a full dense diagonalization, unlike
\texttt{thermal\_vev\_ex}'s own partial-diagonalization
truncation-safety check, which this doesn't need). See
\texttt{examples/finite\_temperature/dynamical\_metts\_VS\_ED}
(\texttt{itensor\_version} in \texttt{("python", 3)}) and
\texttt{examples/finite\_temperature/dynamical\_metts\_julia\_VS\_ED}
(\texttt{itensor\_version="julia\_live"}) for a
cross-check of \texttt{metts\_dynamical\_correlator}
against this exact ED reference, evaluated directly in the time domain,
on a small Heisenberg chain.

\section{Topological invariants}

\textbf{Many-body Berry phase.} For a ground state that depends on an
adiabatic parameter $k$ threaded around a closed loop (e.g.\ inserted
flux), the discretized Berry phase

\begin{equation}
\gamma=\arg\prod_{k}\langle\psi(k)|\psi(k+\delta k)\rangle
\end{equation}

quantized (typically in units of $\pi$) signals a topologically
nontrivial ground state / obstruction to adiabatic continuity --- the
many-body generalization of a Zak/Berry phase, computed by running DMRG
at a discrete set of parameter points and chaining ground-state
overlaps.

\textbf{Single-particle Berry phase / Wilson loop.} For a
translationally invariant single-particle Hamiltonian $H(k)$ with a set
of occupied bands, the non-Abelian Wilson loop over one Brillouin-zone
circuit is

\begin{equation}
W=\det\Big[\textstyle\prod_k U(k,k+\delta k)\Big],\qquad
U_{mn}(k,k+\delta k)=\langle u_m(k)|u_n(k+\delta k)\rangle
\end{equation}

with $m,n$ running over occupied bands; $\gamma=\arg W$ is the
polarization/Zak phase of that band manifold (and, combined with a scan
over a second momentum direction, the ingredient for a Chern number).

\section{Mean-field decoupling}

For a spin Hamiltonian with exchange couplings $J_{ij}$, a
self-consistent mean-field (Weiss-field) decoupling replaces the
two-body exchange term with a one-body field,

\begin{equation}
H_{\rm MF}=\sum_i\mathbf h_i\cdot\mathbf S_i,\qquad
\mathbf h_i=(1-p)\sum_jJ_{ij}\langle\mathbf S_j\rangle
\end{equation}

iterated to self-consistency: solve for $\langle\mathbf S_i\rangle$
under $H_{\rm MF}$ (by DMRG), rebuild $\mathbf h_i$ from the new
expectation values, mix old/new fields, and repeat until
$\max_i|\Delta\langle\mathbf S_i\rangle|$ falls below a tolerance. The
parameter $p\in[0,1]$ interpolates between pure mean-field theory
($p=0$) and keeping a fraction $p$ of the original many-body exchange
treated exactly alongside the self-consistent field (useful for hybrid
mean-field-plus-fluctuations treatments of, e.g., magnetically ordered
phases where pure MF overestimates order).

\begin{lstlisting}
from dmrgpy import meanfield
meanfield.spinchain_meanfield(sc, p=0.0)
\end{lstlisting}

\section{Fidelity susceptibility and quantum phase transitions}

For a Hamiltonian $H(\lambda)=H_0+\lambda H_1$ depending on a tuning
parameter $\lambda$, the fidelity susceptibility measures how sharply
the ground state changes as $\lambda$ varies --- and diverges at a
quantum phase transition, where the gap to the first excited state
closes:

\begin{equation}
\chi(\lambda)=\sum_{n\neq0}\frac{|\langle 0|H_1|n\rangle|^2}{(E_n-E_0)^2+\delta^2}
\qquad(\text{perturbative form, }\texttt{fmode="PT"})
\end{equation}

\texttt{get\_fidelity}'s \textbf{default} is not this perturbative form
but a non-perturbative estimator (\texttt{fmode="derivative"}) that
computes $\chi$ directly from finite differences of the ground-state
overlap matrix between nearby $\lambda$ values, which also handles a
(near-)degenerate ground-state manifold via a smooth gauge choice:

\begin{lstlisting}
from dmrgpy import fidelity
chi = fidelity.get_fidelity(sc, h0, h1, lam, n=3) # fmode="derivative" (default)
chi_pt = fidelity.get_fidelity(sc, h0, h1, lam, n=3, fmode="PT") # perturbative form above
\end{lstlisting}

Scanning $\lambda$ and plotting $\chi(\lambda)$, a peak (sharpening and
diverging with system size) locates a quantum critical point without
needing to know its universality class in advance.

\section{Ground-state degeneracy}

Exact and near (e.g.\ symmetry-protected, or finite-size-split
topological) ground-state degeneracies are estimated by a narrow,
super-Gaussian-broadened level count around the ground-state energy
(\texttt{degeneracy.py}'s \texttt{gs\_degeneracy\_simple}/
\texttt{eigenvalue\_degeneracy}):

\begin{equation}
g(E_0)\approx\sum_i\exp\!\left[-\left(\frac{(E_i-E_0)^2}{\delta}\right)^2\right]
\end{equation}

Note the quartic falloff in $(E_i-E_0)$ (not a plain Gaussian) --- this
makes the window considerably narrower than $\delta$ would suggest for a
standard Gaussian, so \texttt{delta} needs to be picked accordingly
(typically larger than the target energy resolution) to count
near-degenerate levels rather than only exactly-degenerate ones.

summed over a growing number of low-lying computed eigenstates $E_i$
until the count converges --- a value near an integer $g$ signals a
$g$-fold degenerate (or near-degenerate, at the working precision
$\delta$) ground-state manifold, as expected e.g.\ for the four
edge-state-split ground states of an open Haldane chain, or a
symmetry-broken ordered phase.

\begin{lstlisting}
from dmrgpy import degeneracy
g = sc.get_gs_degeneracy()
\end{lstlisting}

\section{Reduced density matrices and operator distributions}

\textbf{Reduced density matrix.} The single-site reduced density matrix
$\rho_i=\mathrm{Tr}_{j\neq i}|\mathrm{GS}\rangle\langle\mathrm{GS}|$ is
available directly (\texttt{sc.get\_rdm(i=0)}), the basic object
entanglement entropies, local observables, and further post-processing
(e.g.\ local susceptibilities) are built from.

\textbf{Operator distributions.} More generally, the full probability
distribution of an arbitrary Hermitian operator $X$ (not just its
expectation value) in the ground state,

\begin{equation}
P(x)=\langle\mathrm{GS}|\,\delta(X-x)\,|\mathrm{GS}\rangle
\end{equation}

is computed via the same KPM machinery as the dynamical correlators of
Section~\ref{sec:dynamical} (\texttt{sc.get\_distribution}), or
reconstructed from a finite set of raw moments $\langle X^k\rangle$ via
maximum entropy (\texttt{get\_distribution\_maxent}). Useful for e.g.\
full counting statistics of a conserved charge, or distinguishing a
sharply peaked (well-defined quantum number) ground state from a
broadly spread one.

\section{Post-processing tools}

\begin{itemize}
  \item \textbf{Analytic continuation} (\texttt{analyticcontinuation.py}):
    Pad\'e continuation of a correlator known on the
    imaginary/complex-frequency axis (e.g.\ from a Matsubara-like or
    complex-shifted CVM calculation, Section~\ref{sec:dynamical}'s
    \texttt{submode="CVMimag"}) to the real frequency axis, where the
    physical spectral function lives.
  \item \textbf{Function fitting} (\texttt{functionfit.py}): a generic
    multi-start Powell minimizer used e.g.\ to fit the Calabrese-Cardy
    entropy formula in Section~4.
  \item \textbf{Finite-size extrapolation} (\texttt{extrapolate.py}):
    polynomial extrapolation in $1/L$ of a size-dependent quantity
    $y(L)$ toward the thermodynamic limit $L\to\infty$ --- standard
    practice for extracting bulk quantities (energy density, order
    parameters, gaps) from finite DMRG chains.
  \item \textbf{Maximum-entropy reconstruction} (\texttt{reconstruct.py}):
    reconstructs a positive spectral function from a truncated moment
    expansion, underlying both the \texttt{"maxent"} dynamical-correlator
    submode and \texttt{get\_distribution\_maxent}.
\end{itemize}

\section{Worked-example cookbook}
\label{sec:cookbook}

\textbf{Central charge of a critical transverse-field Ising chain}
($H=\sum_iS_i^zS_{i+1}^z+\tfrac12\sum_iS_i^x$ at the critical field,
expected $c=\tfrac12$):

\begin{lstlisting}
sc.maxm = 200          # larger bond dimension: needed at criticality
wf = sc.get_gs()
print(wf.get_CFT_central_charge())
\end{lstlisting}

\textbf{Haldane gap of the spin-1 chain} (see Section~4):

\begin{lstlisting}
spins = ["S=1" for i in range(n)]
sc = spinchain.Spin_Chain(spins)
# ... build Heisenberg h ...
es = sc.get_excited(n=6)
print("Haldane gap:", es[4]-es[0])
\end{lstlisting}

\textbf{Charge gap of a Hubbard chain} (see Section~4):

\begin{lstlisting}
print("Single-particle gap:", fc.get_gap())
print("Charge (pair) gap:", fc.get_charge_gap(d=2))
\end{lstlisting}

\textbf{Fidelity susceptibility across the Ising transition} (see
Section~12):

\begin{lstlisting}
h0 = 4*sum(sc.Sz[i]*sc.Sz[i+1] for i in range(n-1))   # Ising coupling
h1 = 2*sum(sc.Sx[i] for i in range(n))                # transverse field
for lam in lambdas:
    chis.append(fidelity.get_fidelity(sc, h0, h1, lam, n=3))
\end{lstlisting}

\textbf{Momentum- and frequency-resolved dynamical structure factor}
$S(q,\omega)$ (see Section~\ref{sec:dynamical}), by combining the
site-resolved KPM correlator with a lattice Fourier transform:

\begin{lstlisting}
Sqw = {}
for i in range(n):
    for j in range(n):
        x, y = sc.get_dynamical_correlator(submode="KPM", name=(sc.Sz[i], sc.Sz[j]))
        Sqw[(i, j)] = (x, y)          # combine with sum_ij e^{iq(i-j)} S_ij(w) offline
\end{lstlisting}

\section{STM/Kondo tunneling spectra (third-order perturbation theory)}
\label{sec:kondospectrum}

\texttt{Spin\_Chain.get\_kondo\_spectrum} computes the differential
tunneling conductance $dI/dV(eV)$ of an STM tip coupled to one site of a
spin chain, following the weak-coupling (Kondo-scattering) perturbation
theory of Ternes, \emph{New J. Phys.} \textbf{17}, 063016 (2015),
\href{https://arxiv.org/abs/1505.04430}{arXiv:1505.04430}. This is a
different observable from the dynamical correlators of
Section~\ref{sec:dynamical}: instead of a retarded Green's function of
the spin system alone, it is the full Fermi's-golden-rule tunneling
current through tip+spin+sample, expanded to third order in the
tip-sample tunneling amplitude.

Two backends are available via \texttt{mode=}:

\begin{itemize}
\item \texttt{mode="ED"} (default): full exact diagonalization of the
  chain's Hamiltonian (every eigenstate is needed as a possible virtual
  intermediate state, not just the low-energy ones), independent of the
  chain's own \texttt{itensor\_version}/mode setting. Works at any
  $T\geq0$.
\item \texttt{mode="DMRG"}: \texttt{itensor\_version=3} throughout, never
  diagonalizing beyond the ground state -- only $T=0$ is supported. See
  "T=0 and the DMRG backend" below.
\end{itemize}

\begin{lstlisting}
sc = spinchain.Spin_Chain(["1/2"])
sc.set_hamiltonian(g*muB*B*sc.Sz[0])   # Zeeman-split S=1/2 impurity
eV, dIdV = sc.get_kondo_spectrum(eV_grid, site=0, Jrho_s=-0.05, U=0.25,
                                  T=1.0, order=3)
\end{lstlisting}

\textbf{Second order} (\texttt{order=2}) is the plain
spin-flip/potential-scattering Fermi golden rule result,

\begin{equation}
\frac{\partial I}{\partial V}(eV)\propto\sum_{i,f}p_i\Big[\tfrac12|\langle f|S_-|i\rangle|^2+\tfrac12|\langle f|S_+|i\rangle|^2+|\langle f|S_z|i\rangle|^2\Big]\,\Theta(eV-\epsilon_{if})+4U^2
\end{equation}

summed over both tunneling directions, with $p_i$ the Boltzmann
occupation of eigenstate $i$ at temperature $T$ and $\Theta$ a
temperature-broadened step function. This reproduces the textbook
inelastic-tunneling spin-flip steps at $eV=\pm(\epsilon_f-\epsilon_i)$
(e.g.\ the Zeeman step of a single $S=1/2$ impurity).

\textbf{Third order} (\texttt{order=3}, the default) adds two
corrections that require summing over \emph{all} eigenstates as virtual
intermediate states $m$ (energy conservation is not required for $m$): a
Kondo term (a Levi-Civita triple product of spin matrix elements
$\langle i|S|f\rangle$, $\langle f|S|m\rangle$, $\langle m|S|i\rangle$,
weighted by a temperature-broadened logarithmic function
$F(eV-\epsilon_m,T)$ that produces the characteristic zero-bias
Kondo-like resonance, splitting into two peaks under a Zeeman field),
and -- when \texttt{U!=0} -- a potential-scattering interference term
responsible for a bias-asymmetric lineshape.

\textbf{Both tunneling directions} are summed at every order, since the
measured $dI/dV$ is for the net current $I=I^{t\to s}-I^{s\to t}$. For
unpolarized tip and sample the matrix elements are direction
independent, so writing $g(eV)$ for the $t\to s$ expression the measured
terms are

\begin{equation}
\frac{\partial I}{\partial V}\Big|_{\rm 2nd,\,Kondo}=g(eV)+g(-eV),\qquad
\frac{\partial I}{\partial V}\Big|_{U\text{-}M}=g(eV)-g(-eV).
\end{equation}

The relative sign is fixed by the paper's own worked $S=1/2$ example:
the purely Kondo-like processes contribute with the \emph{same} sign in
both directions, while the potential-scattering ones change sign when
the tunneling direction is inverted. So the second-order and
third-order Kondo terms are \textbf{even} in bias (the zero-field Kondo
resonance sits exactly at $eV=0$), and the potential-interference term
is \textbf{odd} (zero at $eV=0$) -- that term is the sole source of bias
asymmetry in the spectrum.

\textbf{Scope and known limitations}, worth reading before trusting
specific numbers:

\begin{itemize}
\item Only a single chain site couples to the tip (\texttt{site=}); the
  paper's own model allows several sites with independent tip couplings,
  not implemented here.
\item The paper's own closed-form equations for two numerical building
  blocks -- the temperature-broadened step $\Theta(x)$ and the
  temperature-broadened Kondo log function $F(\epsilon,T)$ -- do not
  reproduce the physics the paper itself describes for them (checked
  directly: the printed $\Theta(x)$ diverges rather than saturating, and
  the printed $F$ closed form drops its own temperature broadening).
  \texttt{kondospectrumtk/stepfunctions.py} uses corrected forms
  instead, re-derived from the paper's own unambiguous defining
  integrals and verified against digitized values from the paper's own
  figures (see that module's docstring, and
  \texttt{examples/kondo/kondo\_spectrum\_VS\_paper/}).
\item The potential-interference term's general-spin closed form
  (\texttt{U!=0} in \texttt{order=3}) is an extrapolation from the
  paper's own worked $S=1/2$ example (only that special case is spelled
  out in closed form in the paper). Its overall normalization is
  nonetheless pinned down, together with every other prefactor in the
  spectrum, by the paper's absolutely scaled Figs.\ 3b/3d -- see below.
\item \textbf{Normalization is checked against absolute figure values.}
  Figs.\ 3b and 3d of the paper are plotted in absolute units
  ($e^2T_0^2/h$), so their zero-bias peak heights (1.13 at $U=0$, 1.39
  at $U=0.25$, for a single $S=1/2$ at $B=0$, $T=1\,$K,
  $J\rho_s=-0.05$, $\omega_0=20\,$meV) fix every prefactor at once.
  \texttt{examples/kondo/kondo\_spectrum\_VS\_paper/} and
  \texttt{tests/test\_kondo\_spectrum.py} assert on them directly. In
  particular they fix the spin-average normalization (``SA factor''):
  the paper's spin-averaged transition matrix element is twice the plain
  electron-spin trace, so the third-order Levi-Civita coefficient is
  $\mathrm{Im}[X]/2$ and the elastic potential channel is $4|U|^2$
  rather than the bare $|U|^2$ of the printed
  $|\mathcal{M}^{(1)}|^2$ equation -- see
  \texttt{kondospectrumtk/conductance.py}'s module docstring for the
  full bookkeeping.
\end{itemize}

\textbf{T=0 and the DMRG backend} (\texttt{mode="DMRG"}). At $T=0$ only
the ground state is thermally populated, which simplifies both terms
enough to avoid diagonalizing beyond the ground state entirely -- the
actual motivation for supporting \texttt{mode="DMRG"} in the first
place, since DMRG cannot enumerate excited states the way
\texttt{mode="ED"} does:

\begin{itemize}
\item \textbf{Second order} reduces to a $\Theta_0$-weighted (the exact,
  closed-form Heaviside limit of $\Theta$) cumulative integral of the
  ordinary $T=0$ dynamical structure factor
  $S_{\alpha\alpha}(\omega)=\sum_f|\langle f|S_\alpha|\mathrm{GS}\rangle|^2\delta(\omega-\epsilon_{f0})$,
  which \texttt{get\_dynamical\_correlator} already computes
  (\texttt{submode="KPM"} or \texttt{"CVM"}) without any excited-state
  enumeration (\texttt{kondospectrumtk/secondorder\_dc.py}).
\item \textbf{Third order} cannot be reduced to a two-operator correlator
  (it is a three-vertex object, with two \emph{different} intermediate
  states each weighted by a different function, $\Theta_0$ and $F_0$) --
  instead it is built from a Heisenberg three-point function
  $G(t_2,\tau)=\langle\mathrm{GS}|S_l(t_2+\tau)S_k(t_2)S_j(0)|\mathrm{GS}\rangle$,
  obtained via real-time TDVP evolution with a "checkpoint-and-branch"
  construction (evolve $S_j|\mathrm{GS}\rangle$ forward/backward in
  $t_2$, apply $S_k$ at each checkpoint, evolve each branch further in
  $\tau$, overlap with a fixed $S_l|\mathrm{GS}\rangle$ reference at
  every step), then extracted via two closed-form time-domain kernels
  (derived by inverse-Fourier-transforming $\Theta_0$ and $F_0$) rather
  than by evaluating those functions pointwise on a discrete frequency
  grid, which does not converge robustly for this construction -- see
  \texttt{kondospectrumtk/twotime.py}'s module docstring for the full
  derivation and the numerical pitfalls it was built to avoid
  ($\Theta_0$'s kernel is a Cauchy principal value, computed via an
  FFT-based Hilbert transform for machine-precision accuracy). This is
  the expensive part: cost scales with the number of $t_2$ checkpoints,
  each its own short TDVP trajectory
  (\texttt{kondospectrumtk/dmrgtwotime.py}).
\item \textbf{Potential-interference term} (\texttt{U!=0}, part of
  \texttt{order=3}) is also supported: its own $T=0$ limit collapses the
  excited-state sum to a convolution of the \emph{same} $T=0$ dynamical
  structure factor against the $F_0$ kernel instead of $\Theta_0$'s
  cumulative-sum weighting, so it reuses \texttt{get\_dynamical\_correlator}
  exactly like the second-order term above, needing no excited-state
  enumeration either (\texttt{kondospectrumtk/potentialdc.py}). Carries
  the same general-spin extrapolation caveat as
  \texttt{conductance.third\_order\_potential\_dIdV} (see that function's
  docstring). All three DMRG-side routines sum over both tunneling
  directions exactly as their \texttt{mode="ED"} counterparts do -- only
  the closed-form kernels depend on $eV$, so the $s\to t$ direction
  costs nothing beyond evaluating them a second time at $-eV$ (the
  expensive dynamical correlator / $G(t_2,\tau)$ is still built once).
\end{itemize}

Further \texttt{mode="DMRG"} limitations beyond the general ones above:
the second-order term's \texttt{es} frequency-grid parameter has no safe
default and must be supplied explicitly (it needs to cover every
relevant transition energy, which is a property of the chain's spectrum,
not of the \texttt{eV} sweep range -- see \texttt{second\_order\_dIdV\_dc}'s
docstring); the third-order term's \texttt{dt2}, \texttt{n\_t2\_half},
\texttt{dtau}, \texttt{n\_tau\_half} time-grid parameters likewise have
no safe default and must be supplied explicitly (a grid fine/wide enough
for the default $\omega_0$/$\Gamma_0$ needs $\sim10^5$--$10^6$ $t_2$
checkpoints, each its own real TDVP trajectory -- infeasible as a silent
default -- while a small, fast default is wildly under-resolved and
returns a finite but silently wrong result instead of erroring, confirmed
directly; see \texttt{two\_time\_kondo\_term\_dmrg}'s docstring); chains
need at least 3 sites (a 1-site chain hits an internal ITensor v3 error
unrelated to this feature, building the Hamiltonian MPO).

\texttt{kondospectrumtk/dmrgtwotime.py} was written against this
codebase's existing, verified DMRG API and validated once a compiled
ITensor v3 backend became available: $G(t_2,\tau)$ matches the ED
reference to $\sim10^{-9}$--$10^{-10}$ pointwise, and the swept
third-order Kondo term matches a grid-consistent ED reference to
$\sim10^{-10}$ (\texttt{tests/test\_kondo\_spectrum\_dmrgtwotime.py},
skipped automatically when no compiled \texttt{itensor\_version=3}
backend is available). Getting there surfaced three real bugs, none of
which showed up in the ED-only testing this module was originally
written against -- see that module's own docstring for the details (in
short: \texttt{tdvp\_step} silently renormalizes every step to unit norm,
discarding $S_j|\mathrm{GS}\rangle$'s true amplitude unless corrected for
explicitly; a forward/backward time-stepping bug meant "backward"
checkpoints never actually reached negative times; and a naive per-chunk
trapezoidal integral is exactly 0 for the single-$t_2$-point chunks
real-time evolution necessarily produces, silently zeroing the entire
term). The second-order term (\texttt{submode="KPM"}) was spot-checked
too, agreeing to within a few tens of percent at thresholds, consistent
with the expected $\delta$-broadening/moment-truncation error on top of
what the ED path already has.

\subsection{Orbital-resolved IETS of a magnetic atom}

\lstinline|dmrgpy.atom| describes a \emph{single} multi-orbital magnetic
atom rather than a chain: \lstinline|generate_atom(orbs, tij, U, J, soc, B, Js, Ne)|
builds a \lstinline|Spinful_Fermionic_Chain| whose ``sites'' are the
atom's five $d$ orbitals, with a crystal field \lstinline|tij|,
spin-orbit coupling \lstinline|soc|, Hund's coupling \lstinline|J|, an
$-U S^2-J L^2$ interaction, an external field \lstinline|B| acting as
$2\mathbf{S}+\mathbf{L}$, and the electron count \lstinline|Ne| fixed by
a Lagrange multiplier. Two inelastic-tunneling (IETS) observables are
then available:

\begin{itemize}
\item \lstinline|get_spinflip(fc, es=..., iorb=...)| --- the spin-flip
  contribution, $\sum_{a=x,y,z} S_{aa}(\omega)$ built from the
  fluctuation operators $S^a_{i}-\langle S^a_{i}\rangle$ on orbital
  \lstinline|iorb|, returned over both bias polarities.
\item \lstinline|get_orbital_cotunneling(fc, es=..., iorb=...)| --- the
  orbital cotunneling contribution, summing correlators of
  $c^\dagger_{\texttt{iorb}\sigma}c_{j\sigma'}$ over every \emph{other}
  orbital $j$ and both spin channels, again for both bias polarities.
\end{itemize}

\lstinline|iorb| selects the orbital the STM tip couples to (e.g.\ $d_{z^2}$
for a tip directly above the atom), and it matters: with a crystal field
on, different orbitals give different spectral weights \emph{and}
different peak positions. An out-of-range \lstinline|iorb| raises
\lstinline|ValueError|. Both routines run in ED by default
(\lstinline|mode="ED"|, \lstinline|submode="ED"|), so the
\lstinline|dex|/\lstinline|T| caveat of the dynamical-correlator section
applies directly to them --- the low-lying multiplet of a magnetic atom
is precisely the near-degenerate manifold that discussion is about. See
\texttt{examples/kondo/atom\_iets\_orbital\_resolved}.

\textbf{Note (behaviour change).} Before 2026-08-21 both entry points
accepted \lstinline|iorb| and then silently overwrote it with \lstinline|0|,
so every call returned orbital 0's spectrum regardless of what was
requested. Results obtained with \lstinline|iorb!=0| on an earlier
version need regenerating.

\section{Infinite chains (iDMRG)}

Every method above works on a \emph{finite} chain of \texttt{n} sites.
\texttt{Infinite\_Many\_Body\_Chain}/\texttt{Infinite\_Spin\_Chain}
(\texttt{infinitechain.py}) instead describe a translationally-invariant
chain that repeats a fixed-size \textbf{unit cell} forever in both
directions, and solve it with infinite DMRG (iDMRG, White's growing
algorithm generalized to a multi-site unit cell) rather than sweeping a
fixed-length system. Currently limited to a 1- or 2-site unit cell
(\texttt{n\_uc<=2}) -- a $>$2-site unit cell raises
\texttt{NotImplementedError} at construction, see
\texttt{pyitensor/idmrg.py}'s module docstring for why.

\texttt{itensor\_version="python"} (default) or \texttt{3} are both
supported (pass it as a constructor keyword, e.g.\
\texttt{Infinite\_Spin\_Chain(["1/2"], itensor\_version=3)}) -- \texttt{2}
and \texttt{"julia\_live"} have no iDMRG port and also raise
\texttt{NotImplementedError}. \texttt{ic.gs\_method} (\texttt{"vumps"} by
default, on EITHER backend since 2026-08-08 -- see below) picks the
ground-state solver. Both solvers, on both backends, support
\texttt{vev}/\texttt{correlator} (see "Static correlators" below): the v3
C++ backend's \texttt{gs\_method="idmrg"}
(\texttt{mpscpp3/chain\_session.h}'s \texttt{Chain::idmrg\_ground\_state})
is a line-by-line port of \texttt{pyitensor/idmrg.py}'s own growing
algorithm, and the three things the Python side had gained since that
port -- McCulloch's wavefunction prediction, the gauge-consistent
unit-cell extraction described below, and the per-site energy baseline --
are ported too, along with the static observables built on them
(\texttt{Chain::idmrg\_onsite\_expectation},
\texttt{Chain::idmrg\_two\_point\_correlator} and
\texttt{Chain::idmrg\_local\_excitation\_gap}). Relative speed depends
on the model: v3's own \texttt{idmrg\_ground\_state} is $\sim$2.6--2.7x
slower than \texttt{"python"}'s for a gapless (critical) chain (the
local 2-site solve needs close to its full Krylov dimension every
macro-iteration), but can be \emph{faster} for a gapped model, where
that same solve converges quickly -- see
\texttt{docs/documentation.tex}'s iDMRG section for benchmark numbers.

\textbf{Building the Hamiltonian: \texttt{L}/\texttt{C}/\texttt{R}-suffixed
operators.} Instead of one operator list per absolute site index, an
infinite chain exposes each operator three times per unit-cell site
\texttt{i} (\texttt{i=0..n\_uc-1}): \texttt{SxC[i]} (site \texttt{i} of
the \emph{central} cell -- ordinary intra-cell use), \texttt{SxR[i]}
(site \texttt{i} of the \emph{next} cell), and \texttt{SxL[i]} (site
\texttt{i} of the \emph{previous} cell, provided purely so a coupling can
be phrased in whichever direction reads most naturally). A term may
touch at most two \emph{adjacent} unit cells (any subset of
\texttt{\{L,C\}} or of \texttt{\{C,R\}}, never both \texttt{L} and
\texttt{R} at once) -- \texttt{set\_hamiltonian} raises
\texttt{ValueError} otherwise:

\begin{lstlisting}
from dmrgpy import infinitechain
ic = infinitechain.Infinite_Spin_Chain(["1/2"])       # n_uc=1, uniform chain
h = ic.SxC[0]*ic.SxR[0] + ic.SyC[0]*ic.SyR[0] + ic.SzC[0]*ic.SzR[0]  # NN Heisenberg
ic.set_hamiltonian(h)
\end{lstlisting}

A bare, single-site operator (e.g.\ \texttt{ic.SzC[0]}, no product) is
also a valid term -- a Zeeman-field-style onsite Hamiltonian, either on
its own or added alongside bond terms.

A two-site unit cell can express a dimerized (alternating-bond) chain,
and a coupling can skip over an intermediate site by reaching straight
from \texttt{C} to \texttt{R}:

\begin{lstlisting}
ic = infinitechain.Infinite_Spin_Chain(["1/2", "1/2"])
h = (1.0*(ic.SxC[0]*ic.SxC[1] + ic.SyC[0]*ic.SyC[1] + ic.SzC[0]*ic.SzC[1])   # strong intra-cell bond
     + 0.4*(ic.SxC[1]*ic.SxR[0] + ic.SyC[1]*ic.SyR[0] + ic.SzC[1]*ic.SzR[0]))  # weak inter-cell bond
ic.set_hamiltonian(h)
\end{lstlisting}

\textbf{Ground-state energy density.} Since the chain is infinite, the
physically meaningful quantity is the energy \emph{per site}, not a
total:

\begin{lstlisting}
density = ic.gs_energy()   # converged energy density (or the best value reached after ic.maxiter)
ic.converged                # True iff the density stabilized below ic.etol
\end{lstlisting}

\texttt{ic.maxm}/\texttt{ic.cutoff} cap the bond dimension/SVD truncation
exactly like a finite chain's \texttt{maxm}/\texttt{cutoff};
\texttt{ic.maxiter}/\texttt{ic.etol} are iDMRG-specific: \texttt{maxiter}
macro-iterations (each adds one full unit cell to each side) are run,
stopping early once the energy-density finite difference between
consecutive macro-iterations drops below \texttt{etol}. For a
\emph{gapless} model (e.g.\ the uniform Heisenberg chain above) this
finite-difference convergence is only power-law in the number of
iterations at any fixed \texttt{maxm}, not exponential, so it may not
trip \texttt{etol} within a practical \texttt{maxiter} even though the
energy density itself is already accurate to several digits -- check the
returned value directly rather than relying solely on
\texttt{ic.converged}. A \emph{gapped} model (e.g.\ the dimerized chain
above, or any chain with \texttt{j\_weak} bonds well below the uniform
point) has a finite correlation length and converges both faster and
more reliably.

\textbf{Fermionic infinite chains.}
\texttt{Infinite\_Many\_Body\_Chain}'s first argument is a list of
site-type codes, so an infinite chain can be fermionic as well as a spin
chain: \texttt{0} is a spinless fermion site
(\texttt{C}/\texttt{Cdag}/\texttt{N}/\texttt{F}), \texttt{1} a native
spinful (Electron/Hubbard) site carrying both flavours at once
(\texttt{Cup}/\texttt{Cdagup}/\texttt{Cdn}/\texttt{Cdagdn}/\texttt{Nup}/\texttt{Ndn},
local dimension 4). The native spinful code is what makes a two-orbital
model fit the \texttt{n\_uc<=2} limit --- a
\texttt{Spinful\_Fermionic\_Chain}-style representation would need two
tensor-network sites per orbital, i.e.\ four for a two-orbital cell:

\begin{lstlisting}
ic = infinitechain.Infinite_Many_Body_Chain([1, 1])   # two native spinful sites per cell
Cup  = [ic.get_operator("Cup", i, "C") for i in range(2)]
Cdup = [ic.get_operator("Cdagup", i, "C") for i in range(2)]
CupR = [ic.get_operator("Cup", i, "R") for i in range(2)]
ic.set_hamiltonian(Cdup[0]*CupR[0] + ...)
\end{lstlisting}

Both backends thread the Jordan-Wigner string themselves, \emph{locally}
between each term's own two endpoints, rather than from some absolute
origin --- an infinite chain has no site 1 for a finite-chain-style
string to start at. Two consequences for what a Hamiltonian may contain:

\begin{itemize}
\item The ``at most 2 distinct sites per term'' rule is counted in
      \emph{sites}, not operator factors: a three-operator product like
      \texttt{(N\_f - 1/2) * (Cdag\_f * C\_c)} touches two sites and is
      fine. A term whose two endpoints have sites strictly between them
      (e.g.\ \texttt{Cdag} at cell site 0 and \texttt{C} at site 0 of the
      \emph{next} cell, for \texttt{n\_uc=2}) is also fine --- the string
      across the intervening sites is inserted automatically.
\item A term with \textbf{odd total fermion parity} (e.g.\ a bare
      \texttt{Cdag}) is rejected with \texttt{ValueError} from
      \texttt{set\_hamiltonian}: its string would have to run to infinity
      in both directions. Such a term is not parity-conserving, so it
      cannot appear in a physical Hamiltonian anyway.
\end{itemize}

\textbf{Fermionic correlators} work the same way:
\texttt{ic.correlator("Cdag", 0, "C", r)} returns the physical fermionic
correlator, i.e.\ $\langle C^\dagger_0 (\prod_{0<k<r} F_k) C_r\rangle$
with the Jordan-Wigner string threaded across every site strictly between
the two operators --- not the stringless product of two bare matrices,
which is a different quantity whose error \emph{grows with separation}
(so it corrupts exactly the decay rate such a correlator is usually
measured for). Supported on all three paths:
\texttt{itensor\_version="python"} with \texttt{gs\_method="idmrg"} or
\texttt{"vumps"}, and \texttt{itensor\_version=3} with
\texttt{gs\_method="vumps"}. The endpoint matrices and the decision of
whether a string is open at all come from the same helper that builds the
Hamiltonian's own 2-site terms, so a correlator and a Hamiltonian term
written with the same operator names cannot disagree about the
convention. Parity-even operators (\texttt{N}, \texttt{Nup},
\texttt{Sz}, \ldots) get no string, exactly as before.

A pair with \textbf{odd total fermion parity} (e.g.\
\texttt{("Cdag", "N")}) raises \texttt{ValueError}: its string can never
close on an infinite chain, and the quantity vanishes identically in any
parity-conserving state anyway.

See \texttt{examples/idmrg/fermionic\_infinite\_chain/main.py}, which
checks both the energy density (against the exact free-fermion band
integral) and $\langle C^\dagger_0 C_r\rangle$ (against the exact
one-body density matrix) on every backend/solver combination, for a chain
including a hopping that skips a site so a real string is involved.

\textbf{Static correlators.} After \texttt{gs\_energy()} (called
automatically by \texttt{vev}/\texttt{correlator} if not already run),
one- and two-point expectation values of the converged infinite chain are
available via the standard infinite-MPS transfer-matrix formalism.
All four combinations work: \texttt{itensor\_version}
\texttt{"python"} or \texttt{3}, each under \texttt{gs\_method="vumps"}
(the default, see below) or \texttt{gs\_method="idmrg"} (described in this
paragraph). The two backends are cross-checked directly against each
other (\texttt{tests/test\_idmrg\_correlator\_v3.py},
\texttt{examples/idmrg/idmrg\_correlator\_python\_VS\_v3}), agreeing to
$\sim$1e-8 or tighter on a gapped chain:

\begin{lstlisting}
ic.vev("Sz", 0)                    # <Sz> at site 0 of the unit cell
ic.correlator("Sz", 0, "Sz", r)    # <Sz(0) Sz(0+r)>, r measured in physical sites, r>=0
\end{lstlisting}

These are reconstructed \emph{after} convergence, from the
gauge-consistent unit cell the growing algorithm extracts from a single
micro-step's own two-site wavefunction (\texttt{pyitensor/idmrg.py}'s
\texttt{\_theta\_cell}, \texttt{Chain::idmrg\_theta\_cell} on the C++
side). Fermionic operators are handled: the Jordan-Wigner string is
threaded across every site strictly between the two endpoints, so
\texttt{ic.correlator("Cdag", 0, "C", r)} is the physical fermionic
correlator, not a stringless product of two bare matrices (a pair of odd
total fermion parity, whose string could never close, raises
\texttt{ValueError}). Two cheap diagnostics are worth knowing about.
$\langle H_{uc}\rangle$ -- the sum of every bond's correlator inside one
unit cell -- must equal \texttt{n\_uc * density} exactly for any
translationally invariant state, so it is a model-agnostic check on
whether a given \texttt{maxm}/\texttt{maxiter} is enough, needing no
reference value. \texttt{ic.\_result.state\_overlap} (on the raw
\texttt{pyitensor.idmrg.IDMRGResult}) is the second: the overlap between
each macro-iteration's converged local state and the prediction it
started from, which approaches 1 once the state has genuinely stopped
changing.

Both were, until recently, routinely far from their ideal values, and
correlators built on \texttt{n\_uc=1} in particular could come out with
the wrong \emph{sign}. That is fixed: the growing algorithm now carries
the state across iterations with McCulloch's wavefunction prediction, and
extracts the unit cell in a single, self-consistent gauge. Measured
against exactly solvable references,
$\langle H_{uc}\rangle - \texttt{n\_uc}\cdot\text{density}$ now lands at
$10^{-15}$--$10^{-9}$ (it previously missed by up to $0.12$),
\texttt{state\_overlap} reaches $1-10^{-13}$ (it previously plateaued
around 0.5--0.65 for \texttt{n\_uc=1}), and the XX chain's
$\langle S^z\rangle$ comes out at $10^{-13}$ against an exact 0.
Correlators still converge more slowly in \texttt{maxm} than the energy
density does -- that is ordinary finite-bond-dimension physics, most
visible for a gapless model -- so the $\langle H_{uc}\rangle$ check
remains the right thing to run before trusting a number.

\textbf{The default ground-state solver: VUMPS.}
\texttt{ic.gs\_method = "vumps"} (the default since 2026-08-08) switches
\texttt{gs\_energy()} from the growing algorithm above to VUMPS
(Variational Uniform Matrix Product States, Zauner-Stauber et al.,
arXiv:1701.07035; see \texttt{pyitensor/vumps.py}'s own module
docstring for the algorithm) -- instead of growing a finite window and
truncating it down to \texttt{maxm} at every step (the
\texttt{gs\_method="idmrg"} growing algorithm above), VUMPS solves
directly, in the thermodynamic limit, for the actual
\texttt{maxm}-dimensional variational optimum (\texttt{ic.maxm} sets
VUMPS's own target bond dimension $D$ here too). Both
\texttt{itensor\_version="python"} (\texttt{pyitensor/vumps.py}) and
\texttt{itensor\_version=3} (\texttt{mpscpp3/chain\_session.h}'s
\texttt{Chain::vumps\_ground\_state}, a C++ port of the same algorithm
-- built from plain dense arrays closed over LAPACK rather than ITensor
tensor-network objects, since the bond/physical dimensions this feature
targets are always small; see that method's own doc comment) support
\texttt{gs\_method="vumps"}; the two are cross-checked directly against
each other to $\sim$$10^{-10}$ or tighter on TFIM/Heisenberg at
$D=1,2,3$ (\texttt{tests/test\_vumps\_v3.py}). Explicitly set
\texttt{ic.gs\_method = "idmrg"} instead for
\texttt{local\_excitation\_gap}/\texttt{td\_dynamical\_correlator} (no
VUMPS equivalent, see their own sections below) or for VUMPS's
documented $D>1$ convergence-robustness gap (see below) if the growing
algorithm's own more battle-tested behavior is preferred:

\begin{lstlisting}
ic = infinitechain.Infinite_Spin_Chain(["1/2"])
ic.gs_method = "vumps"
ic.maxm = 4        # VUMPS's own target bond dimension D
ic.maxiter = 800   # VUMPS outer-iteration cap per bond dimension in its own D-ramp
ic.vumps_nrestarts = 6   # independent random-restart attempts per bond dimension
h = ic.SxC[0]*ic.SxR[0] + ic.SyC[0]*ic.SyR[0] + ic.SzC[0]*ic.SzR[0]
ic.set_hamiltonian(h)
density = ic.gs_energy()
\end{lstlisting}

\texttt{vev}/\texttt{correlator} also work under
\texttt{gs\_method="vumps"}, on BOTH backends:
\texttt{pyitensor.vumps.onsite\_expectation}/\texttt{two\_point\_correlator}
for \texttt{itensor\_version="python"}, and
\texttt{Chain::vumps\_onsite\_expectation}/\texttt{vumps\_two\_point\_correlator}
(a line-for-line C++ port of the same formula) for
\texttt{itensor\_version=3} -- cross-checked directly against each other
to $\sim10^{-14}$ or tighter on TFIM at $D=2,3$
(\texttt{tests/test\_vumps\_correlator\_v3.py}). Both are computed
directly from the converged mixed-gauge \texttt{\{AC, AR\}} rather than
\texttt{pyitensor.idmrg}'s dominant-right-fixed-point eigenproblem:
\texttt{AC} is already the exactly-normalized single-(super)site reduced
state by construction of the mixed canonical gauge (Vanderstraeten,
Haegeman, Verstraete, ``Tangent-space methods for uniform matrix product
states'', arXiv:1810.07006, Eq.(34)), and \texttt{AR}'s exact
right-orthonormality lets a two-point correlator spanning multiple unit
cells close by a direct trace with no eigenproblem either (the
mixed-gauge analogue of that same review's Eq.(37)--(39)). Unlike the
growing-algorithm's own reconstructed-from-the-last-macro-iteration
correlators above, these carry no \texttt{maxiter}/$\langle
H_{uc}\rangle$-self-consistency caveat -- VUMPS solves directly at the
target bond dimension in the thermodynamic limit, so once
\texttt{.converged} is \texttt{True} the correlator is exact for that
converged \texttt{\{AL,AR,C\}}, only limited by the bond dimension $D$
itself (same caveat as the energy density, see below).
\texttt{local\_excitation\_gap} is the one method that still requires
\texttt{gs\_method="idmrg"} specifically (it re-diagonalizes the growing
algorithm's own final 2-site effective Hamiltonian, which has no VUMPS
equivalent), on either backend at \texttt{window=0}
(\texttt{Chain::idmrg\_local\_excitation\_gap} for
\texttt{itensor\_version=3}); its \texttt{window>0} variant is
\texttt{itensor\_version="python"}-only, being an explicit prototype
rather than stable API; conversely
\texttt{excitation\_energies}/\texttt{excitation\_gap} (below)
\emph{require} \texttt{gs\_method="vumps"} (the default) -- they need
\texttt{VUMPSResult}'s own mixed-gauge \texttt{\{AL,AR,C,GL,GR\}}, which
the growing algorithm's \texttt{IDMRGResult}
(\texttt{gs\_method="idmrg"}) has no equivalent of.

\begin{lstlisting}
ic.gs_method = "vumps"
ic.vev("Sz", 0)                    # <Sz> at site 0, from the converged VUMPSResult
ic.correlator("Sz", 0, "Sz", r)    # <Sz(0) Sz(0+r)>, same signature as gs_method="idmrg"
\end{lstlisting}

\textbf{A real, scoped reliability caveat.} \texttt{D=1} converges
reliably and exactly for every exactly-solvable (product-state-like)
model tried (a pure field, a fully-decoupled Heisenberg dimer). Already
at \texttt{D>1}, though -- confirmed directly on the transverse-field
Ising model, a simple gapped model with no special critical/symmetry
complications -- independent calls at the same \texttt{D} can land on
noticeably different converged energies (most within a percent or two
of the exact answer, but occasionally further off), even with
\texttt{vumps\_ground\_state}'s own built-in $D$-ramp, multi-restart,
and variational-principle safety-net machinery (see
\texttt{pyitensor/vumps.py}'s own ``Convergence robustness'' docstring
section for the full, numerically-confirmed account). A caller needing
a reliable \texttt{D>1} result should call \texttt{gs\_energy()} a few
times independently (e.g.\ on freshly-constructed chains) and keep the
lowest reported density, the same ``rerun and take the best''
discipline \texttt{idmrg\_ground\_state}'s own unseeded-random-MPS
initialization already recommends elsewhere in this section. See
\texttt{examples/idmrg/vumps\_TFIM/main.py} for a worked example
sweeping \texttt{D} and comparing VUMPS against both
\texttt{gs\_method="idmrg"} and the transverse-field Ising model's own
exact (free-fermion) energy density.

Entanglement/entropy are not implemented for infinite chains yet.

\textbf{Excited states: the tangent-space/quasiparticle excitation
ansatz.} Unlike a finite chain, an infinite chain's excitations form a
momentum-resolved band $E(k)$, not a single discrete state, so ``the
first excited state'' here means the standard single-mode/quasiparticle
ansatz (Haegeman et al.), built on top of a \textbf{VUMPS} ground
state's own mixed-gauge \texttt{\{AL,AR,C\}} representation
(\texttt{ic.gs\_method = "vumps"} -- required, see above): a
tangent-space vector with one excitation tensor $B$ (a function of a
free matrix $X$) inserted at every unit-cell position, weighted by
momentum $k$. \texttt{ic.excitation\_energies(k, n=1)} returns the
lowest \texttt{n} excitation energies (above the ground state) at
momentum \texttt{k} (radians, per unit cell);
\texttt{ic.excitation\_gap(ks=None)} scans \texttt{k} (default
\texttt{numpy.linspace(-pi, pi, 41)}) and returns the minimum -- the
scalar ``gap'', mirroring the finite-chain \texttt{get\_gap()} naming
(\S4):

\begin{verbatim}
ic.gs_method = "vumps"             # required for excitation_energies
ic.excitation_energies(0.0, n=1)   # lowest excitation energy at k=0
ic.excitation_gap()                # min_k E(k), the scalar gap
\end{verbatim}

\textbf{Scope.} Both \texttt{itensor\_version="python"}
(\texttt{pyitensor/idmrg\_excitations.py}) and
\texttt{itensor\_version=3} (\texttt{mpscpp3/chain\_session.h}'s
\texttt{Chain::vumps\_excitation\_energies}, a C++ port of the same
algorithm, same dense-array/LAPACK approach as
\texttt{vumps\_ground\_state} above) are supported -- both require
\texttt{gs\_method="vumps"} (\texttt{NotImplementedError} otherwise, or
if a different \texttt{itensor\_version} is used). Any converged bond
dimension \texttt{D>=1} is supported -- including a genuinely entangled
ground state (\texttt{D>1}, e.g.\ the transverse-field Ising or a
dimerized Heisenberg chain), which used to be an explicit, rejected
scope limit here (see \texttt{pyitensor/idmrg\_excitations.py}'s own
module docstring, ``History'' section, for the eight-pass investigation
that limit came from and how it was eventually resolved by rewriting
the ansatz from scratch on top of VUMPS's own mixed-gauge state,
mirroring MPSKit.jl's own architecture). For \texttt{D=1} the computed
dispersion matches the exact free-fermion single-magnon dispersion of a
field-polarized XX chain to $\sim$14 digits across the whole Brillouin
zone (\texttt{examples/idmrg/excitation\_gap\_xx/main.py}); for
\texttt{D=2} it matches an independently-converged MPSKit.jl
transverse-field Ising result to 6 significant figures
(\texttt{examples/idmrg/excitation\_gap\_tfim/main.py}), and
$H_{\mathrm{eff}}(k)$ is Hermitian to machine precision at every
\texttt{D} tried. \texttt{itensor\_version=3} is cross-checked directly
against \texttt{itensor\_version="python"} across a full momentum scan,
matching to $\sim$$10^{-10}$ or tighter on the gapped TFIM case and to a
looser $\sim$$10^{-4}$--$10^{-7}$ on the gapless/critical Heisenberg
case (both backends' own non-convex VUMPS restart search can land on
slightly different local optima there, not a discrepancy in the ported
algorithm itself) -- see \texttt{tests/test\_vumps\_excitations\_v3.py}
and \texttt{examples/idmrg/vumps\_excitation\_v3\_VS\_python/main.py}. A
deliberately constructed longer-range term (spanning more than 2
adjacent unit cells after \texttt{n\_uc}-grouping) is rejected with
\texttt{NotImplementedError}, raised from \texttt{gs\_energy()} itself
(\texttt{vumps\_ground\_state}'s own reach-1 check) rather than from the
excitation machinery.

\textbf{Cost.} The eigenproblem \texttt{excitation\_energies} solves has
dimension $D^2(d_g-1)$, so its cost grows quickly with the converged bond
dimension. On \texttt{itensor\_version="python"} two things bound that:
the momentum-dependent channel resolvents are built once per momentum and
cached (they do not depend on the excitation tensor, so rebuilding them
per matrix element was pure waste), and above a size threshold the
eigenproblem is solved by Lanczos on the effective Hamiltonian's action
rather than by assembling that matrix at all. Neither changes any returned
number -- the two solver paths are cross-checked against each other in
\texttt{tests/test\_infinite\_chain.py} -- but a momentum scan at a large
$D$ is substantially cheaper than it was. \texttt{itensor\_version=3}
still assembles the matrix (and still rebuilds its resolvents per matrix
element), so the two backends now differ in cost structure for this call,
in the \texttt{"python"} backend's favour at large $D$.

\textbf{Spectral weights: $S(k,\omega)$ directly in the thermodynamic
limit.} Each branch the ansatz returns is a genuine momentum eigenstate,
so its contribution to a dynamical correlator is an exact $\delta$-peak
--- not a broadened, windowed approximation to one.
\texttt{ic.spectral\_weights(opname, k, p=0, n=1)} returns that peak's
position \emph{and} its residue:

\begin{verbatim}
ic.gs_method = "vumps"                        # required, as above
e, w = ic.spectral_weights("Sx", 0.7, n=1)    # energies and weights
e, w, total = ic.spectral_weights("Sx", 0.7, n=1, return_total=True)
ks, es, S = ic.dynamical_structure_factor("Sx", delta=0.05)
\end{verbatim}

with $w_a = |\langle k,a|O(k)|\Psi\rangle|^2$ for the normalized
quasiparticle state $|k,a\rangle$ and $O(k) = N^{-1/2}\sum_m e^{ikm}
O_m$, i.e.
\[
  S(k,\omega) \;=\; \sum_a w_a\,\delta(\omega - e_a).
\]
\texttt{opname} is a local operator on sub-site \texttt{p} of the unit
cell, and $k$ is per unit cell. This is the same physical object
\texttt{kpm\_finite}/\texttt{td\_dynamical\_correlator} estimate,
computed a completely different way: those embed a \emph{finite} window
in the infinite chain and therefore carry a window-size error plus a
KPM/time-truncation broadening, while these peaks are exact in both
momentum and energy. What they do not contain is multi-particle
continuum weight --- which is exactly what \texttt{return\_total}
measures, so the two families are complementary rather than one
superseding the other.
\texttt{ic.dynamical\_structure\_factor(opname, ks=None, energies=None,
delta=0.05, p=0, n=1)} is the convenience wrapper that assembles those
peaks into a Lorentzian-broadened $(k,\omega)$ grid ready to plot as a
heat map; the broadening there is cosmetic only.

\textbf{The sum rule, and what \texttt{return\_total} is for.} A one-site
operator applied to a uniform MPS lands \emph{exactly} inside the same
variational tangent space the excitations live in ($B = O\!\cdot\!A$ is
itself a valid excitation tensor), so the total weight summed over
\textbf{every} branch is exactly the per-site connected static structure
factor $\sum_r e^{ikr}(\langle O_0 O_r\rangle - \langle O\rangle^2)$.
\texttt{return\_total=True} returns that number --- free, since it is the
squared norm of the same source vector each weight is an inner product
against --- and \texttt{w.sum()/total} is then the fraction of the
momentum-resolved response the branches actually returned account for.
That is the practical test of whether a single-mode picture is adequate
at that momentum, and it is a sharp one: on the paramagnetic
transverse-field Ising chain, $\sigma^x$ (parity-odd, exactly one
quasiparticle) puts $>99.9\%$ of its weight on the lowest branch, while
$\sigma^z$ (parity-even) puts \emph{exactly zero} there --- measured at
$\sim$$10^{-21}$, an exact selection rule --- and all of it on higher
branches lying inside the exact two-particle continuum. See
\texttt{examples/idmrg/dynamical\_structure\_factor\_tfim/main.py}, which
plots all three.

The connectedness is automatic: no $\langle O\rangle$ is subtracted
anywhere, and none needs to be, because the disconnected piece lives
entirely along the gauge direction the left gauge-fixing projector
annihilates (see \texttt{pyitensor/idmrg\_excitations.py}'s
\texttt{\_spectral\_source\_vector}/\texttt{\_spectral\_resolvent}). The
$k=0$, $\langle\sigma^z\rangle=-1$ product-state case --- where the
disconnected part is 1, not small --- is pinned by a test.

\textbf{Validation and scope.} Four independent references, since a
spectral weight has no single golden number. The sharpest is the
\textbf{AKLT} point of the spin-1 chain ($H = S\!\cdot\!S +
\tfrac13(S\!\cdot\!S)^2$), whose ground state is \emph{exactly} a $D=2$
MPS --- so there is no variational error at all, and three closed forms
are reproduced to machine precision ($\sim$$5\times10^{-15}$): the static
structure factor summed from $\langle S_0\!\cdot\!S_r\rangle =
4(-1/3)^r$; Arovas--Auerbach--Haldane's single-mode dispersion
$\tfrac{5}{27}(5+3\cos k)$ [PRL \textbf{60}, 531 (1988)], which the
\emph{first moment} must reproduce exactly since $S^z_k|\Psi\rangle$ lies
entirely inside the tangent space; and the SU(2) content, which emerges
rather than being imposed --- the eight branches split into a magnon
triplet and a quintuplet at every momentum, $S^z$ reaches only the
triplet (quintuplet weight $\sim$$10^{-23}$), and the two multiplets
cross near $k\approx0.9$. On the pure \textbf{Haldane} chain the gap at
$k=\pi$ converges monotonically from below to the literature $0.4104789$
($0.2132/0.4074/0.4094/0.4098$ at $D=4/8/12/16$) while the sum rule holds
identically at every $D$ --- including $D=4$, where the gap is 50\%
wrong, which is the point: it is an identity of the ansatz, not of the
ground state's accuracy. The magnon triplet exhausts 97.3--97.6\% of the
$k=\pi$ sum rule and falls to $\sim$68\% at small $k$. See
\texttt{examples/idmrg/haldane\_structure\_factor/main.py}.

The other three, on spin-1/2: the exactly solvable $J=0$
Ising product state ($\sigma^x$/$\sigma^y$ weight exactly 1 at every
momentum, $\sigma^z$ exactly 0); the static sum rule above, checked
against a real-space sum of \texttt{ic.correlator} --- machinery sharing
no code with the excitation ansatz --- matching to $\sim$$10^{-13}$ at
$D=2$; and the f-sum rule $\sum_a e_a w_a = \tfrac12 \langle
[O_k^\dagger,[H,O_k]]\rangle$, which brings the excitation \emph{energies}
in too and converges with the ground state's own bond dimension
($2\times10^{-4}$ relative at $D=2$, $10^{-7}$ at $D=4$). See
\texttt{tests/test\_infinite\_chain\_spectral.py}. Requires
\texttt{itensor\_version="python"} --- unlike
\texttt{excitation\_energies}, the mpscpp3 port is energies-only and has
neither the eigenvectors nor the mixed-transfer source vector this needs,
deferred for the same reason it has not picked up the iterative
eigensolver (\texttt{docs/idmrg\_improvement\_plan.md}) --- plus
\texttt{gs\_method="vumps"} and \texttt{n\_uc} $\le 2$ as above. A
fermionic (parity-odd) \texttt{opname} is rejected: its Jordan-Wigner
string would have to be closed at infinity, the same reason
\texttt{correlator} rejects an odd-parity operator pair.

\textbf{Cost.} Independent of how many branches are asked for, and of
$D$, $d_g$ and the operator: exactly \textbf{two} linear solves per
momentum, both cached on the chain across repeated calls at the same $k$.
The naive way to build the weight --- leaving the excitation tensor's legs
open and solving once per basis element --- would cost $D^2 d_g$ solves
per momentum instead; transposing each geometric series moves its solve
to the operator end of the contraction instead. So the eigensolve for the
branches themselves dominates, and adding weights to an existing
dispersion scan is close to free.

\textbf{Unit cells of more than 2 sites.}
\texttt{Infinite\_Many\_Body\_Chain([...])} accepts a cell of any length,
and \texttt{gs\_method="vumps"} (the default) then runs the
\textbf{sequential multi-site} algorithm --- the state is a list of
per-site tensors \texttt{AL[n]}/\texttt{AR[n]}/\texttt{C[n]}, and one
iteration sweeps the cell solving a one-site eigenproblem
\texttt{H\_AC[n]} at each site and a zero-site \texttt{H\_C[n]} at each
bond (\texttt{pyitensor/vumps\_ms.py}). Nothing is grouped, so the cost is
\textbf{linear} in the cell size.

That distinction is the whole point. The older way to support a multi-site
cell is to fold it into one supersite of dimension $\prod_p d_p$ and run
the single-site algorithm --- exact, but exponential: a 4-site spinful cell
is $d_g = 256$. Nietner, Vanhecke, Verstraete, Eisert and Vanderstraeten
(arXiv:2003.01142) state it directly --- ``the cost of a naive application
of the VUMPS algorithm would scale exponentially with the size of the unit
cell'' --- and give as their key property ``a computational effort that
scales linearly rather than exponentially in the size of the unit cell''.
The sequential algorithm is the one production codes implement (TeNPy's
\texttt{SingleSiteVUMPSEngine}/\texttt{TwoSiteVUMPSEngine} sweep over the
MPS unit cell; MPSKit likewise), and it originates with the multi-site
VUMPS of Zauner-Stauber, Vanderstraeten, Fishman, Verstraete and Haegeman
(PRB 97, 045145 (2018), arXiv:1701.07035).

\texttt{vev} and \texttt{correlator} work at any \texttt{n\_uc}. Two things
do not, and say so rather than failing obscurely:
\texttt{gs\_method="idmrg"} still requires \texttt{n\_uc <= 2} (its growth
loop pairs sublattice $m$ with $n_{uc}-1-m$, which are only genuinely
adjacent for \texttt{n\_uc <= 2}), and so do
\texttt{excitation\_energies}/\texttt{excitation\_gap}, whose tangent-space
ansatz is still written against the grouped single-supersite gauge.
\texttt{n\_uc <= 2} keeps using the grouped VUMPS path, so its values are
unchanged; the two agree to machine precision where both apply.

\textbf{Product-state traps, and the noise that breaks them.} Every
particle-number-conserving Hamiltonian has the vacuum (and the filled state)
as an \emph{exact} eigenstate, and an exact eigenstate is an absorbing fixed
point of the growing algorithm: the local solve warm-started there returns it
immediately, its Schmidt rank is 1, so every subsequent truncation keeps a
single singular value and nothing grows the bond dimension back. The run then
reports a product state with \texttt{converged=True} in a fraction of a
second --- a silently wrong ground state, not a visible failure. Both
backends do this; it is a property of the algorithm, not of either port.

\texttt{ic.noise} (default \texttt{1e-4}) applies White's density-matrix
perturbation to break it, together with a matching random admixture on the
local solve's start vector --- both halves are needed, since enlarging the
basis alone leaves the solve pinned to an exact eigenvector it cannot leave.
The schedule is \textbf{demand-driven}: noise arms only while the state
genuinely is a product state (keyed on the purity of the noise-free reduced
state) and runs a few iterations past that, so an already-entangled model
never sees it and is bit-for-bit unaffected. \texttt{ic.noise\_iters}
(default 40) caps the total, so a model whose ground state genuinely
\emph{is} a product state --- a field-polarized chain --- stops re-arming it
and still ends on a clean tail; such models come back exact. Convergence is
never declared while noise is active. Set \texttt{ic.noise = 0} to disable
the mechanism entirely and recover the previous numerics exactly. See
\texttt{docs/known\_issue\_idmrg\_product\_state\_collapse.md} for the
measured before/after against an exact band integral.

\textbf{A cheaper, cruder alternative: the local superblock gap.}
\texttt{ic.local\_excitation\_gap(niter=200)} re-diagonalizes the growing
algorithm's own final, converged 2-site effective Hamiltonian -- already
solved once for the ground state -- for its \emph{second}-lowest
eigenvalue instead, and returns the difference. This is the direct
infinite-chain analogue of a well-known finite-DMRG trick: at the last
sweep, ask the same local effective Hamiltonian for its two lowest Ritz
pairs rather than just the ground state. It is also the natural place a
Lagrange-multiplier/orthogonality-penalty idea (as used by finite DMRG's
own dedicated excited-state method, \texttt{get\_excited}/\texttt{get\_gap},
\S4) would show up here -- except no penalty weight is needed: since there
is no separate re-sweep (the ``state to stay orthogonal to'' is just the
local ground vector already found, in the very same local Hilbert space),
the constraint is enforced exactly via deflation (projecting out the
ground vector), which is what a penalty method converges to anyway as its
weight $\to\infty$. Unlike \texttt{excitation\_gap}, this requires
\texttt{gs\_method="idmrg"} (not \texttt{"vumps"}, the default) and
carries no momentum label -- it is a single number, not a dispersion:

\begin{verbatim}
ic.local_excitation_gap()   # a single scalar, no k dependence
\end{verbatim}

\textbf{Accuracy caveat.} This is a genuinely cruder notion of ``gap'':
it reuses the ground state's own \texttt{HL}/\texttt{HR} environments
unmodified (never letting them relax for whatever the second local
eigenstate actually represents), so it is \emph{not} guaranteed to match
the true minimum-momentum gap the tangent-space ansatz targets. Measured
directly on the two cases \texttt{excitation\_gap} and this method were
both cross-checked against: for the exactly-solvable field-polarized XX
chain (\texttt{D=1}), it comes out $\sim$10\% too high (5.5 vs.\ the exact
5.0); for a genuinely entangled (\texttt{D>1}) dimerized Heisenberg chain,
it lands within $\sim$0.5\% of a large finite open chain's own
extrapolated ED gap (see
\texttt{examples/idmrg/local\_excitation\_gap/main.py}). Prefer
\texttt{excitation\_gap} for a physically principled answer (it now
handles any \texttt{D}, see above); \texttt{local\_excitation\_gap}
remains useful as a cheap, order-of-magnitude cross-check, or when
\texttt{gs\_method="idmrg"} is what's already been run for other reasons
(\texttt{vev}/\texttt{correlator}). \textbf{Both of those calibrations were
measured on spin models}, where the charge channel described below does not
exist at all; on a fermionic model with \texttt{ConserveQNs=false} there is
no reason those error bars carry over, and the estimate should be
cross-checked against a finite chain, the correlator decay, or an $n(\mu)$
plateau before it is relied on.

\textbf{Why the excited solve shifts rather than only projects.} The second
eigenvalue is the lowest one of the \emph{same} stored operator restricted
to its ground state's orthogonal complement. Writing that as the bare
projector $P = I - |\psi_0\rangle\langle\psi_0|$ and diagonalizing $PHP$ is
wrong, and used to make this method return large \textbf{negative} ``gaps'':
$PHP$ agrees with $H$ on the complement, but it also carries $\psi_0$ itself
as an eigenvector with eigenvalue \emph{exactly zero}. That is harmless only
while the rest of the spectrum lies below zero --- the moment the stored
operator's own ground eigenvalue is \textbf{positive}, zero becomes $PHP$'s
smallest eigenvalue and is exactly what a smallest-eigenvalue solver is
asked to find (deflation keeps $\psi_0$ out of the Krylov space in exact
arithmetic only; rounding regrows the component and a restarted solver locks
onto it). The reported gap is then $0 - e_0$, i.e.\ precisely minus the
stored superblock energy --- reported from a spinful (c,f) Kondo chain at
\texttt{maxm=16} as $-358.003$ meV against a stored energy of
$+0.358002587681$, digit for digit. The sign of that stored energy is
nobody's choice: the per-site energy baseline leaves it as a small residual
boundary term of either sign, which is why the same model was right at
\texttt{maxm=8} (energy $-1.05$, negative) and wrong at \texttt{maxm=16}.
Both backends now use $PHP + \sigma|\psi_0\rangle\langle\psi_0|$ with
$\sigma$ far above the bottom of the spectrum, so $\psi_0$'s own eigenvalue
can never be the answer while the complement is untouched. Both eigenvalues
are also re-solved from the stored superblock at the same solver strength
rather than the ground one being read back from the growing algorithm, and
if the re-solve lands \emph{below} the growing algorithm's own value a
\texttt{RuntimeWarning} is raised --- the returned number is then still a
genuine spectral gap of the stored operator, but no longer a gap measured
above the state whose observables \texttt{vev}/\texttt{correlator} report.

\textbf{Charge excitations, and why a uniform on-site shift moves this
number.} Both backends build every site with \texttt{ConserveQNs=false}
(\texttt{mpscpp3/get\_sites.h}; \texttt{pyitensor} likewise), so the local
superblock spectrum contains particle-number-\emph{changing} excitations
against the frozen environment, with nothing to confine the deflated solve
to the ground state's own charge sector. On a fermionic model the gap can
therefore be set by a charge excitation --- and adding a constant
\texttt{mu} to every on-site energy, which in a gapped plateau leaves the
converged state and all its correlators untouched, genuinely changes the
answer. Measured exactly (by dense diagonalization of the stored
superblock) on a gapped SSH spinless-fermion chain at \texttt{maxm=16}: the
gap is 0.631321 at \texttt{mu=0} and 0.431325 / 0.431324 at
\texttt{mu=$\pm$0.2}, i.e.\ exactly $|\mathtt{mu}|$ lower, because the two
$\pm 1$-electron excitations that are degenerate at \texttt{mu=0} split;
the stored ground vector is the exact ground state at every point (overlap
1.0000), so no solver is at fault. This is a property of the estimator, not
a defect to fix: ``the lowest state orthogonal to the ground state'' simply
is not a fixed-charge quantity here. Read it as an order-of-magnitude
cross-check, and cross-check a fermionic gap against a finite-chain
calculation, the correlator decay, or an $n(\mu)$ plateau before relying on
it.

\textbf{Tightening it further: \texttt{window=}.}
\texttt{ic.local\_excitation\_gap(window=w, niter=200)} (default
\texttt{w=0}, exactly the behavior above) grows the local diagonalization
block by \texttt{w} extra \emph{free} physical sites on each side of the
original 2, re-solving both the ground state and the deflated first
excited state fresh within this larger block (rather than reusing the
growing algorithm's own ground vector) -- i.e.\ it lets the excitation
spread across more of the chain instead of only ever living on the same
frozen 2 sites, at the cost of an exponentially larger local Hilbert space
(\texttt{d**(2*w)} more, \texttt{d} = the physical dimension). This is
\emph{not} the same thing as throwing more Krylov vectors at the original
2-site problem -- \texttt{local\_excitation\_gap} solves that one to Krylov
convergence already, so more iterations there do not change the answer
(with the caveat above that the \emph{ground} eigenvalue it deflates
against must be the right one, which is why both are re-solved together). Measured directly: on the field-polarized XX chain the error drops
from 10\% at \texttt{w=0} to 3.8\%/2.0\%/0.81\%/0.44\% at
\texttt{w=1/2/4/6}; on a gapped, genuinely entangled (\texttt{D>1})
transverse-field Ising chain, \texttt{w=3} (an 8-site local block) matches
an 18-site open finite chain's own ED gap to $<$1\%, converging at least
as fast as growing the finite chain itself does. The improvement is not
free, and its \emph{rate} depends on the model's correlation length: on
the \texttt{S=1} Heisenberg chain (the Haldane gap, correlation length
$\sim$6 sites), \texttt{w=0/1/2} only move the estimate from
$\sim$29\%/24\%/21\% too high -- still improving, but far more slowly,
since a handful of extra sites barely dents a correlation length that
long, and the physical dimension \texttt{d=3} makes each extra site pair
9$\times$ more expensive rather than XX/TFIM's 4$\times$. Only
\texttt{n\_uc=1} is supported for \texttt{w>0} -- widening needs to know
which sublattice position each extra site takes, which is not tracked for
\texttt{n\_uc=2} yet (raises \texttt{NotImplementedError}).
\texttt{w>0} is also \texttt{itensor\_version="python"}-only: it is a
prototype rather than stable API, so it was deliberately left out of the
v3 C++ port (which does cover \texttt{w=0}).

\textbf{Dynamical correlators (finite-window KPM).}
\texttt{ic.kpm\_finite(opname\_i, p\_i, opname\_j, r,
n\_window, window\_chain\_kwargs=None, **kwargs)} computes
$\langle \text{opname\_i}(\text{site } p_i)\,
\text{opname\_j}(\text{site } p_i+r)\rangle(\omega)$ of the infinite
chain, reusing the existing finite-chain KPM machinery
(\texttt{kpmdmrg.get\_dynamical\_correlator}) unmodified rather than a
new Chebyshev-recursion implementation. Named \texttt{kpm\_finite} (not
\texttt{get\_dynamical\_correlator}, the finite-chain method it wraps)
to flag up front that it is the finite-window \emph{approximation}
described below, not an exact infinite-size calculation. There is no
infinite-chain analogue of \texttt{vev}/\texttt{correlator}'s
transfer-matrix formalism for a \emph{dynamical} quantity -- the
Hamiltonian is extensive/unbounded in the thermodynamic limit, so a
literal Chebyshev expansion of the full $H$ has no meaning (unlike
\texttt{apply\_mpo}/\texttt{imps\_sum}'s bounded-operator scope above).
Instead, this method builds an ordinary finite, open-boundary chain of
\texttt{n\_window} repeats of the unit cell (Hamiltonian =
\texttt{h\_intra} tiled onto every cell plus \texttt{h\_inter} tiled
onto every adjacent pair of cells -- one bond fewer than a periodic ring
at the two open ends), places the two operators at the window's
\emph{central} unit cell (as far as possible from both ends), and
delegates directly to the same KPM code path an ordinary finite
\texttt{Spin\_Chain}/\texttt{Fermionic\_Chain} uses:

\begin{lstlisting}
es, ys = ic.kpm_finite("Sz", 0, "Sz", 0, n_window=16,
        window_chain_kwargs=dict(maxm=30, nsweeps=15, kpmmaxm=60),
        delta=0.3, es=np.linspace(-1, 5, 200))
\end{lstlisting}

\texttt{n\_window} has no default -- see the convergence caveat below.
\texttt{window\_chain\_kwargs} is an optional dict of attribute
overrides applied to the temporary finite \texttt{Many\_Body\_Chain}
(\texttt{maxm}, \texttt{nsweeps}, \texttt{kpmmaxm}, \texttt{kpm\_scale},
...), independent of \texttt{ic}'s own \texttt{maxm}/etc.; remaining
\texttt{**kwargs} (\texttt{delta}, \texttt{kernel}, \texttt{es},
...) are forwarded to
\texttt{kpmdmrg.get\_dynamical\_correlator} unchanged.

\textbf{Scope restriction -- a finite-window approximation, read before
use.} This is \emph{not} an exact infinite-size method: results carry
finite-size/open-boundary corrections that must be checked by
convergence in \texttt{n\_window}, exactly as a static
\texttt{vev}/\texttt{correlator} caller would check
\texttt{maxm}/\texttt{etol} convergence of the original iDMRG ground
state. One Chebyshev moment corresponds to one application of the
(nearest-neighbor) window Hamiltonian, so it can only move information
by $\sim$1 site per moment (a Lieb-Robinson-style bound) -- but KPM's
own moment count scales with the \emph{window's own extensive
bandwidth} divided by the requested \texttt{delta} (an ordinary finite
chain's KPM already has this property, nothing new here), so a
genuinely fine \texttt{delta} can require a moment count comparable to
(or larger than) \texttt{n\_window} itself, at which point
open-boundary reflections contaminate the result regardless of how
large \texttt{n\_window} is. Prefer a coarser \texttt{delta}, or check
that the correlator has visibly converged with growing
\texttt{n\_window}, for quantitative work (especially near a gapless
point, where a fine \texttt{delta} is most tempting). Unlike
\texttt{vev}/\texttt{correlator}, this does not need \texttt{ic.\_result}
(no dependency on a previously converged \texttt{IDMRGResult}, or even
on \texttt{ic.itensor\_version}), so it works regardless of which
backend \texttt{gs\_energy()} itself used. See
\texttt{examples/idmrg/dynamical\_correlator\_finite\_window/main.py}
for a worked example sweeping \texttt{n\_window}.

\textbf{A genuinely infinite-chain dynamical correlator, via real-time
TDVP (\texttt{td\_dynamical\_correlator}).} \texttt{kpm\_finite}'s own
open-boundary window has a real error source no amount of
\texttt{n\_window} alone can fix: an open chain's own ground state
carries boundary artifacts (e.g. Friedel-oscillation-like features)
that contaminate even the \emph{central} region, not just the two
edges. \texttt{ic.td\_dynamical\_correlator(opname\_i, p\_i, opname\_j,
n\_window, dt=0.1, nt=200, x\_values=None, maxdim=60, cutoff=1e-10,
niter=50, connected=True, **kwargs)} fixes this by capping the window's
two ends with the \emph{converged} iDMRG growth environment
(\texttt{idmrg\_ground\_state}'s own \texttt{HL}/\texttt{HR}, already
computed during growth and exposed on \texttt{IDMRGResult}) instead of
plain open boundaries -- infinite boundary conditions (IBC), following
Milsted/Vanderstraeten et al., ``Infinite boundary conditions for
response functions and limit cycles in iDMRG'' (arXiv:1804.09163) --
and evolves the perturbed window in real time via two-site TDVP rather
than expanding in Chebyshev moments. Supports both
\texttt{itensor\_version="python"} and \texttt{itensor\_version=3}
(calls \texttt{gs\_energy()} automatically if needed, like
\texttt{vev}/\texttt{correlator}, though \texttt{vev}/\texttt{correlator}
themselves still require \texttt{itensor\_version="python"} regardless):

\begin{lstlisting}
ks, es, Skw = ic.td_dynamical_correlator(
        "Sz", 0, "Sz", n_window=14, dt=0.05, nt=200,
        maxdim=60, cutoff=1e-10, x_values=range(-6, 7),
        ks=np.linspace(-np.pi, np.pi, 41), delta=0.1, window=[-1, 6])
\end{lstlisting}

The native ITensor v3 backend (\texttt{Chain::td\_dynamical\_correlator\_window},
\texttt{mpscpp3/chain\_session.h}) reuses the vendored ITensorTDVP
library's own boundary-tensor \texttt{tdvp(psi,H,t,LH,RH,sweeps,args)}
overload directly against a tiled window MPS/MPO -- unlike the
\texttt{"python"} backend, which has to hand-roll its own window-aware
TDVP sweep (pyitensor's generic TDVP infers a site's Link via a
same-Index chain-neighbor lookup that cannot see a window's extra
boundary legs). Two scope differences versus \texttt{"python"}: (1)
\texttt{x\_values} may not extend beyond the window's own explicit range
(\texttt{center+x} must stay within the window, i.e. increase
\texttt{n\_window} instead of relying on padding) -- the
\texttt{"python"} backend pads beyond the window with extra unevolved
unit-cell copies, not ported here; (2) as of this writing,
\texttt{itensor\_version=3}'s own \texttt{idmrg\_ground\_state} has a
known, pre-existing convergence bug for Hamiltonians with an onsite
(``field'') term (energy diverging every macro-iteration) -- unrelated
to \texttt{td\_dynamical\_correlator} itself, but it means a v3
\texttt{td\_dynamical\_correlator} call inherits that limitation for
such models; a purely bond-coupled Hamiltonian (e.g. plain Heisenberg)
is unaffected.

\texttt{opname\_j} is applied at sublattice position \texttt{p\_i} and
evolved forward in time under the window's own ground-state-energy-
shifted Hamiltonian ($e^{-i(H-E_{\rm GS})t}$, matching this codebase's
own established real-time correlator convention, e.g.\
\texttt{mpscpp3::quench\_tdvp}'s $H_{\rm shift}=H-E_{\rm GS}\cdot
\text{Id}$); \texttt{opname\_i} is inserted at the shifted
position (bra side, not itself evolved) -- this is the paper's own
headline efficiency result: every \texttt{x} (and, via the spatial
Fourier transform below, every \texttt{k}) comes from this \emph{one}
window evolution, not one run per distance the way a naive real-time
approach (or \texttt{kpm\_finite}'s own one-run-per-\texttt{r} KPM
calls) would need. Cross-checked against an exact non-interacting
(free-fermion) reference for the XX chain, on systems far larger than
many-body ED could reach -- see \texttt{docs/documentation.tex}'s own
architecture-level notes on this method for the full derivation and
what that check found. \texttt{connected=True} (default) subtracts the
disconnected background \texttt{<opname\_i><opname\_j>} before the
spatial Fourier transform -- turning it off produces a spurious,
dominant $k=0$ contribution with no discernible dispersion (the raw
correlator approaches \texttt{<opname\_i><opname\_j>}, not $0$, at
large separation). \texttt{Skw} (shape \texttt{(len(ks), len(es))}) is
obtained via a spatial DFT ($S(k,t)=\sum_x e^{-ikx}S(x,t)$) followed by
the \emph{same} damping/FFT convention (\texttt{delta} ->
Lorentzian broadening) every other dynamical-correlator submode in this
codebase already uses.

\textbf{Fermionic operators, and the vacuum normalization} (both since
2026-08-29). A parity-odd pair (\texttt{"Cdag"}/\texttt{"C"}, either
order, on any fermionic site type) computes the \emph{physical} fermionic
correlator: the Jordan-Wigner string is threaded across the window on the
ket (before the evolution, so it is evolved along with the perturbation)
and across the bra at measurement time, matching the convention
\texttt{correlator} already used for the static case. The two are pinned
to each other by an exact identity --- $S(x,t{=}0)$ equals
\texttt{correlator(...)} to machine precision at every $x$, both signs,
across the window's own padding --- and against the exact free-fermion
Green function $\langle c^\dagger_x(t) c_0\rangle = \sum_l [e^{iht}]_{xl}
P[l,0]$ at $t>0$ (\texttt{tests/test\_idmrg\_window\_fermionic.py},
\texttt{examples/idmrg/fermionic\_dynamical\_correlator/main.py}). Note
the anticommutation sign this implies: for two parity-odd operators at
different sites $\langle A_x B_0\rangle = -\langle B_0 A_x\rangle$, so a
fermionic $S(x,t)$ and a static \texttt{correlator} written in the
opposite site order differ by a minus sign. A pair with odd \emph{total}
parity (one fermionic operator against a parity-even one) raises rather
than returning a number, exactly as \texttt{correlator} does: its string
can never close. \texttt{connected=True} subtracts nothing for a
parity-odd pair, whose disconnected background is zero by symmetry.
Before this, the path applied a bare \texttt{C}/\texttt{Cdag} matrix with
no string on either backend --- a different number entirely, not a less
converged one (measured: $+0.203$ at $x=2$ against an exact $-0.001$).
Independently, every $S(x,t)$ --- spin included --- is now divided by the
vacuum amplitude $\langle\psi|\psi(t)\rangle$ measured through the
identical contraction, which cancels a spurious global factor that had
been inflating results on both backends (the shift \texttt{eshift} that
keeps the evolution phase-stationary is measured with the window's
boundary legs traced, while correlators are measured with them closed by
the transfer-matrix fixed points, and the two see different energies). On
the dimerized XX chain that alone took the residual against the exact
free-fermion answer from $\sim$0.07 to $\sim$1e-5.

\textbf{Both backends are exact on the $S(x,0)$ identity.} $S(x,t{=}0)$
must equal the chain's own static \texttt{correlator(...)} exactly (no
evolution has happened yet), and both \texttt{itensor\_version="python"}
and \texttt{itensor\_version=3} satisfy it to $\sim$1e-15. That was not
always true: until 2026-08-29 the v3 window tiled the raw per-micro-step
iDMRG factors rather than the gauge-consistent unit cell every other v3
static observable uses, and missed the identity by up to $\sim$1.7e-1 on
a plain Heisenberg chain whose energy density the two backends agreed on
to 6.7e-11 --- shape and finiteness looked fine throughout, which is how
it went unnoticed. Fixing it also removed this backend's
$\exp(+i\,\texttt{eshift}\,t)$ phase correction in favour of the exact
vacuum normalization the Python backend already used, and a latent
index-stride bug that only showed up when the converged unit cell is not
square. See
\texttt{examples/idmrg/td\_dynamical\_correlator\_python\_VS\_v3/main.py},
which checks the identity at every $x$ on both backends and then compares
the two $S(x,t)$ trajectories.

\textbf{Scope}: this simplifies the paper's own Eq.~7 to $t_1=0$ (the
ground state is perturbed by \texttt{opname\_j} and evolved forward
only; \texttt{opname\_i} is never itself time-evolved) rather than the
full two-branch trick (evolving a \emph{second}, independent window
backward in time too, which doubles the accessible total time for the
same TDVP cost) -- a documented, straightforward follow-up.
\texttt{n\_window} and \texttt{x\_values} are two \emph{separate}
convergence axes to check (not one): \texttt{n\_window} controls how
much environment margin surrounds the perturbation, \texttt{x\_values}
how far the spatial sum/Fourier transform reaches -- growing
\texttt{x\_values} together with \texttt{n\_window} can keep changing
results (a slowly decaying connected-correlator tail keeps contributing
as the range widens), so converge each independently. See
\texttt{examples/idmrg/td\_dynamical\_correlator/main.py} for a worked
\texttt{n\_window}-convergence sweep and a cross-check against
\texttt{kpm\_finite}.

\textbf{Applying an operator/gate to the converged chain (advanced).}
\texttt{pyitensor.idmrg.apply\_mpo(result, W\_bulk, cutoff=..., maxdim=...)}
is the infinite-chain analogue of the finite backends' \texttt{applyMPO}:
it contracts a periodic MPO onto every site of the converged unit cell
and re-canonicalizes/truncates the grown bond dimension back down via
the standard two-sided fixed-point infinite-MPS canonicalization
procedure, returning a new \texttt{pyitensor.idmrg.PeriodicMPS} that
\texttt{onsite\_expectation}/\texttt{two\_point\_correlator} accept
exactly like an \texttt{IDMRGResult}. There is no
\texttt{Infinite\_Many\_Body\_Chain}-level wrapper yet, so it is reached
by working with \texttt{pyitensor.idmrg} directly against
\texttt{ic.\_result} (only meaningful for
\texttt{itensor\_version="python"}, same restriction as
\texttt{vev}/\texttt{correlator}):

\begin{lstlisting}
from dmrgpy.pyitensor import idmrg
from dmrgpy.pyitensor.index import Index
from dmrgpy.pyitensor.tensor import ITensor

sites_uc = ic._result.sites_uc
d = sites_uc.dim(1)
pauli_x = 2 * sites_uc.site_type(1).matrix("Sx")           # a chi_W=1 operator
link_l, link_r = Index(1, tags="Link"), Index(1, tags="Link")
s = sites_uc.si(1)
W = [ITensor((link_l, s, s.prime(1), link_r), pauli_x.reshape(1, d, d, 1))]

new_state = idmrg.apply_mpo(ic._result, W, cutoff=1e-12, maxdim=None)
idmrg.onsite_expectation(new_state, "Sz", 0)               # <Sz> of the new state
\end{lstlisting}

\textbf{Scope restriction -- bounded operators only.} \texttt{W\_bulk}
must represent a \emph{bounded} (non-extensive) periodic operator: the
same tensor reused at every unit cell, with no unconditional ``keep
accumulating forever'' self-loop -- single-site products (as above),
gates tiled once per unit cell (an SVD-split 2-site gate embedded with
identity everywhere else), symmetry operators, and the like.
\texttt{pyitensor/idmrg.py}'s own Hamiltonian automaton (built internally
by \texttt{\_build\_periodic\_mpo} for \texttt{gs\_energy()}) is
deliberately the \emph{other} kind -- its accumulator channel needs a
genuine chain boundary to correctly represent an extensive sum, so
feeding it into \texttt{apply\_mpo}'s boundary-less periodic contraction
does not compute ``H\textbar psi\textrangle'' (see
\texttt{pyitensor/idmrg.py}'s own ``Applying a (bounded) MPO to the
converged iMPS'' section docstring for the specific failure mode). This
makes \texttt{apply\_mpo} the natural building block for e.g.\ an
iTEBD-style real/imaginary-time evolution step (repeatedly apply a local
Trotter gate + truncate) -- not yet implemented as a public feature, but
\texttt{apply\_mpo} is exactly the primitive it would use.

\texttt{W\_bulk}'s physical Indices must be the \emph{same objects} as
\texttt{result.sites\_uc}'s own (\texttt{sites\_uc.si(i+1)}), so a custom
operator automaton must be built against \texttt{result.sites\_uc}
directly (as in the example above), not a freshly constructed
\texttt{SiteX}. The truncation/regauging step's own numerical
conditioning degrades the further the \emph{raw} grown bond dimension
(\texttt{chi\_A * chi\_W}, before any truncation) exceeds the state's
real entanglement at that cut -- keep \texttt{maxm}/\texttt{maxdim}
modest relative to \texttt{chi\_W} for the best-conditioned result.

\textbf{Overlap/fidelity between two converged infinite MPS (advanced).}
\texttt{pyitensor.idmrg.imps\_overlap(result\_a, result\_b,
normalize=True)} computes the per-unit-cell overlap between two
\texttt{IDMRGResult}/\texttt{PeriodicMPS} objects -- the infinite-chain
notion of a finite MPS inner product $\langle\phi|\psi\rangle$. A literal
$\langle\phi|\psi\rangle$ over an infinite chain is not, in general, a
finite number (it scales as $\eta^N$ over $N$ unit cells, $N \to \infty$,
where $\eta$ is the dominant eigenvalue of the mixed transfer matrix
built from the two states' own tensors), so by default
\texttt{imps\_overlap} returns the always-finite \emph{per-site fidelity}
instead -- magnitude 1 iff the two iMPS represent the same physical state
(any gauge or normalization convention on the raw tensors), magnitude
$<1$ otherwise:

\begin{lstlisting}
from dmrgpy.pyitensor import idmrg

idmrg.imps_overlap(ic._result, ic._result)          # 1 -- same state
idmrg.imps_overlap(ic._result, new_state)            # cross-check apply_mpo didn't change the physical state
idmrg.imps_overlap(ic._result, some_other_result)    # < 1 in magnitude for a genuinely different state
\end{lstlisting}

\texttt{result\_a}/\texttt{result\_b} must share the same
\texttt{n\_uc} and, at every sublattice position, the same local
physical dimension -- \texttt{imps\_overlap} raises \texttt{ValueError}
otherwise. The two states' bond dimensions need \emph{not} match (e.g.\
comparing a ground state against an \texttt{apply\_mpo} output truncated
to a different \texttt{maxdim}). Pass \texttt{normalize=False} for the
raw, un-normalized mixed-transfer eigenvalue instead -- mainly a
diagnostic, analogous to \texttt{apply\_mpo}'s own returned
\texttt{.eta}.

\textbf{Direct sum of two converged infinite MPS (advanced).}
\texttt{pyitensor.idmrg.imps\_sum(result\_a, result\_b, cutoff=...,
maxdim=...)} is the periodic-chain analogue of the finite backends'
\texttt{mpsalgebra.sum}: a block-diagonal-in-the-bond-space direct sum at
every unit-cell cut (there is no open boundary to instead concatenate
along the way a finite chain's two ends do), re-canonicalized/truncated
via the same \texttt{apply\_mpo}-style two-sided fixed-point procedure,
returning a new \texttt{PeriodicMPS}:

\begin{lstlisting}
from dmrgpy.pyitensor import idmrg

result = idmrg.imps_sum(result_a, result_b, cutoff=1e-12, maxdim=None)
\end{lstlisting}

\textbf{Scope restriction -- read before use.} Tiled to the
thermodynamic limit, this ``+'' does \emph{not} represent a literal
Hilbert-space vector sum the way the finite-chain construction does. It
is only well-posed (a single dominant branch surviving -- the
mathematically correct infinite-volume answer, not a truncation
artifact) when \texttt{result\_a}/\texttt{result\_b} have a genuine
per-site norm mismatch, i.e.\ different self-overlap transfer eigenvalues
(\texttt{eta}, in the sense of \texttt{imps\_overlap}'s own
\texttt{eta\_aa}/\texttt{eta\_bb}). This mismatch never happens between
two \emph{ordinary} \texttt{IDMRGResult}s -- every one is individually
normalized to \texttt{eta=1} exactly by left-canonical SVD construction
-- so summing two separately-converged ground states (the most natural
reason to want this operation, e.g.\ to combine two symmetry-related
solutions of the same Hamiltonian) hits a genuine tie every time. That
tied case is a ``cat state'' superposition of two macroscopically
distinct branches, which this module's single-fixed-point machinery
cannot represent as one canonical periodic MPS -- \texttt{imps\_sum}
raises \texttt{RuntimeError} there (via \texttt{pyitensor.idmrg}'s
internal degeneracy check on the combined transfer matrix's dominant
eigenvalue) rather than silently collapsing to one arbitrary branch.
Correctly evaluating observables on that common, physically meaningful
case (e.g.\ via the thermodynamic-limit local-observable identity for
two orthogonal, equally-weighted branches) needs correlator machinery
this module does not have yet -- a documented, deliberate scope limit,
not something silently gotten wrong. See \texttt{pyitensor/idmrg.py}'s
own ``Summing two converged iMPS'' section docstring for the full
derivation and empirical confirmation.

\textbf{Direct sum of two converged VUMPS iMPS (advanced).}
\texttt{pyitensor.vumps.imps\_sum(result\_a, result\_b, cutoff=...,
maxdim=...)} is the VUMPS-mixed-gauge analogue of \texttt{idmrg.imps\_sum}
immediately above -- same construction in spirit (block-diagonal direct
sum, re-canonicalized/truncated via \texttt{idmrg.\_canonicalize\_periodic}),
but working on \texttt{VUMPSResult}'s own grouped-supersite \texttt{AL}
tensor (bond dimension \texttt{D\_a+D\_b}) rather than
\texttt{idmrg.py}'s per-sublattice periodic \texttt{U\_list}, and
returning a new \texttt{pyitensor.vumps.UniformMPS} (accepted directly by
\texttt{vumps.onsite\_expectation}/\texttt{two\_point\_correlator}, same
duck typing as \texttt{idmrg.PeriodicMPS}) whose full mixed gauge
$\{A_L,A_R,C,A_C\}$ is then reconstructed via
\texttt{vumps.\_complete\_mixed\_gauge} -- the standard ``bringing a
uniform MPS to canonical form'' procedure (Vanderstraeten, Haegeman,
Verstraete, ``Tangent-space methods for uniform matrix product states'',
arXiv:1810.07006, Eq.(9)-(17): factor the truncated $A_L$'s own dominant
right transfer-matrix fixed point $r=CC^\dagger$, then
$A_R := C^{-1}A_L C$ and $A_C := A_L C$):

\begin{lstlisting}
from dmrgpy.pyitensor import vumps

result = vumps.imps_sum(result_a, result_b, cutoff=1e-12, maxdim=None)
\end{lstlisting}

Same physical scope restriction as \texttt{idmrg.imps\_sum} above, for the
same reason: every converged \texttt{VUMPSResult} has $A_L$ exactly
left-canonical \emph{and} $A_R$ exactly right-canonical by construction of
the mixed gauge, so both its left and right self-overlap transfer
eigenvalues are exactly \texttt{eta=1} -- summing two ordinary
\texttt{VUMPSResult}s therefore always hits the same
degenerate-dominant-eigenvalue tie \texttt{idmrg.imps\_sum} does, and
\texttt{imps\_sum} raises \texttt{RuntimeError} there rather than silently
returning one arbitrary branch (arXiv:1810.07006's own Sec.~2.1
independently makes the same point: non-injective/``cat state'' MPS
tensors are exactly the ones with a degenerate dominant transfer-matrix
eigenvalue). Only two states with a genuine per-site norm mismatch (e.g.\
one deliberately rescaled, exactly mirroring \texttt{idmrg.imps\_sum}'s
own worked example) have a well-posed sum -- see
\texttt{pyitensor/vumps.py}'s own ``Summing two converged VUMPS iMPS''
section docstring for the full derivation, and
\texttt{examples/idmrg/vumps\_imps\_sum/main.py} for a worked example of
both cases plus a cross-check of the surviving branch's
\texttt{onsite\_expectation}/\texttt{two\_point\_correlator} at genuinely
entangled $D>1$.

Also ported to \texttt{itensor\_version=3} (\texttt{Chain::vumps\_imps\_sum},
\texttt{mpscpp3/chain\_session.h}), a dense-array translation reusing the
same \texttt{vx\_canonicalize\_n1}/\texttt{vumps\_complete\_mixed\_gauge}
machinery \texttt{apply\_mpo} below shares -- same scope restriction and
degeneracy behavior as the pyitensor version above. There is no
\texttt{Infinite\_Many\_Body\_Chain}-level wrapper on this backend either
(matching pyitensor's own scope), so it is reached directly against a
Chain's own converged VUMPS snapshot:

\begin{lstlisting}
D, d_g, AL, AR, C, AC, eta = ic._session3.vumps_imps_sum(D_b, AL_b, cutoff=1e-12, maxdim=0)
ic._session3.vumps_load_uniform_state(D, d_g, AL.flatten().tolist(),
                                       AR.flatten().tolist(), C.flatten().tolist())
\end{lstlisting}

\texttt{AL\_b} is the second state's own flat $(D_b,d_g,D_b)$ $A_L$ array
(row-major) -- e.g.\ obtained from another Chain's own snapshot via
\texttt{Chain::vumps\_get\_snapshot()} (which returns
$(D,d_g,A_L,A_R,C)$ for \texttt{this} Chain's own current state);
\texttt{maxdim<=0} means no cap (Python's \texttt{maxdim=None}).
\texttt{vumps\_load\_uniform\_state} writes the result back into a
Chain's own snapshot so \texttt{vumps\_onsite\_expectation}/
\texttt{vumps\_two\_point\_correlator} (and, since
\texttt{ic.\_session3\_has\_vumps} is untouched by any of this,
\texttt{ic.vev}/\texttt{ic.correlator} too) see it. See
\texttt{tests/test\_vumps\_imps\_sum\_v3.py} for the full cross-check
against \texttt{itensor\_version="python"}.

\textbf{Applying an operator/gate to a converged VUMPS iMPS (advanced).}
\texttt{pyitensor.vumps.apply\_mpo(result, W\_bulk, cutoff=...,
maxdim=...)} is the VUMPS-mixed-gauge analogue of \texttt{idmrg.apply\_mpo}
above: it groups \texttt{W\_bulk} into a single grouped-supersite MPO
tensor (the same \texttt{\_group\_automaton} routine that groups VUMPS's
own Hamiltonian automaton), grows the converged $A_L$ by it via
\texttt{idmrg.grow\_by\_mpo}, re-canonicalizes/truncates via
\texttt{idmrg.\_canonicalize\_periodic} (the identical two-sided
fixed-point procedure \texttt{idmrg.apply\_mpo} uses), and completes the
resulting truncated left-canonical tensor to the full mixed gauge
$\{A_L,A_R,C,A_C\}$ via \texttt{vumps.\_complete\_mixed\_gauge} (the same
completion \texttt{imps\_sum} immediately above uses). \texttt{W\_bulk}
takes \emph{exactly} \texttt{idmrg.apply\_mpo}'s own convention -- a list
of \texttt{n\_uc} rank-4 (Left, in, out, Right) ITensors, one per
unit-cell sublattice site -- so the identical \texttt{W\_bulk} list built
for one backend can be fed to the other's own \texttt{apply\_mpo} to
cross-check both against each other on the same operator:

\begin{lstlisting}
from dmrgpy.pyitensor import vumps

new_state = vumps.apply_mpo(result, W, cutoff=1e-12, maxdim=None)
\end{lstlisting}

Returns a new \texttt{pyitensor.vumps.UniformMPS}, accepted directly by
\texttt{vumps.onsite\_expectation}/\texttt{two\_point\_correlator} exactly
like a \texttt{VUMPSResult}. Same scope restriction as
\texttt{idmrg.apply\_mpo} above -- \texttt{W\_bulk} must represent a
\emph{bounded} (non-extensive) periodic operator; the Hamiltonian's own
automaton (built internally by \texttt{vumps\_ground\_state} for
\texttt{gs\_energy()}) is out of scope, for the identical reason described
there. There is no \texttt{Infinite\_Many\_Body\_Chain}-level wrapper yet
either, so \texttt{pyitensor.vumps.apply\_mpo} is reached directly,
working against \texttt{ic.\_vumps\_result} (only meaningful for
\texttt{itensor\_version="python"} with \texttt{gs\_method="vumps"}). See
\texttt{examples/idmrg/vumps\_apply\_mpo/main.py} for a worked example: a
\texttt{chi\_W=1} unitary single-site flip cross-checked exactly against
\texttt{idmrg.apply\_mpo} at the exact $D=1$ field-polarized point, the
same flip's $\langle S_z\rangle$/$\langle S_z(0)S_z(r)\rangle$/\texttt{eta}
invariants checked at a genuinely entangled $D>1$ TFIM ground state, and a
genuinely bond-growing \texttt{chi\_W>1} two-site gate tiled once per an
\texttt{n\_uc=2} unit cell.

Also ported to \texttt{itensor\_version=3} (\texttt{Chain::vumps\_apply\_mpo},
\texttt{mpscpp3/chain\_session.h}): grows a Chain's own converged $A_L$ by
a caller-supplied \texttt{W\_bulk} (grouped via
\texttt{vumps\_group\_automaton}, reused unmodified from the
Hamiltonian-automaton path since the grouping contraction itself doesn't
care what it's grouping), then shares \texttt{vumps\_imps\_sum}'s own
\texttt{vx\_canonicalize\_n1}/\texttt{vumps\_complete\_mixed\_gauge}
machinery. Same bounded-operator scope restriction, and the same ``no
\texttt{Infinite\_Many\_Body\_Chain}-level wrapper'' caveat as
\texttt{vumps\_imps\_sum} above -- reached directly against a Chain's own
converged VUMPS snapshot:

\begin{lstlisting}
W_bulk_flat = [W0.flatten().tolist(), W1.flatten().tolist()]  # one per unit-cell site
D, d_g, AL, AR, C, AC, eta = ic._session3.vumps_apply_mpo(
    W_bulk_flat, Dw_left, Dw_right, cutoff=1e-12, maxdim=0)
ic._session3.vumps_load_uniform_state(D, d_g, AL.flatten().tolist(),
                                       AR.flatten().tolist(), C.flatten().tolist())
\end{lstlisting}

\texttt{W\_bulk\_flat[p]} is a dense, row-major (Left, in, out, Right)
array (size $\text{Dw\_left}[p] \cdot d_p \cdot d_p \cdot \text{Dw\_right}[p]$)
-- the same convention as \texttt{idmrg.apply\_mpo}'s own ITensor list,
just flattened. See \texttt{tests/test\_vumps\_apply\_mpo\_v3.py} and
\texttt{examples/idmrg/vumps\_apply\_mpo\_v3\_VS\_python/main.py} for the
full cross-check against \texttt{itensor\_version="python"}, including the
same three cases ($D=1$ exact, $D>1$ unitary invariants, \texttt{chi\_W>1}
bond growth at \texttt{n\_uc=2}).

\section{Running the pure-Python backend on a GPU}

\texttt{itensor\_version="python"} can put its tensors on a GPU instead of
in host memory. It is one process-wide switch, set before building a chain:

\begin{verbatim}
from dmrgpy.pyitensor import backend
backend.set_backend("jax")          # "numpy" (default) puts them back
print(backend.device_info())        # "jax: cuda:0" if a device was found

from dmrgpy import spinchain
sc = spinchain.Spin_Chain(["S=1/2"]*30, itensor_version="python")
...                                  # everything else is unchanged
\end{verbatim}

It needs \texttt{jax} with its CUDA plugin installed; nothing else changes
about your script, and results agree with the host run to $\sim 10^{-11}$
or better.

\textbf{Whether it is worth using depends on one number: your bond
dimension.}

\begin{center}
\begin{tabular}{ll}
\hline
bond dimension & what to expect \\
\hline
below $\sim$120 & the GPU is \emph{slower} (0.1--0.5$\times$); use \texttt{"numpy"} \\
$\sim$120--160 & break-even \\
240 & $\sim 5\times$ (ground states and KPM correlators alike) \\
480 & $\sim 20\times$ on a ground state that needs it \\
\hline
\end{tabular}
\end{center}

The reason is not the calculation type but the arithmetic-to-overhead
ratio: each device operation costs $\sim$0.35\,ms no matter how small, so
many small tensors lose and few large ones win. Note this also means
longer chains do not help by themselves (more operations, same size each)
--- bond dimension does.

Ground states, static and KPM dynamical correlators, and \textbf{real-time
evolution} (TDVP, \texttt{timedependent.evolve\_and\_measure} and the
\texttt{submode="TD"} correlators) all run on the device. Time evolution
is the natural fit: entanglement grows with time, so bond dimension
climbs into the paying range by construction, while a 1D ground state's
does not (see the warning at the end of this section).

Three things to set when you use it:

\begin{itemize}
  \item \texttt{backend.set\_pad\_bonds(K)} (with $K$ your bond
    dimension) if the script does \emph{one} calculation and exits: it
    makes every tensor shape identical so the GPU compiles each kernel
    once, worth 1.4--3.1$\times$ on such a run. Skip it for long sweeps
    inside one process, where it costs 6--31\%.
  \item \texttt{backend.set\_jit(True)} fuses the engine's hot inner
    kernels into one compiled kernel each, which is what lowers the
    $\sim$0.35\,ms-per-operation floor. The default \texttt{"auto"}
    turns it on exactly when \texttt{set\_pad\_bonds} is set, because
    compiling only pays off once the tensor shapes stop changing --- and
    measured on an H200, the two together are worth
    \textbf{6.4--12.1$\times$ on a cold run} (a script that starts,
    computes and exits), where either alone is worth only
    2.4--3.8$\times$. Once everything is compiled the trade reverses:
    padding then costs arithmetic on known-zero blocks, so for a long
    sweep inside one process set \texttt{set\_jit(True)}
    \emph{without} padding, which measured 1.3--1.6$\times$ warm.
  \item for a KPM correlator, set \texttt{kpmmaxm} equal to
    \texttt{maxm}, so the ground-state solve and the moment recursion
    share one set of shapes.
\end{itemize}

\texttt{docs/gpu\_cpu\_performance.md} has the full measured comparison,
including which models are and are not meaningful GPU benchmarks --- a
uniform Heisenberg chain's ground state converges at $\chi \sim 60$ and
cannot show a speedup at any \texttt{maxm}.

\section{Performance: BLAS threads}

DMRG spends its time in a great many \emph{small} dense linear-algebra
calls rather than a few large ones -- a two-site tensor at \texttt{maxm=30}
on spin-1/2 sites is a $60\times60$ matrix, and that is what hits
\texttt{svd}/\texttt{eigh}/\texttt{matmul} tens of thousands of times in one
ground-state solve. At that size a multithreaded BLAS can lose more to
thread barriers than it gains: measured with MKL, going from one thread to
two made a $60\times60$ complex \texttt{svd} 1.6x slower and \texttt{eigh}
2.3x slower.

End to end that alone is minor ($\sim$1.13x for
\texttt{itensor\_version="python"}, nothing measurable for \texttt{3}). It
matters when the machine is \textbf{oversubscribed} -- a shared cluster
node, or several dmrgpy runs launched in parallel -- because each process's
BLAS assumes it owns every core. On a 14-core host with another job holding
10 of them, letting MKL use its default thread count turned a 9.4\,s
\texttt{"python"} solve into 28.2\,s, and a 1.6\,s
\texttt{itensor\_version=3} solve into 26.7\,s.

The most reliable fix is to pin threads before numpy is imported:

\begin{lstlisting}[language=bash]
MKL_NUM_THREADS=1 OMP_NUM_THREADS=1 OPENBLAS_NUM_THREADS=1 python3 myscript.py
\end{lstlisting}

For finer control there is an opt-in context manager (needs
\texttt{threadpoolctl}, which is not a declared dependency):

\begin{lstlisting}
from dmrgpy import blasthreads

print(blasthreads.current_hint())     # how threads are configured now
with blasthreads.limit(1):
    e0 = sc.gs_energy()               # restored on exit, including on error
\end{lstlisting}

dmrgpy never changes thread counts on its own: threading does pay off once
the tensors get large enough, and a script may have chosen its setting
deliberately. Treat the numbers above as indicative -- they come from a busy
shared host -- and measure on your own machine before tuning around them.
See \texttt{dmrgpy/blasthreads.py} for the full measurements.


\section{What raises, and what changed in the 2026-08 audit}

A cross-backend audit in August 2026 (\texttt{docs/audit\_2026\_08\_hole\_hunt.md},
which records every reproduction) went looking for calls that silently did
something other than what was asked. Most of what it found is now either
fixed or loud. Two of the fixes \textbf{change numbers} and are worth
knowing about if you have results from before them:

\begin{itemize}
\item \textbf{\texttt{submode="TD"} and \texttt{submode="TDZ"} were a factor
of $\pi$ too large.} The Fourier transform behind them omitted the $1/\pi$
of the convention
$S_{AB}(\omega)=\frac1\pi\mathrm{Re}\int_0^T\!dt\,e^{i\omega t}C(t)w_\delta(t)$
that this guide documents and that every other submode already followed,
and gave the $t=0$ sample full rather than half weight. On a 4-site chain
the exact sum rule $\int S(\omega)d\omega=\langle AB\rangle=0.25$ came out
as 1.28. Peak positions and widths were never affected -- the error is
uniform in $\omega$ -- so only absolute spectral weight changes.
\texttt{sxt\_to\_skomega} and the infinite-chain $S(k,\omega)$ reduction
share the same transform and inherit the fix.
\item \textbf{\texttt{itensor\_version=3} real-time evolution from a
non-unit-norm state.} \texttt{evolution\_ABA} (whose start state is
$A|\mathrm{gs}\rangle$) and \texttt{evolve\_and\_measure(wf=...)} were
rescaled by $1/\|\psi\|^2$ for every $t>0$ under
\texttt{tevol\_method="TDVP"}/\texttt{"TDVP\_GSE"}, while $C(0)$ stayed
correct. Against ED on a 5-site chain the error went from 0.176 to 1.7e-9.
\texttt{"TEBD"}, \texttt{"MPO"}, \texttt{itensor\_version=2} and
\texttt{"python"} were never affected.
\end{itemize}

Beyond those, the following now raise where they used to be silent:

\begin{center}
\begin{tabular}{@{}p{0.34\textwidth}p{0.30\textwidth}p{0.29\textwidth}@{}}
\toprule
\textbf{Call} & \textbf{Was} & \textbf{Now} \\
\midrule
\texttt{sc.maxm = m} with \texttt{m<1} & SIGABRT (v2/v3), wrong energy (\texttt{"python"}) & \texttt{ValueError} \\
changing \texttt{maxm}/\texttt{nsweeps} between \texttt{gs\_energy()} calls & returned the first, less converged energy & re-solves \\
\texttt{sc.tevol\_method = "tdvp"} (any unrecognized name) & ran the legacy MPO-Taylor integrator & \texttt{ValueError} \\
\texttt{submode=} on a non-Hermitian Hamiltonian & every submode returned the CVM resolvent & dispatches; \texttt{NotImplementedError} for the Hermitian-only ones \\
unknown keyword to \texttt{submode="KPM"} & silently discarded & \texttt{TypeError} \\
\texttt{name="ZZ"} with \texttt{submode="KPM"}/\texttt{"EX"}, or any \texttt{mode="ED"} & \texttt{RuntimeError: No active exception to reraise} & works \\
\texttt{submode="CVMimag"}/\texttt{"maxent"} & \texttt{exit()} -- terminated your process & \texttt{NotImplementedError} \\
\texttt{kpm\_energy\_truncate=True} on \texttt{itensor\_version=2} & silently ignored & \texttt{NotImplementedError} \\
\texttt{kpm\_energy\_truncate=True} with far too small a \texttt{kpm\_scale} & plausible spectrum at the wrong energy & \texttt{RuntimeError} in the worst regime \\
an operator on an out-of-range site (\texttt{"python"}) & treated as the identity & \texttt{IndexError} \\
\texttt{get\_excited\_states(n=...)} beyond the Hilbert space & ground on for minutes & \texttt{ValueError} \\
\texttt{Thermal\_Spin\_Chain(..., T<0)} & treated as \texttt{T=0} & \texttt{ValueError} \\
\texttt{Thermal\_Spin\_Chain(..., itensor\_version=...)} & dropped & honored \\
\texttt{get\_distribution*(mode="ED")} & \texttt{AttributeError} deep inside & dispatches, or \texttt{NotImplementedError} \\
\texttt{vev(C[i]*C[i])} & \texttt{AttributeError} on every DMRG backend & \texttt{0} (as ED always answered) \\
\texttt{Many\_Body\_Chain.evolution()} & \texttt{AttributeError} & removed \\
\bottomrule
\end{tabular}
\end{center}

And these combinations now work where they used to fail:

\begin{itemize}
\item \textbf{conserved-sector mode + \texttt{get\_rdm()}} on
\texttt{itensor\_version=3}, which returned uninitialized heap memory
(non-Hermitian, different every run);
\item \textbf{conserved-sector mode + a repeated-site operator product}
such as the four-point correlator's \texttt{Cdag\_0 C\_1 Cdag\_0 C\_1},
which aborted the whole process inside ITensor and now evaluates to 0 as
it should;
\item \textbf{conserved-sector mode + site/pair entropy, mutual
information, and the default \texttt{get\_correlation\_matrix()}}, all of
which raised;
\item \textbf{\texttt{get\_rdm(i=ns-1)}} (the last site) on
\texttt{itensor\_version="python"}.
\end{itemize}

One thing the audit's own \texttt{name=} fix broke, since fixed:
centralizing \texttt{name=} resolution made every submode go through the
same normalizer, whose explicit-pair branch admitted only
\texttt{MultiOperator}s --- so \texttt{name=(A,B)} with the already-built
operators \texttt{toMPO()} returns started raising \texttt{ValueError} on
the paths that had always accepted them (\texttt{mode="ED"}, and
\texttt{submode="EX"}). They are accepted again, per element and
including products of them; see \emph{What \texttt{name=} accepts} in
\S\ref{sec:dynamical} for which submodes can and cannot take one.

\textbf{The ED Lehmann sum is now exact (2026-09-02).}
\texttt{mode="ED"}, \texttt{submode="ED"} --- the reference every other
submode is validated against --- truncated its final-state sum to
\texttt{emu < max(es)}, comparing \emph{absolute} eigenvalues against a
frequency measured from the ground state. Three consequences, all
reproduced: adding a constant to $H$ changed the answer (1.1e-3 at
\texttt{H+3*Id} on a 6-site Heisenberg chain) although a constant can move
no pole; a spectrum sitting entirely above \texttt{max(es)} truncated
itself away and raised numba's \texttt{"zero-size array to reduction"}
from deep inside the sum; and the result depended on how wide an
\texttt{es} window the caller happened to ask for.
\texttt{dynamical\_correlator\_finite\_T} (the $T>0$ route, and the
reference for the METTS correlator) carried the same line.

Both now sum over the full spectrum. \textbf{For almost every existing
caller this changes nothing measurable}: whenever the ground-state energy
is comfortably negative --- the ordinary case, and why this survived since
2025-11 --- the old threshold retained everything that mattered, and the
difference is float noise (1.4e-17 on the 5-orbital IETS atom of
\texttt{tests/test\_atom\_iets.py}, 8.3e-13 on a 6-site Heisenberg chain).
It differs only where the old code was wrong. Nor is it slower: the $T=0$
sum reads only \texttt{nex} rows of one matrix-element array and
\texttt{nex} columns of the other, so only those are now built, which is
cheaper than the full $n \times n$ products the truncated version was
computing. The $T>0$ route \emph{does} cost more, and unavoidably so: it
sums over every initial state with its own Boltzmann weight, so every row
and column is genuinely read and no such shortcut exists. That is
$O(\mathrm{dim}^2)$ in the sum, the same order as the dense
diagonalization it already performs --- 0.47\,s on a 1024-dimensional
Hilbert space, against 0.004\,s for a 6-site chain.

\end{document}
